\documentclass[11pt,a4paper]{article}
\usepackage[utf8]{inputenc}
\usepackage[margin=0.75in]{geometry}
\usepackage{amsmath,amssymb,amsfonts,mathtools,extarrows,esint}

\usepackage{graphicx}
\usepackage[export]{adjustbox}
\usepackage{tcolorbox}
\usepackage{enumitem}
\usepackage{xcolor}
\usepackage{hyperref}
\usepackage{booktabs}
\usepackage{array}
\usepackage{tocloft}

\definecolor{primary}{RGB}{31, 78, 121}
\definecolor{secondary}{RGB}{47, 109, 26}
\definecolor{accent}{RGB}{204, 51, 0}
\definecolor{boxbg}{RGB}{245, 247, 250}

\hypersetup{
    colorlinks=true,
    linkcolor=primary,
    urlcolor=primary,
}

\tcbset{
    solbox/.style={
        colback=boxbg,
        colframe=primary,
        fonttitle=\bfseries,
        coltitle=white,
        boxrule=0.8pt,
        arc=4pt,
        left=8pt,
        right=8pt,
        top=6pt,
        bottom=6pt
    }
}

% Clean Table of Contents styling
\cftsetrmarg{3.5em}
\cftsetpnumwidth{2.5em}
\renewcommand{\cftsecleader}{\cftdotfill{\cftdotsep}}
\renewcommand{\cftsecfont}{\small\normalfont}
\renewcommand{\cftsecpagefont}{\small\bfseries\color{primary}}

\title{\Huge\bfseries GATE EC: Electromagnetics \& Transmission Lines\\[0.3em] \Large Previous Year Solved Questions Archive}
\author{\textbf{Athena Project Archive}}
\date{\today}

\begin{document}
\maketitle
\tableofcontents
\newpage


\section*{Question 1: Transmission Lines (GATE EC 2026)}
\addcontentsline{toc}{section}{Question 1: Transmission Lines (GATE EC 2026)}
\noindent \textbf{\color{primary}MCQ} \hfill \textbf{\color{secondary}2 Marks}


\noindent A complex load (in $\Omega$) is represented as $\Gamma_L = 0.5\angle30^\circ$ on the Smith chart. A co-axial cable with a characteristic impedance of $50 \, \Omega$ is connected to the load. The new input impedance of the load now moves to a diametrically opposite point on the same $\Gamma$ circle on the Smith chart. 

Which option is the nearest input impedance of the cable connected load (in $\Omega$)?


\begin{enumerate}[label=(\Alph*)]
    \item $20.7 - j5.1$
    \item \textbf{(Correct)} $17.7 - j11.8$
    \item $97.5 - j65.0$
    \item $97.5 + j65.0$
\end{enumerate}

\noindent \textbf{Correct Answer:} (B)

\begin{tcolorbox}[solbox, title=Detailed Solution -- Question 1]
\textbf{Given:}

Reflection coefficient of load: $\Gamma_L=0.5\angle30^\circ$

Characteristic impedance: $Z_0=50\Omega$

The input impedance moves to the \textbf{diametrically opposite point} on the same Smith chart circle.

On a Smith chart, diametrically opposite point means:

\[
\Gamma_{in}=0.5\angle(30^\circ+180^\circ)
=
0.5\angle210^\circ
\]

\textbf{Step 1: Convert reflection coefficient into rectangular form}

\[
\Gamma_{in}
=
0.5(\cos210^\circ+j\sin210^\circ)
\]

\[
=
0.5(-0.866-j0.5)
\]

\[
\Gamma_{in}
=
-0.433-j0.25
\]

\textbf{Step 2: Use impedance-reflection coefficient relation}

Substituting values,

\[
Z_{in}
=
50\cdot
\frac{1+(-0.433-j0.25)}
{1-(-0.433-j0.25)}
\]

\[
=
50\cdot
\frac{0.567-j0.25}
{1.433+j0.25}
\]

Multiplying numerator and denominator by conjugate:

\[
Z_{in}
=
50\cdot
\frac{(0.567-j0.25)(1.433-j0.25)}
{(1.433)^2+(0.25)^2}
\]

After simplification, 
\[
Z_{in}
\approx
17.7-j11.8\Omega
\]
\end{tcolorbox}

\vspace{1.5em}
\hrule
\vspace{1.5em}

\section*{Question 2: Waveguides \& Optical Fibers (GATE EC 2026)}
\addcontentsline{toc}{section}{Question 2: Waveguides \& Optical Fibers (GATE EC 2026)}
\noindent \textbf{\color{primary}NAT} \hfill \textbf{\color{secondary}1 Mark}


\noindent The cutoff frequency (in GHz) for the dominant $TE_{10}$ mode of an air-filled rectangular waveguide of inner dimension $0.28 \text{ inch} \times 0.14 \text{ inch}$ is \underline{\hspace{1.5cm}}. (rounded off to two decimal places)


\begin{tcolorbox}[solbox, title=Detailed Solution -- Question 2]
\textbf{Given:}

Inner dimensions: $a = 0.28 \text{ inch},\quad b = 0.14 \text{ inch}$

Air-filled waveguide, dominant mode: $TE_{10}$

\textbf{Step 1: Convert dimensions to meters}

1 inch = 0.0254 m

$a = 0.28 \times 0.0254 = 7.112 \times 10^{-3} \text{ m}$

\textbf{Step 2: Apply cutoff frequency formula for $TE_{10}$}

For $TE_{mn}$ mode, cutoff frequency is:

$f_c = \frac{c}{2}\sqrt{\left(\frac{m}{a}\right)^2 + \left(\frac{n}{b}\right)^2}$

For $TE_{10}$: $m=1, n=0$

$f_c = \frac{c}{2a} = \frac{3 \times 10^8}{2 \times 7.112 \times 10^{-3}}$

\textbf{Step 3: Calculate}

$f_c = \frac{3 \times 10^8}{1.4224 \times 10^{-2}} = 21.09 \times 10^9 \text{ Hz}$
\end{tcolorbox}

\vspace{1.5em}
\hrule
\vspace{1.5em}

\section*{Question 3: Basics of Electromagnetics \& Maxwell's Eqns (GATE EC 2026)}
\addcontentsline{toc}{section}{Question 3: Basics of Electromagnetics \& Maxwell's Eqns (GATE EC 2026)}
\noindent \textbf{\color{primary}MSQ} \hfill \textbf{\color{secondary}1 Mark}


\noindent Which option(s) represents/represent the dielectric loss tangent of a substrate?


\begin{enumerate}[label=(\Alph*)]
    \item \textbf{(Correct)} Ratio of the real to imaginary parts of the total displacement current
    \item \textbf{(Correct)} $(\omega\epsilon'' + \sigma)/(\omega\epsilon')$
    \item Ratio of the electric susceptibility to permittivity
    \item Ratio of the polarization vector $\vec{P}$ to the displacement vector $\vec{D}$
\end{enumerate}

\noindent \textbf{Correct Answer:} (A, B)

\begin{tcolorbox}[solbox, title=Detailed Solution -- Question 3]
\textbf{The loss tangent} $\tan\delta$ is defined as the ratio of the total loss current (conduction + displacement loss) to the lossless displacement current:

$\tan\delta = \frac{\omega\epsilon'' + \sigma}{\omega\epsilon'}$

(A) "Ratio of real to imaginary parts of displacement current", this is exactly $\tan\delta$

(B) $\frac{\omega\epsilon'' + \sigma}{\omega\epsilon'}$, this is the standard formula

(C) "Ratio of electric susceptibility to permittivity", $\chi/\epsilon$, this is NOT the loss tangent

(D) "Ratio of $\vec{P}$ to $\vec{D}$" , $P/D = (\epsilon-\epsilon_0)/\epsilon$, this is NOT the loss tangent

\textbf{Answer:}

Options (A) and (B) correctly represent the dielectric loss tangent.
\end{tcolorbox}

\vspace{1.5em}
\hrule
\vspace{1.5em}

\section*{Question 4: Uniform Plane Waves (GATE EC 2026)}
\addcontentsline{toc}{section}{Question 4: Uniform Plane Waves (GATE EC 2026)}
\noindent \textbf{\color{primary}MCQ} \hfill \textbf{\color{secondary}1 Mark}


\noindent The electric field of a monochromatic plane wave travelling in a lossless isotropic and homogenous medium is given by

 $\vec{E}(z,t) = E_0[\hat{x}\cos(\omega t - kz) + \hat{y}\sin(\omega t - kz)]$

 in a right-handed orthogonal co-ordinate system. 

Which of the following is the correct polarization of the electromagnetic wave?


\begin{enumerate}[label=(\Alph*)]
    \item \textbf{(Correct)} Right-handed circularly polarized
    \item Left-handed circularly polarized
    \item Linearly polarized
    \item Linearly polarized with $-45^\circ$ angle to $\hat{x}$
\end{enumerate}

\noindent \textbf{Correct Answer:} (A)

\begin{tcolorbox}[solbox, title=Detailed Solution -- Question 4]
Given, 

\[
\vec{E}(z,t)=E_0\left[\cos(\omega t-kz)\hat{a}_x+\sin(\omega t-kz)\hat{a}_y\right]
\]

At $z=0$, 

$\vec{E}(t)=E_0\cos\omega t , \hat{a}_x + E_0\sin\omega t , \hat{a}_y$

Here, 

$|E_{x0}|=|E_{y0}|=E_0$

and the phase difference is, 

$|\phi|=90^\circ$

Therefore, the wave is \textbf{circularly polarized}. 

Now, at $t=0$, $\omega t=0$, 

$\vec{E}=E_0\hat{a}_x$

At $t=\dfrac{T}{4}$, $\omega t=\dfrac{\pi}{2}$, 

$\vec{E}=E_0\hat{a}_y$

Thus, the electric field rotates from $\hat{a}_x$ toward $\hat{a}_y$ with time. 

Hence, it is a \textbf{Right-handed circularly polarized} wave.
\end{tcolorbox}

\vspace{1.5em}
\hrule
\vspace{1.5em}

\section*{Question 5: Miscellaneous Topics (GATE EC 2026)}
\addcontentsline{toc}{section}{Question 5: Miscellaneous Topics (GATE EC 2026)}
\noindent \textbf{\color{primary}MCQ} \hfill \textbf{\color{secondary}1 Mark}


\noindent Consider the Friis' transmission equation $P_R = \frac{P_T G_T G_R \lambda^2}{(4\pi D)^2}$, where $P_R$ and $P_T$ are the received and the transmitted powers, respectively. 

$G_T$ and $G_R$ are the gain of transmitting and receiving antennas, respectively, $D$ is the distance between the transmitting and receiving antennas, and $\lambda$ is the wavelength in free space.

Given: $G_T = G_R = 1.0$, $\lambda = 0.30 \text{ m}$ and $P_T = +10 \text{ dBm}$.

 Choose the distance ($D$), in km, from the following options at which the received power, $P_R = -90 \text{ dBm}$?


\begin{enumerate}[label=(\Alph*)]
    \item $\frac{15}{4\pi}$
    \item \textbf{(Correct)} $\frac{15}{2\pi}$
    \item $\frac{75}{2\pi}$
    \item $\frac{3}{4\pi}$
\end{enumerate}

\noindent \textbf{Correct Answer:} (B)

\begin{tcolorbox}[solbox, title=Detailed Solution -- Question 5]
\textbf{Given:}

$G_T=G_R=1$, $\lambda=0.30\text{ m}$, $P_T=+10\text{ dBm}$, $P_R=-90\text{ dBm}$

\textbf{Step 1: Power ratio}

$\frac{P_R}{P_T} = \frac{-90-10}{10}\text{ dB} = -100\text{ dB} = 10^{-10}$

\textbf{Step 2: Friis equation}

\[
\frac{P_R}{P_T} = G_T G_R \left(\frac{\lambda}{4\pi D}\right)^2 = 1\times1\times\left(\frac{0.3}{4\pi D}\right)^2
\]

\textbf{Step 3: Solve for D}

$10^{-10} = \left(\frac{0.3}{4\pi D}\right)^2$

$10^{-5} = \frac{0.3}{4\pi D}$

\[
D = \frac{0.3}{4\pi \times 10^{-5}} = \frac{0.3}{4\pi}\times10^5 = \frac{30000}{4\pi}\text{ m} = \frac{30}{4\pi}\text{ km} = \frac{15}{2\pi}\text{ km}
\]
\end{tcolorbox}

\vspace{1.5em}
\hrule
\vspace{1.5em}

\section*{Question 6: Miscellaneous Topics (GATE EC 2025)}
\addcontentsline{toc}{section}{Question 6: Miscellaneous Topics (GATE EC 2025)}
\noindent \textbf{\color{primary}NAT} \hfill \textbf{\color{secondary}2 Marks}


\noindent Two resistors are connected in a circuit loop of area $5 \mathrm{~m}^{2}$, as shown in the figure below. The circuit loop is placed on the $x-y$ plane.

When a time-varying magnetic flux, with fluxdensity $\mathrm{B}(\mathrm{t})=0.5 \mathrm{t}$ (in Tesla), is applied along the positive z-axis, the magnitude of current I (in Amperes, rounded off to two decimal places) in the loop is \underline{\hspace{1.5cm}} . 

\begin{center}\includegraphics[max width=0.85\linewidth]{images/img\_13.png}\end{center}

\begin{center}\includegraphics[max width=0.85\linewidth]{images/img\_13.png}\end{center}


\begin{tcolorbox}[solbox, title=Detailed Solution -- Question 6]
Given: 

Loop area, $A = 5\ \text{m}^2$ 

Magnetic field, $B(t) = 0.5t\ \text{T}$ 

Total resistance in the loop = $2\ \Omega + 2\ \Omega = 4\ \Omega$ 

We use Faraday's Law of electromagnetic induction:
$\mathcal{E} = -\frac{d\Phi}{dt}$
where $\Phi = B(t) \cdot A$ is the magnetic flux.

So,
$\Phi(t) = B(t) \cdot A = 0.5t \cdot 5 = 2.5t$ 

$\frac{d\Phi}{dt} = \frac{d}{dt}(2.5t) = 2.5\ \text{V}$ 

So, the induced EMF is:
$\mathcal{E} = -2.5\ \text{V}$  

(The negative sign indicates direction, but we are interested in magnitude of current)

 

Now, using Ohm's Law:
$I = \frac{|\mathcal{E}|}{R} = \frac{2.5}{4} = 0.625\ \text{A}$
\end{tcolorbox}

\vspace{1.5em}
\hrule
\vspace{1.5em}

\section*{Question 7: Transmission Lines (GATE EC 2025)}
\addcontentsline{toc}{section}{Question 7: Transmission Lines (GATE EC 2025)}
\noindent \textbf{\color{primary}NAT} \hfill \textbf{\color{secondary}2 Marks}


\noindent A $50 \Omega$ lossless transmission line is terminated with a load $\mathrm{Z}_{\mathrm{L}}$ of $(50-\mathrm{j} 75) \Omega$.

If the average incident power on the line is 10 mW , then the average power delivered to the load (in mW , rounded off to one decimal place) is \underline{\hspace{1.5cm}} .


\begin{tcolorbox}[solbox, title=Detailed Solution -- Question 7]
Step 1: Characteristic Impedance
Given:
$Z_0 = 50\ \Omega$ (lossless line)

$Z_L = 50 - j75\ \Omega$ 

Incident power = $P_{\text{inc} = 10\ \text{mW}$ 

Step 2: Reflection Coefficient
Reflection coefficient $\Gamma = \frac{Z_L - Z_0}{Z_L + Z_0}$ 
First compute:
$Z_L + Z_0 = (50 - j75) + 50 = 100 - j75$ 

$Z_L - Z_0 = (50 - j75) - 50 = -j75$ 

So:
$\Gamma = \frac{-j75}{100 - j75}$ 
Multiply numerator and denominator by the complex conjugate of the denominator:
$\Gamma = \frac{-j75 \cdot (100 + j75)}{(100 - j75)(100 + j75)}$ 
Numerator:
$-j75(100 + j75) = -j7500 - j^2 75 \cdot 75 = -j7500 + 5625$ 
Denominator:
$100^2 + 75^2 = 10000 + 5625 = 15625$ 
So:
$\Gamma = \frac{5625 - j7500}{15625} = 0.36 - j0.48$ 

Step 3: Reflected Power
$|\Gamma|^2 = (0.36)^2 + (0.48)^2 = 0.1296 + 0.2304 = 0.36$ 

So, reflected power = 
\[
P_{\text{ref} = |\Gamma|^2 \cdot P_{\text{inc} = 0.36 \cdot 10 = 3.6\ \text{mW}}
\]
 

Step 4: Delivered Power

\[
P_{\text{delivered} = P_{\text{inc} - P_{\text{ref} = 10 - 3.6 = \boxed{6.4\ \text{mW}}}}
\]
\end{tcolorbox}

\vspace{1.5em}
\hrule
\vspace{1.5em}

\section*{Question 8: Basics of Electromagnetics \& Maxwell's Eqns (GATE EC 2025)}
\addcontentsline{toc}{section}{Question 8: Basics of Electromagnetics \& Maxwell's Eqns (GATE EC 2025)}
\noindent \textbf{\color{primary}MCQ} \hfill \textbf{\color{secondary}2 Marks}


\noindent An electric field of $0.01 \mathrm{~V} / \mathrm{m}$ is applied along the length of a copper wire of circular cross-section with diameter 1 mm . Copper has a conductivity of $5.8 \times 10^{7} \mathrm{~S} / \mathrm{m}$.

The current (in Amperes, rounded off to two decimal places) flowing through the wire is \underline{\hspace{1.5cm}} .


\begin{enumerate}[label=(\Alph*)]
    \item \textbf{(Correct)} 0.46
    \item 1.82
    \item 0.58
    \item 1.12
\end{enumerate}

\noindent \textbf{Correct Answer:} (A)

\begin{tcolorbox}[solbox, title=Detailed Solution -- Question 8]
$\mathrm{J}=\sigma \mathrm{E}=0.01 \times 5.8 \times 10^{7}$

Current,

\[
\begin{aligned}
I &=J \pi r^{2} \\
& =0.01 \times 5.8 \times 10^{7} \times \pi \times\left(0.5 \times 10^{-3}\right)^{2} \\
& =0.045 \times 10^{7} \times 10^{-6}=0.46 \mathrm{~A}
\end{aligned}
\]
\end{tcolorbox}

\vspace{1.5em}
\hrule
\vspace{1.5em}

\section*{Question 9: Basics of Electromagnetics \& Maxwell's Eqns (GATE EC 2025)}
\addcontentsline{toc}{section}{Question 9: Basics of Electromagnetics \& Maxwell's Eqns (GATE EC 2025)}
\noindent \textbf{\color{primary}MCQ} \hfill \textbf{\color{secondary}2 Marks}


\noindent A square metal sheet of $4 \mathrm{~m} \times 4 \mathrm{~m}$ is placed on the $x-y$ plane as shown in the figure below.

If the surface charge density (in $\mu \mathrm{C} / \mathrm{m}^{2}$ ) on the sheet is $\rho_{\mathrm{s}}(\mathrm{x}, \mathrm{y})=4|\mathrm{y}|$, then the total charge (in $\mu \mathrm{C}$, rounded off to the nearest integer) on the sheet is

\underline{\hspace{1.5cm}} . 

\begin{center}\includegraphics[max width=0.85\linewidth]{images/img\_76.png}\end{center}

\begin{center}\includegraphics[max width=0.85\linewidth]{images/img\_76.png}\end{center}


\begin{enumerate}[label=(\Alph*)]
    \item 16
    \item 85
    \item \textbf{(Correct)} 64
    \item 256
\end{enumerate}

\noindent \textbf{Correct Answer:} (C)

\begin{tcolorbox}[solbox, title=Detailed Solution -- Question 9]
\[
\begin{aligned}
\mathrm{Q} &=\iint 4|\mathrm{y}| \mathrm{dx} \mathrm{dy}=4 \iint 4 \mathrm{ydx} \mathrm{dy} \\
& =16 \int \mathrm{dx} \int \mathrm{ydy}=16[\mathrm{x}]_{0}^{2}\left[\frac{\mathrm{y}^{2}}{2}\right]_{0}^{2} \\
& =16 \times 2 \times \frac{2^{2}}{2}=16 \times 4=64 \mu \mathrm{C}
\end{aligned}
\]
\end{tcolorbox}

\vspace{1.5em}
\hrule
\vspace{1.5em}

\section*{Question 10: Transmission Lines (GATE EC 2024)}
\addcontentsline{toc}{section}{Question 10: Transmission Lines (GATE EC 2024)}
\noindent \textbf{\color{primary}NAT} \hfill \textbf{\color{secondary}2 Marks}


\noindent A lossless transmission line with characteristic impedance $Z_{0}=50 \Omega$ is terminated with an unknown load. The magnitude of the reflection co-efficient is $|\Gamma|=0.6$. As one moves towards the generator from the load, the maximum value of the input impedance magnitude looking towards the load (in $\Omega$ ) is \underline{\hspace{1.5cm}}


\begin{tcolorbox}[solbox, title=Detailed Solution -- Question 10]
As we know that

 $Z_{\max }=S Z_{0}$

Now,

 
\[
\begin{aligned}
S & =\frac{1+|\Gamma|}{1-|\Gamma|}=\frac{1+0.6}{1-0.6}=\frac{1.6}{0.4}=4 \\
\therefore \quad Z_{\text {max }} & =4 \times 50=200 \Omega .
\end{aligned}
\]
\end{tcolorbox}

\vspace{1.5em}
\hrule
\vspace{1.5em}

\section*{Question 11: Uniform Plane Waves (GATE EC 2024)}
\addcontentsline{toc}{section}{Question 11: Uniform Plane Waves (GATE EC 2024)}
\noindent \textbf{\color{primary}MCQ} \hfill \textbf{\color{secondary}2 Marks}


\noindent A uniform plane wave with electric field $\vec{E}(x)=A_{y} \hat{a}_{y} e^{-j \frac{2 \pi x}{3}} \mathrm{~V} / \mathrm{m}$ is travelling in the air (relative permittivity, $\epsilon_{r}=1$ and relative permeability, $\mu_{r}=1$ ) in the $+x$ direction $\left(A_{y}\right.$ is a positive constant, $\hat{a}_{y}$ is the unit vector along the $y$ axis). It is incident normally on an ideal electric conductor (conductivity, $\sigma=\infty$ ) at $x=0$. The position of the first null of the total magnetic field in the air (measured from $x=0$, in metres) is \underline{\hspace{1.5cm}}


\begin{enumerate}[label=(\Alph*)]
    \item \textbf{(Correct)} $-\frac{3}{4}$
    \item $-\frac{3}{2}$
    \item $-6$
    \item $-3$
\end{enumerate}

\noindent \textbf{Correct Answer:} (A)

\begin{tcolorbox}[solbox, title=Detailed Solution -- Question 11]
\begin{center}\includegraphics[max width=0.85\linewidth]{images/img\_56.png}\end{center}

\begin{center}\includegraphics[max width=0.85\linewidth]{images/img\_56.png}\end{center}

At perfect conductor,

 
\[
\begin{array}{ll}
& \Gamma=-1 \\
\Rightarrow & \Gamma=1 \mid \pi
\end{array}
\]

As we know that $H_{\min }$ occurs at $E_{\max }$.Hence,

 
\[
\begin{aligned}
& E_{\max }=E_{0}[1+|\Gamma|] \\
& H_{\min }=\frac{E_{0}}{\eta}[1-|\Gamma|] \text { at } 2 \beta x_{\max }=2 n \pi+\theta_{\Gamma}
\end{aligned}
\]

So,

 $2 \beta x_{\max }=2 n \pi+\theta_{\Gamma}$

$\Rightarrow \quad \frac{4 \pi}{\lambda} x_{\max }=2 n \pi+\pi$

 $\Rightarrow \quad x_{\max }=(2 n+1) \frac{\lambda}{4} ; n=0,1,2,3 \ldots$

So, for $1^{\text {st }}$ value of $x$ for $\mathrm{H}$-field to be zero is

 $x_{\max }=\frac{\lambda}{4} \quad[\text { at } n=0]$

Now, from the equation given,

 
\[
\begin{array}{rlrl}
\beta & =\frac{2 \pi}{3} \\
\Rightarrow & \frac{2 \pi}{\lambda} & =\frac{2 \pi}{3} \\
\Rightarrow \quad & \lambda & =3
\end{array}
\]

Hence,

 $x_{\max }=\frac{3}{4}$

If the reference is not at interference then $x=-\frac{3}{4}$.
\end{tcolorbox}

\vspace{1.5em}
\hrule
\vspace{1.5em}

\section*{Question 12: Basics of Electromagnetics \& Maxwell's Eqns (GATE EC 2024)}
\addcontentsline{toc}{section}{Question 12: Basics of Electromagnetics \& Maxwell's Eqns (GATE EC 2024)}
\noindent \textbf{\color{primary}MCQ} \hfill \textbf{\color{secondary}1 Mark}


\noindent Let $\hat{i}$ and $\hat{j}$ be the unit vectors along $x$ and $y$ axes, respectively and let $A$ be a positive constant. Which one of the following statements is true for the vector fields $\vec{F}_{1}=A(\hat{i} y+\hat{j} x)$ and $\vec{F}_{2}=A(\hat{i} y-\hat{j} x) ?$


\begin{enumerate}[label=(\Alph*)]
    \item Both $\vec{F}_{1}$ and $\vec{F}_{2}$ are electrostatic fields.
    \item \textbf{(Correct)} Only $\vec{F}_{1}$ is an electrostatic field.
    \item Only $\vec{F}_{2}$ is an electrostatic field.
    \item Neither $\vec{F}_{1}$ nor $\vec{F}_{2}$ is an electrostatic field.
\end{enumerate}

\noindent \textbf{Correct Answer:} (B)

\begin{tcolorbox}[solbox, title=Detailed Solution -- Question 12]
For an electrostatic field, $\nabla \times \bar{F}=0$

 
\[
\begin{aligned}
& \bar{F}_{1}=A[y \hat{i}+x \hat{j}] \\
& \Rightarrow \\
& \Delta \times \bar{F}_{1}=A\left|\begin{array}{ccc}
\hat{i} & \hat{j} & \hat{k} \\
\partial / \partial x & \partial / \partial y & \partial / \partial z \\
y & x & 0
\end{array}\right| \\
& =A[0 \hat{i}-0 \hat{j}+(1-1) \hat{k}] \\
& =0 \\
& \bar{F}_{2}=A[y \hat{i}-x \hat{j}] \\
& \Rightarrow \\
& \Delta \times \bar{F}_{2}=A\left|\begin{array}{ccc}
\hat{i} & \hat{j} & \hat{k} \\
\partial / \partial x & \partial / \partial y & \partial / \partial z \\
y & -x & 0
\end{array}\right| \\
& =A[0 \hat{i}-0 \hat{j}+(-1-1) \hat{k}] \\
& =-2 A \hat{k}
\end{aligned}
\]

Hence, $\bar{F}_{1}$ is electrostatic, $\bar{F}_{2}$ is not electrostatic.
\end{tcolorbox}

\vspace{1.5em}
\hrule
\vspace{1.5em}

\section*{Question 13: Transmission Lines (GATE EC 2024)}
\addcontentsline{toc}{section}{Question 13: Transmission Lines (GATE EC 2024)}
\noindent \textbf{\color{primary}MCQ} \hfill \textbf{\color{secondary}1 Mark}


\noindent Consider a lossless transmission line terminated with a short circuit as shown in the figure below. As one moves towards the generator from the load, the normalized impedances $z_{\text {in } A}, z_{\text {in } B}, z_{\text {in } C}$ and $z_{\text {in } D}$ (indicated in the figure) are \underline{\hspace{1.5cm}} 

\begin{center}\includegraphics[max width=0.85\linewidth]{images/img\_70.png}\end{center}

\begin{center}\includegraphics[max width=0.85\linewidth]{images/img\_70.png}\end{center}


\begin{enumerate}[label=(\Alph*)]
    \item \[
z_{\text {in } A}=+0.4 j \Omega, z_{\text {in } B}=\infty, z_{\text {in } C}=-0.4 j \Omega, z_{\text {in } D}=0
\]
    \item \[
z_{\text {in } A}=\infty, z_{\text {in } B}=+0.4 j \Omega, z_{\text {in } C}=0, z_{\text {in } D}=+0.4 j \Omega
\]
    \item \[
z_{\text {in } A}=-1 j \Omega, z_{\text {in } B}=0, z_{\text {in } C}=+1 j \Omega, z_{\text {in } D}=\infty
\]
    \item \textbf{(Correct)} \[
z_{\text {in } A}=+1 j \Omega, z_{\text {in } B}=\infty, z_{\text {in } C}=-1 j \Omega, z_{\text {in } D}=0
\]
\end{enumerate}

\noindent \textbf{Correct Answer:} (D)

\begin{tcolorbox}[solbox, title=Detailed Solution -- Question 13]
\begin{center}\includegraphics[max width=0.85\linewidth]{images/img\_21.png}\end{center}

\begin{center}\includegraphics[max width=0.85\linewidth]{images/img\_21.png}\end{center}

$Z_{s / c}=j Z_{0} \tan \beta l$.

$\Rightarrow$ Normalized impedance,

 
\[
\begin{aligned}
& \bar{Z}_{S / C}=\frac{Z_{S / C}}{Z_{0}}=j \tan \beta l \\
& \bar{Z}_{\text {in } A}=j \tan \left(\frac{2 \pi}{\lambda} \cdot \frac{\lambda}{8}\right)=j \tan (\pi / 4)=j 1 \Omega \\
& \bar{Z}_{\text {in } B}=j \tan \left(\frac{2 \pi}{\lambda} \cdot \frac{\lambda}{4}\right)=\infty \\
& \bar{Z}_{\text {in } C}=j \tan \left(\frac{2 \pi}{\lambda} \cdot \frac{3 \lambda}{8}\right)=-j 1 \Omega \\
& \bar{Z}_{\text {in } D}=j \tan \left(\frac{2 \pi}{\lambda} \cdot \frac{\lambda}{2}\right)=0
\end{aligned}
\]
\end{tcolorbox}

\vspace{1.5em}
\hrule
\vspace{1.5em}

\section*{Question 14: Basics of Electromagnetics \& Maxwell's Eqns (GATE EC 2023)}
\addcontentsline{toc}{section}{Question 14: Basics of Electromagnetics \& Maxwell's Eqns (GATE EC 2023)}
\noindent \textbf{\color{primary}NAT} \hfill \textbf{\color{secondary}2 Marks}


\noindent A transparent dielectric coating is applied to glass $\left(\epsilon_{r}=4, \mu_{r}=1\right)$ to eliminate the reflection of red light $\left(\lambda_{0}=0.75 \mu \mathrm{m}\right)$. The minimum thickness of the dielectric coating, in $\mu \mathrm{m}$, that can be used is \underline{\hspace{1.5cm}} (rounded off to two decimal places).


\begin{tcolorbox}[solbox, title=Detailed Solution -- Question 14]
For no reflection, impedance must be matched.

Hence, $\eta_{2}$ acts like a quarter wave impedance transformer.

\begin{center}\includegraphics[max width=0.85\linewidth]{images/img\_89.jpg}\end{center}

\begin{center}\includegraphics[max width=0.85\linewidth]{images/img\_89.jpg}\end{center}

So,

\[
(i) \quad \eta_{2}=\sqrt{\eta_{1} \cdot \eta_{3}} \Rightarrow \epsilon_{r_{2}}=\sqrt{\epsilon_{r_{1}} \cdot \epsilon_{r_{3}}} \Rightarrow \epsilon_{r_{2}}=2
\]

(ii) For impedance matching,

\[
\begin{aligned}
& d=(2 n+1) \frac{\lambda}{4} ; n=0,1,2 \ldots \\
& \lambda=\frac{\lambda_{0}}{\sqrt{\epsilon_{r}}}=\frac{\lambda_{0}}{\sqrt{\epsilon_{r_{2}}}} \\
Here \;\;& \lambda=\frac{0.75 \times 10^{-6}}{\sqrt{2}}=0.53 \times 10^{-6}
\end{aligned}
\]

Hence, for minimum distance, $n=0$

So,$d=\frac{\lambda}{4}=\frac{0.53 \times 10^{-6}}{4}=0.133 \mu \mathrm{m}$
\end{tcolorbox}

\vspace{1.5em}
\hrule
\vspace{1.5em}

\section*{Question 15: Transmission Lines (GATE EC 2023)}
\addcontentsline{toc}{section}{Question 15: Transmission Lines (GATE EC 2023)}
\noindent \textbf{\color{primary}MSQ} \hfill \textbf{\color{secondary}2 Marks}


\noindent The following circuit(s) representing a lumped element equivalent of an infinitesimal section of a transmission line is/are

\begin{center}\includegraphics[max width=0.85\linewidth]{images/img\_18.jpg}\end{center}

\begin{center}\includegraphics[max width=0.85\linewidth]{images/img\_18.jpg}\end{center}


\begin{enumerate}[label=(\Alph*)]
    \item A
    \item \textbf{(Correct)} B
    \item \textbf{(Correct)} C
    \item \textbf{(Correct)} D
\end{enumerate}

\noindent \textbf{Correct Answer:} (B, C, D)


\vspace{1.5em}
\hrule
\vspace{1.5em}

\section*{Question 16: Uniform Plane Waves (GATE EC 2023)}
\addcontentsline{toc}{section}{Question 16: Uniform Plane Waves (GATE EC 2023)}
\noindent \textbf{\color{primary}MSQ} \hfill \textbf{\color{secondary}2 Marks}


\noindent The electric field of a plane electromagnetic wave is 

\[
E=a_{x} C_{1 x} \cos (\omega t-\beta z)+a_{y} C_{1 y} \cos (\omega t-\beta z+\theta) \mathrm{V} / \mathrm{m}
\]
 

 Which of the following combination(s) will give rise to a left handed elliptically polarized (LHEP) wave?


\begin{enumerate}[label=(\Alph*)]
    \item \textbf{(Correct)} $C_{1 x}=1, C_{1 y}=1, \theta=\pi / 4$
    \item \textbf{(Correct)} $C_{1 x}=2, C_{1 y}=1, \theta=\pi / 2$
    \item $C_{1 x}=1, C_{1 y}=2, \theta=3 \pi / 2$
    \item \textbf{(Correct)} $C_{1 x}=2, C_{1 y}=1, \theta=3 \pi / 4$
\end{enumerate}

\noindent \textbf{Correct Answer:} (A, B, D)

\begin{tcolorbox}[solbox, title=Detailed Solution -- Question 16]
Given, 
\[
\quad \vec{E}=\hat{a}_{x} C_{1 x} \cos (\omega t-\beta z)+\hat{a}_{y} C_{1 y} \cos (\omega t-\beta z+\theta)
\]

at $z=0$

\[
\vec{E}=C_{1 x} \cos \omega t \hat{a}_{x}+C_{1 y} \cos (\omega t+\theta) \hat{a}_{y}
\]

Going by options,

Option (A) $\quad \vec{E}=\cos \omega t \hat{a}_{x}+\cos (\omega t+\pi / 4) \hat{a}_{y}$

at $t=0, \omega t=0, \vec{E}=\hat{a}_{x}+\frac{1}{\sqrt{2}} \hat{a}_{y}$

at $t=T / 4, \omega t=\pi / 2, \vec{E}=0-\frac{1}{\sqrt{2}} \hat{a}_{y}$

$\Rightarrow$ Hence, it is LHEP.

\begin{center}\includegraphics[max width=0.85\linewidth]{images/img\_12.jpg}\end{center}

\begin{center}\includegraphics[max width=0.85\linewidth]{images/img\_12.jpg}\end{center}

Option (B) $\quad \vec{E}=2 \cos \omega t \hat{a}_{x}+\cos (\omega t+\pi / 2) \hat{a}_{y}$

at $t=0, \omega t=0, \vec{E}=2 \hat{a}_{x}$

at $t=T / 4, \omega t=\pi / 2, \vec{E}=-1 \hat{a}_{y}$

$\Rightarrow$ Hence, it is LHEP.

\begin{center}\includegraphics[max width=0.85\linewidth]{images/img\_80.jpg}\end{center}

\begin{center}\includegraphics[max width=0.85\linewidth]{images/img\_80.jpg}\end{center}

Option (C) $\quad \vec{E}=\cos \omega t \hat{a}_{x}+2 \cos (\omega t+3 \pi / 2) \hat{a}_{y}$

at $t=0, \quad \quad \omega t=0, \vec{E}=\hat{a}_{x}$

at $t=T / 4, \omega t=\pi / 2, \vec{E}=2 \hat{a}_{y}$

$\Rightarrow$ Hence, it is RHEP.

\begin{center}\includegraphics[max width=0.85\linewidth]{images/img\_85.jpg}\end{center}

\begin{center}\includegraphics[max width=0.85\linewidth]{images/img\_85.jpg}\end{center}

Option (D) $\quad \vec{E}=2 \cos \omega t \hat{a}_{x}+\cos (\omega t+3 \pi / 4) \hat{a}_{y}$

at $t=0, \omega t=0, \vec{E}=2 \hat{a}_{x}-\frac{1}{\sqrt{2}} \hat{a}_{y}$

at 
\[
t=T / 4, \omega t=\pi / 2, \vec{E}=0-\frac{1}{\sqrt{2}} \hat{a}_{y}=\frac{-1}{\sqrt{2}} \hat{a}_{y}
\]

$\Rightarrow$ Hence, it is LHEP.

\begin{center}\includegraphics[max width=0.85\linewidth]{images/img\_30.jpg}\end{center}

\begin{center}\includegraphics[max width=0.85\linewidth]{images/img\_30.jpg}\end{center}

$\therefore$ Option (A), (B) and (D) are correct.
\end{tcolorbox}

\vspace{1.5em}
\hrule
\vspace{1.5em}

\section*{Question 17: Transmission Lines (GATE EC 2023)}
\addcontentsline{toc}{section}{Question 17: Transmission Lines (GATE EC 2023)}
\noindent \textbf{\color{primary}MSQ} \hfill \textbf{\color{secondary}2 Marks}


\noindent The standing wave ratio on a $50 \Omega$ lossless transmission line terminated in an unknown load impedance is found to be 2.0. The distance between successive voltage minima is $30 \mathrm{~cm}$ and the first minimum is located at $10 \mathrm{~cm}$ from the load. $Z_{L}$ can be replaced by an equivalent length $l_{m}$ and terminating resistance $R_{m}$ of the same line. The value of $R_{m}$ and $l_{m}$, respectively, are

\begin{center}\includegraphics[max width=0.85\linewidth]{images/img\_20.jpg}\end{center}

\begin{center}\includegraphics[max width=0.85\linewidth]{images/img\_20.jpg}\end{center}


\begin{enumerate}[label=(\Alph*)]
    \item $R_{m}=100 \Omega, l_{m}=20 \mathrm{~cm}$
    \item \textbf{(Correct)} $R_{m}=25 \Omega, l_{m}=20 \mathrm{~cm}$
    \item \textbf{(Correct)} $R_{m}=100 \Omega, l_{m}=5 \mathrm{~cm}$
    \item $R_{m}=25 \Omega, l_{m}=5 \mathrm{~cm}$
\end{enumerate}

\noindent \textbf{Correct Answer:} (B, C)

\begin{tcolorbox}[solbox, title=Detailed Solution -- Question 17]
Given $S=2, Z_{\min }=10 \mathrm{~cm}, Z_{0}=50 \Omega$

As we know that, $\quad|\Gamma|=\frac{S-1}{S+1}=\frac{2-1}{2+1}=\frac{1}{3}$

Now, distance between successive voltage minima $=30 \mathrm{~cm}$

$\Rightarrow \quad \frac{\lambda}{2}=30 \mathrm{~cm}$

$\Rightarrow \quad \lambda=60 \mathrm{~cm}$

Also, for minima,

$2 \beta Z_{\min }=(2 n+1) \pi+\theta_{\Gamma}$

At $n=0,1 \mathrm{st}$ minima, $Z_{\text {min }}=10 \mathrm{~cm}$

$\frac{4 \pi}{\lambda} Z_{\min }=\pi+\theta_{\Gamma}$

$\Rightarrow \quad \frac{4 \pi}{60} * 10=\pi+\theta_{\Gamma}$

$\Rightarrow \quad \frac{2 \pi}{3}-\pi=\theta_{\Gamma}$

\[
\Rightarrow \quad \theta_{\Gamma}=\frac{-\pi}{3} \quad \therefore \Gamma=\frac{1}{3} \angle-60^{\circ}
\]

Now,$\Gamma=\frac{Z_{L}-Z_{0}}{Z_{L}+Z_{0}}$

\begin{center}\includegraphics[max width=0.85\linewidth]{images/img\_82.jpg}\end{center}

\begin{center}\includegraphics[max width=0.85\linewidth]{images/img\_82.jpg}\end{center}

\[
\begin{array}{ll}\Rightarrow \quad & Z_{L}=Z_{0}\left[\frac{1+\Gamma}{1-\Gamma}\right] \\ \Rightarrow & Z_{L}=50\left[\frac{1+0.33 e^{-j \frac{\pi}{3}}}{1-0.33 e^{-j \frac{\pi}{3}}}\right]\end{array}
\]

$\Rightarrow \quad Z_{L}=67.97 \angle-32.67^{\circ}$

Now 
\[
Z_{\text {in }}=Z_{0}\left[\frac{Z_{L}+j Z_{0} \tan \beta l}{Z_{0}+j Z_{L} \tan \beta l}\right]
\]

\[
\Rightarrow \quad Z_{\text {in }}=50\left[\frac{R_{m}+j 50 \tan \beta l_{m}}{50+j R_{m} \tan \beta l_{m}}\right]
\]

Here, $\quad Z_{\text {in }}=Z_{L}=67.97 \angle-32.67^{\circ}$

Going through options,

\[
\left.\begin{array}{l}R_{m}=100 \Omega \text { and } L_{m}=5 \mathrm{~cm} \\ R_{m}=25 \Omega \text { and } L_{m}=20 \mathrm{~cm}\end{array}\right.
\]
 satisfy this identity, hence option (B) and (C) are correct.
\end{tcolorbox}

\vspace{1.5em}
\hrule
\vspace{1.5em}

\section*{Question 18: Transmission Lines (GATE EC 2023)}
\addcontentsline{toc}{section}{Question 18: Transmission Lines (GATE EC 2023)}
\noindent \textbf{\color{primary}MCQ} \hfill \textbf{\color{secondary}1 Mark}


\noindent A cascade of common-source amplifiers in a unity gain feedback configuration oscillates when


\begin{enumerate}[label=(\Alph*)]
    \item the closed loop gain is less than 1 and the phase shift is less than $180^{\circ}$.
    \item the closed loop gain is greater than 1 and the phase shift is less than $180^{\circ}$.
    \item the closed loop gain is less than 1 and the phase shift is greater than $180^{\circ}$.
    \item \textbf{(Correct)} the closed loop gain is greater than 1 and the phase shift is greater than $180^{\circ}$.
\end{enumerate}

\noindent \textbf{Correct Answer:} (D)

\begin{tcolorbox}[solbox, title=Detailed Solution -- Question 18]
For oscillation,

 Loop gain magnitude $\geq 1$.

Phase of loop gain $=360^{\circ}$

So, correct answer is option (D).

 In options, closes loop gain is mentioned technically it should be loop gain.
\end{tcolorbox}

\vspace{1.5em}
\hrule
\vspace{1.5em}

\section*{Question 19: Miscellaneous Topics (GATE EC 2023)}
\addcontentsline{toc}{section}{Question 19: Miscellaneous Topics (GATE EC 2023)}
\noindent \textbf{\color{primary}MCQ} \hfill \textbf{\color{secondary}1 Mark}


\noindent Consider a narrow band signal, propagating in a lossless dielectric medium $\left(\epsilon_{r}=4, \mu_{r}=1\right)$, with phase velocity $v_{p}$ and group velocity $v_{g}$. Which of the following statement is true? ( $c$ is the velocity of light in vacuum.)


\begin{enumerate}[label=(\Alph*)]
    \item $v_{p} > c, v_{g} > c$
    \item $v_{p} < c, v_{g} > c$
    \item $v_{p} > c, v_{g} < c$
    \item \textbf{(Correct)} $v_{p} < c, v_{g} < c$
\end{enumerate}

\noindent \textbf{Correct Answer:} (D)

\begin{tcolorbox}[solbox, title=Detailed Solution -- Question 19]
Phase velocity, 
\[
V_{p}=\frac{\omega}{\beta}=\frac{\omega}{\omega \sqrt{\mu \epsilon}}=\frac{1}{\sqrt{\mu_{0} \epsilon_{0}}}, \frac{1}{\sqrt{\mu_{r} \epsilon_{r}}}=\frac{C}{\sqrt{\mu_{r} \epsilon_{r}}}
\]

$\therefore \quad V_{p} < C$

Group velocity, 
\[
V_{g}=\frac{d \omega}{d \beta}=\frac{V_{p}}{1-\frac{\omega}{V_{p}} \frac{d V_{p}}{d \omega}}
\]

Here, $\quad V_{p} \neq f(\omega)$

$\therefore \quad V_{g}=V_{p} < C$

Hence, $\quad V_{p} < C$

$\quad \quad V_{g} < C$
\end{tcolorbox}

\vspace{1.5em}
\hrule
\vspace{1.5em}

\section*{Question 20: Basics of Electromagnetics \& Maxwell's Eqns (GATE EC 2022)}
\addcontentsline{toc}{section}{Question 20: Basics of Electromagnetics \& Maxwell's Eqns (GATE EC 2022)}
\noindent \textbf{\color{primary}NAT} \hfill \textbf{\color{secondary}2 Marks}


\noindent In an electrostatic field, the electric displacement density vector, $\vec{D}$, is given by
$\vec{D}(x,y,z)=(x^3\vec{i}+y^3\vec{j}+xy^2\vec{k})C/m^2$
,
where $\vec{i},\vec{j},\vec{k}$ are the unit vectors along x-axis, y-axis, and z-axis, respectively.
Consider a cubical region R centered at the origin with each side of length 1 m, and
vertices at ($\pm 0.5 m, \pm 0.5 m, \pm 0.5 m$). The electric charge enclosed within R is
\underline{\hspace{1.5cm}} C (rounded off to two decimal places).


\begin{tcolorbox}[solbox, title=Detailed Solution -- Question 20]
$\vec{D}(x,y,z)=(x^3\vec{i}+y^3\vec{j}+xy^2\vec{k})c/m^2$

$Q_{enc.}=\int _v \rho _v\cdot dV=\int (\triangledown \cdot \vec{D})dV$

$\triangledown \cdot \vec{D}=3x^2+3y^2$

$dV=dxdydz$

\[
\therefore \; Q_{enc.}=\int _v 3(x^2+y^2)dxdydz =3\left [ \int_{-0.5}^{0.5} x^2 dx  \int_{-0.5}^{0.5}dy \int_{-0.5}^{0.5} dz+\int_{-0.5}^{0.5}dx\int_{-0.5}^{0.5}y^2dy\int_{-0.5}^{0.5}dz \right ]
\]

\[
=3\left [ \frac{x^3}{3}|_{-0.5}^{0.5} \times 1 \times 1 + \frac{y^3}{3}|_{-0.5}^{0.5} \times 1 \times 1\right ] =0.25+0.25=0.5C
\]
\end{tcolorbox}

\vspace{1.5em}
\hrule
\vspace{1.5em}

\section*{Question 21: Waveguides \& Optical Fibers (GATE EC 2022)}
\addcontentsline{toc}{section}{Question 21: Waveguides \& Optical Fibers (GATE EC 2022)}
\noindent \textbf{\color{primary}MSQ} \hfill \textbf{\color{secondary}2 Marks}


\noindent A waveguide consists of two infinite parallel plates (perfect conductors) at a
separation of $10^{-4}$ cm, with air as the dielectric. Assume the speed of light in air to
be $3 \times  10^{8}$ m/s. The frequency/frequencies of TM waves which can propagate in
this waveguide is/are \underline{\hspace{1.5cm}}.


\begin{enumerate}[label=(\Alph*)]
    \item \textbf{(Correct)} $6 \times 10^{15}Hz$
    \item $0.5 \times 10^{12}Hz$
    \item \textbf{(Correct)} $8 \times 10^{14}Hz$
    \item $1 \times 10^{13}Hz$
\end{enumerate}

\noindent \textbf{Correct Answer:} (A, C)

\begin{tcolorbox}[solbox, title=Detailed Solution -- Question 21]
Cut-off frequency.

\[
\begin{aligned}
f_c&=\frac{c}{2a}\;\;\;(m=1)
&=\frac{3 \times 10^8}{2 \times 10^{-4} \times 10^{-2}}\\
&=1.5 \times 10^{14}Hz
\end{aligned}
\]

$f > f_c$ will only propagate
 A and C will propage.
\end{tcolorbox}

\vspace{1.5em}
\hrule
\vspace{1.5em}

\section*{Question 22: Uniform Plane Waves (GATE EC 2022)}
\addcontentsline{toc}{section}{Question 22: Uniform Plane Waves (GATE EC 2022)}
\noindent \textbf{\color{primary}MSQ} \hfill \textbf{\color{secondary}1 Mark}


\noindent Consider the following wave equation,

\[
\frac{\partial^2 f(x,t)}{\partial t^2}=10000\frac{\partial^2 f(x,t)}{\partial x^2}
\]

Which of the given options is/are solution(s) to the given wave equation?


\begin{enumerate}[label=(\Alph*)]
    \item \textbf{(Correct)} $f(x,t)=e^{-(x-100t)^2}+e^{-(x+100t)^2}$
    \item $f(x,t)=e^{-(x-100t)}+0.5e^{-(x+1000t)}$
    \item \textbf{(Correct)} $f(x,t)=e^{-(x-100t)}+\sin (x+100t)$
    \item $f(x,t)=e^{j100 \pi(-100x+t)}+e^{j100 \pi(100x+t)}$
\end{enumerate}

\noindent \textbf{Correct Answer:} (A, C)

\begin{tcolorbox}[solbox, title=Detailed Solution -- Question 22]
As we know, wave equation is given by
$\frac{\partial^2 f(x,t)}{\partial x^2}=\frac{C^2d^2f(x,y)}{dt^2}$
Here, option (A) and (C) are satisfying the above
standard wave equation.
\end{tcolorbox}

\vspace{1.5em}
\hrule
\vspace{1.5em}

\section*{Question 23: Basics of Electromagnetics \& Maxwell's Eqns (GATE EC 2022)}
\addcontentsline{toc}{section}{Question 23: Basics of Electromagnetics \& Maxwell's Eqns (GATE EC 2022)}
\noindent \textbf{\color{primary}MCQ} \hfill \textbf{\color{secondary}1 Mark}


\noindent In a circuit, there is a series connection of an ideal resistor and an ideal capacitor.
The conduction current (in Amperes) through the resistor is $2\sin (t+\pi/2)$.
The displacement current (in Amperes) through the capacitor is \underline{\hspace{1.5cm}}.


\begin{enumerate}[label=(\Alph*)]
    \item $2 \sin (t)$
    \item $2 \sin (t+\pi)$
    \item \textbf{(Correct)} $2 \sin (t+\pi /2)$
    \item 0
\end{enumerate}

\noindent \textbf{Correct Answer:} (C)

\begin{tcolorbox}[solbox, title=Detailed Solution -- Question 23]
\begin{center}\includegraphics[max width=0.85\linewidth]{images/img\_68.jpg}\end{center}

\begin{center}\includegraphics[max width=0.85\linewidth]{images/img\_68.jpg}\end{center}

 In series connection, current pass through each element remain same. Hence, $i_c=i_d$
So, 
\[
i_d=
2\sin (t+\pi/2)
\]
.
\end{tcolorbox}

\vspace{1.5em}
\hrule
\vspace{1.5em}

\section*{Question 24: Antennas \& Radiating Systems (GATE EC 2021)}
\addcontentsline{toc}{section}{Question 24: Antennas \& Radiating Systems (GATE EC 2021)}
\noindent \textbf{\color{primary}NAT} \hfill \textbf{\color{secondary}2 Marks}


\noindent An antenna with a directive gain of $6\text{ dB}$
is radiating a total power of $\text{16 kW}$. The amplitude of the electric field in free space at a distance of $\text{8 km}$
from the antenna in the direction of $\text{6 dB}$
gain (rounded off to three decimal places) is  \underline{\hspace{1.5cm}}  $\text{V} /m$
.


\begin{tcolorbox}[solbox, title=Detailed Solution -- Question 24]
\[
\begin{aligned} G_{d}(\mathrm{~dB}) &=10 \log _{10}\left(G_{d}\right) \\ \Rightarrow\qquad 6 &=10 \log _{10} G_{d} \\ \Rightarrow\qquad G_{d} &=10^{0.6} \\ G_{d} &=3.981 \\ E_{m} &=\frac{\sqrt{60 G_{d} P_{\text {rad }}}}{r} \\ E_{m} &=\frac{\sqrt{60\left(16 \times 10^{3}\right)(3.981)}}{8 \times 10^{3}}=0.244(\mathrm{~V} / \mathrm{m}) \end{aligned}
\]
\end{tcolorbox}

\vspace{1.5em}
\hrule
\vspace{1.5em}

\section*{Question 25: Waveguides \& Optical Fibers (GATE EC 2021)}
\addcontentsline{toc}{section}{Question 25: Waveguides \& Optical Fibers (GATE EC 2021)}
\noindent \textbf{\color{primary}NAT} \hfill \textbf{\color{secondary}2 Marks}


\noindent A standard air-filled rectangular waveguide with dimensions $\text{a} =8 cm, b=4 cm$, operates at $\text{3.4 GHz}$. For the dominant mode of wave propagation, the phase velocity of the signal in $v_{p}$. The value (rounded off to two decimal places) of $v_{p}/c$, where c denotes the velocity of light, is \underline{\hspace{1.5cm}}


\begin{tcolorbox}[solbox, title=Detailed Solution -- Question 25]
\[
f_{c 10}=\frac{c}{2 a}=\frac{3 \times 10^{8}}{2\left(8 \times 10^{-2}\right)}=1.875 \mathrm{GHz}
\]

 Guide phase velocity,$V_{p}=\frac{C}{\sqrt{1-\left(\frac{f_{C 10}}{f}\right)^{2}}}$

\[
\frac{V_{p}}{C}=\frac{1}{\sqrt{1-\left(\frac{f_{C 10}}{f}\right)^{2}}}=\frac{1}{\sqrt{1-\left(\frac{1.875}{3.4}\right)^{2}}}=1.198
\]
\end{tcolorbox}

\vspace{1.5em}
\hrule
\vspace{1.5em}

\section*{Question 26: Basics of Electromagnetics \& Maxwell's Eqns (GATE EC 2021)}
\addcontentsline{toc}{section}{Question 26: Basics of Electromagnetics \& Maxwell's Eqns (GATE EC 2021)}
\noindent \textbf{\color{primary}NAT} \hfill \textbf{\color{secondary}2 Marks}


\noindent For a vector field $D=\rho\cos^{2}\:\varphi \:a_{\rho }+z^{2}\sin^{2}\:\varphi \:a_{\varphi }$
in a cylindrical coordinate system $\left ( \rho ,\varphi ,z \right )$
with unit vectors $a_{\rho },a_{\varphi }$
and $a_{z}$
, the net flux of D leaving the closed surface of the cylinder $\left ( \rho =3, 0\leq z\leq 2 \right )$
(rounded off to two decimal places) is \underline{\hspace{1.5cm}}


\begin{tcolorbox}[solbox, title=Detailed Solution -- Question 26]
Method 1: $\vec{D}=\rho \cos ^{2} \phi \hat{a}_{\rho}+z^{2} \sin ^{2} \phi \hat{a}_{\phi}$
 Electric flux crossing the closed surface is

\[
\begin{aligned} \psi &=\oiint \vec{D} \cdot \overrightarrow{d S}=\iiint(\vec{\nabla} \cdot \vec{D}) d v \\ \vec{\nabla} \cdot \vec{D} &=\frac{1}{\rho} \frac{\partial}{\partial p}\left(\rho D_{\rho}\right)+\frac{1}{\rho} \frac{\partial D_{\phi}}{\partial \phi}+\frac{\partial D_{z}}{\partial z} \\ &=\frac{1}{\rho} \frac{\partial}{\partial p}\left(\rho \rho \cos ^{2} \phi\right)+\frac{1}{\rho} \frac{\partial}{\partial \phi}\left(z^{2} \sin ^{2} \phi\right)+0 \\ &=\frac{1}{\rho}(2 \rho) \cos ^{2} \phi+\frac{1}{\rho} z^{2} 2 \sin \phi \cos \phi=2 \cos ^{2} \phi+\frac{z^{2}}{\rho} \sin 2 \phi\\ \iiint(\vec{\nabla} \cdot \vec{D}) d v &=\iiint 2 \cos ^{2} \phi(\rho d \rho d \phi d z)+\iiint\left(\frac{z^{2}}{\rho} \sin 2 \phi\right) \rho d \rho d \phi d z \\ &=2 \int_{\rho=0}^{3} \rho d \rho \int_{\phi=0}^{2 \pi}\left(\frac{1+\cos 2 \phi}{2}\right) d \phi \int_{z=0}^{2} d z+\int_{\rho=0}^{2} d \rho \int_{\phi=0}^{2 \pi} \sin 2 \phi d \phi \int_{z=0}^{2} z^{2} d z \\ &=2\left(\frac{\rho^{2}}{-2}\right)_{\rho=0}^{3} \frac{1}{2}(2 \pi)(z)_{z=0}^{2}+0 \\ &=2\left(\frac{3^{2}}{2}\right) \pi(2)=18 \pi(\text { Coulomb })=56.55(\text { Coulomb }) \end{aligned}
\]

Method 2: Electric flux crossing the closed surface is
$\psi=\oiint \vec{D} \cdot \overrightarrow{d S}$
 Electric flux crossing $\rho$=3 cylindrical surface is

\[
\begin{aligned} \left.\psi\right|_{\rho=3} &=\oiint\left(\rho \cos ^{2} \phi \hat{a}_{p}\right) \cdot(\rho d \phi d z) \hat{a}_{p} \\ &=3^{2} \int_{\phi=0}^{2 \pi} \cos ^{2} \phi d \phi \int_{z=0}^{2} d z \\ &=9 \frac{1}{2}(2 \pi)(2)=18 \pi(\text { coulomb })=56.55(\text { coulomb }) \end{aligned}
\]
\end{tcolorbox}

\vspace{1.5em}
\hrule
\vspace{1.5em}

\section*{Question 27: Transmission Lines (GATE EC 2021)}
\addcontentsline{toc}{section}{Question 27: Transmission Lines (GATE EC 2021)}
\noindent \textbf{\color{primary}MCQ} \hfill \textbf{\color{secondary}2 Marks}


\noindent The impedance matching network shown n the figure is to match a lossless line having characteristic impedance $Z_{0}= 50 \:\Omega$
with a load impedance $Z_{L}$. A quarter-wave line having a characteristic impedance  $Z_1=75\:\Omega$
is connected to $Z_{L}$. Two stubs having characteristic impedance of $75\:\Omega$
each are connected to this quarter-wave line. One is a short-circuited ($\text{S.C}$) stub of length $0.25\lambda$
connected across PQ and the other one is an open-circuited $(\text{O.C})$
stub of length $0.5\lambda$ connected across RS. 

\begin{center}\includegraphics[max width=0.85\linewidth]{images/img\_92.jpg}\end{center}

\begin{center}\includegraphics[max width=0.85\linewidth]{images/img\_92.jpg}\end{center}

The impedance matching in achieved when the real part of $Z_{L}$
is


\begin{enumerate}[label=(\Alph*)]
    \item \textbf{(Correct)} $112.5\:\Omega$
    \item $75.0\:\Omega$
    \item $50.0\:\Omega$
    \item $33.3\:\Omega$
\end{enumerate}

\noindent \textbf{Correct Answer:} (A)

\begin{tcolorbox}[solbox, title=Detailed Solution -- Question 27]
\[
\begin{aligned} Z_{\text {in } \lambda / 4}&=\frac{Z_{O}^{2}}{Z_{L}}=\frac{(75)^{2}}{0}=\infty \quad\left[\text { for } \frac{\lambda}{4} T_{X} \text { line }\right]\\ \text{Given }Z_{L}\text{ of }\frac{\lambda}{4}\text{ line is 0}(\mathrm{SC})\\ Z_{\text {in }_{\lambda / 2}}&=Z_{L}=\infty\left[\text { for } \frac{\lambda}{2} T_{X} \text { line }\right]\\ \text{Given }Z_{L}\text{ of }\frac{\lambda}{2}\text{ line is }\infty(O . C) \end{aligned}
\]

 The input impedance of $\frac{\lambda}{4}$ transmission line, as well as $\frac{\lambda}{2}$ transmission line is 00 , and they are in parallel with main transmission line, so they are not effective for main.
 Transmission line.
 Final configuration of given line is 

\begin{center}\includegraphics[max width=0.85\linewidth]{images/img\_46.jpg}\end{center}

\begin{center}\includegraphics[max width=0.85\linewidth]{images/img\_46.jpg}\end{center}

 For impedance matching, $Z_{L}=\frac{Z_{1}^{2}}{Z_{0}}=\frac{(75)^{2}}{50}=112.5$
\end{tcolorbox}

\vspace{1.5em}
\hrule
\vspace{1.5em}

\section*{Question 28: Waveguides \& Optical Fibers (GATE EC 2021)}
\addcontentsline{toc}{section}{Question 28: Waveguides \& Optical Fibers (GATE EC 2021)}
\noindent \textbf{\color{primary}NAT} \hfill \textbf{\color{secondary}1 Mark}


\noindent The refractive indices of the core and cladding of an optical fiber are 1.50 and 1.48, respectively. The critical propagation angle, which is defined as the maximum angle that the light beam makes with the axis of the optical fiber to achieve the total internal reflection, (rounded off to two decimal places) is \underline{\hspace{1.5cm}} degree.


\begin{tcolorbox}[solbox, title=Detailed Solution -- Question 28]
Given that
 Refractive index of core $\eta_{1}=1.50$
 Refractive index of clad $\eta_{2}=1.48$
 Critical propagation angle $\left(\theta_{P}\right)$

\[
\theta_{P}=\sin ^{-1}\left[ \frac{ \sqrt{\eta_{1}^{2}-\eta_{2}^{2}}}{\eta_{1}}    \right]=\sin ^{-1} \left[  \frac{  \sqrt{1.5^{2}-1.48^{2}} }{1.5} \right] =9.37
\]
\end{tcolorbox}

\vspace{1.5em}
\hrule
\vspace{1.5em}

\section*{Question 29: Basics of Electromagnetics \& Maxwell's Eqns (GATE EC 2021)}
\addcontentsline{toc}{section}{Question 29: Basics of Electromagnetics \& Maxwell's Eqns (GATE EC 2021)}
\noindent \textbf{\color{primary}NAT} \hfill \textbf{\color{secondary}1 Mark}


\noindent Consider the vector field 
\[
F\:=\:a_{x}\left ( 4y-c_{1}z \right )+\hat{a}_y\left ( 4x + 2z\right )+a_{z}\left ( 2y +z\right )
\]
 in a rectangular coordinate system (x,y,z) with unit vectors $a_{x},\:a_{y}$ and $a_{z}$. If the field F is irrotational (conservative), then the constant $c_{1}$
(in integer) is \underline{\hspace{1.5cm}}


\begin{tcolorbox}[solbox, title=Detailed Solution -- Question 29]
\[
\begin{aligned} \nabla \times \vec{F}&=0\\ \left|\begin{array}{ccc} i & j & k \\ \frac{\partial}{\partial x} & \frac{\partial}{\partial y} & \frac{\partial}{\partial z} \\ 4 y-c_{1} z & 4 x+2 z & 2 y+z \end{array}\right|&=0\\ =i(2-2)-j\left(0+c_{1}\right)+k(4-4)=0 \\ c_{1} &=0 \end{aligned}
\]
\end{tcolorbox}

\vspace{1.5em}
\hrule
\vspace{1.5em}

\section*{Question 30: Waveguides \& Optical Fibers (GATE EC 2021)}
\addcontentsline{toc}{section}{Question 30: Waveguides \& Optical Fibers (GATE EC 2021)}
\noindent \textbf{\color{primary}MCQ} \hfill \textbf{\color{secondary}1 Mark}


\noindent Consider a rectangular coordinate system (x,y,z) with unit vectors $a_{x}\:a_{y}$ and $a_{z}$. A plane wave traveling in the region $z\geq 0$ with electric field vector $E=10\cos\left ( 2\times 10^{8}t\:+\:\beta z\right )a_{y}$ is incident normally on the plane at $z=0$, where $\beta$ is the phase constant. The region $z\geq 0$ is in free space and the region $z < 0$ is filled with a lossless medium (permittivity $\varepsilon \:=\:\varepsilon _{0}$, permeability $\mu \:=\:4\mu _{0}$ , where $\varepsilon _{0}\:=\:8.85\times 10^{-12}\:\text{F} /m$ and $\mu _{0}\:=\:4\pi \times 10^{-7}\:\text{H} /m})$. The value of the reflection coefficient is


\begin{enumerate}[label=(\Alph*)]
    \item \textbf{(Correct)} $\frac{1}{3}$
    \item $\frac{3}{5}$
    \item $\frac{2}{5}$
    \item $\frac{2}{3}$
\end{enumerate}

\noindent \textbf{Correct Answer:} (A)

\begin{tcolorbox}[solbox, title=Detailed Solution -- Question 30]
Given $\vec{E}=10 \cos \left(2 \pi \times 10^{8} t+\beta z\right) \hat{a}_{y}$ for $z \geq 0$ having free space. For $z <  0$ medium has 
$\epsilon_{r 2}=1 ; \mu_{r 2}=4$

\begin{center}\includegraphics[max width=0.85\linewidth]{images/img\_63.jpg}\end{center}

\begin{center}\includegraphics[max width=0.85\linewidth]{images/img\_63.jpg}\end{center}

\[
\Gamma=\frac{\eta_{2}-\eta_{1}}{\eta_{2}+\eta_{1}}=\frac{120 \pi(2)-120 \pi}{120 \pi(2)+120 \pi}=\frac{1}{3}
\]
\end{tcolorbox}

\vspace{1.5em}
\hrule
\vspace{1.5em}

\section*{Question 31: Uniform Plane Waves (GATE EC 2020)}
\addcontentsline{toc}{section}{Question 31: Uniform Plane Waves (GATE EC 2020)}
\noindent \textbf{\color{primary}NAT} \hfill \textbf{\color{secondary}2 Marks}


\noindent The magnetic field of a uniform plane wave in vacuum is given by 

$\vec{H}(x,y,z,t)=(\hat{a}_x+2\hat{a}_y+b\hat{a}_z) \cos (\omega t+3x-y-z)$

 The value of b is \underline{\hspace{1.5cm}}


\begin{tcolorbox}[solbox, title=Detailed Solution -- Question 31]
For uniform plane wave
$\hat{a}_{H}\cdot \hat{a}_{\rho }=0$

$\hat{a}_{H}$is unit vector in magnetic field direction $\hat{a}_{\rho }$ is unit vector in power flow direction

\[
\hat{a}_{H }=\frac{1\hat{a}_{x}+2\hat{a}_{y}+b\hat{a}_{z}}{\sqrt{1^{2}+2^{2}+b^{2}}}
\]

\[
\hat{a}_{\rho }=\frac{-3\hat{a}_{x}+\hat{a}_{y}+\hat{a}_{z}}{\sqrt{3^{2}+1^{2}+1^{2}}}
\]

\[
\hat{a}_{H}\cdot \hat{a}_{\rho }=0 (\hat{a}_{x}+2\hat{a}_{y}+b\hat{a}_{z})\cdot (-3\hat{a}_{x}+\hat{a}_{y}+\hat{a}_{z})=0
\]

$-3+2+b=0$
$b=1$
\end{tcolorbox}

\vspace{1.5em}
\hrule
\vspace{1.5em}

\section*{Question 32: Antennas \& Radiating Systems (GATE EC 2020)}
\addcontentsline{toc}{section}{Question 32: Antennas \& Radiating Systems (GATE EC 2020)}
\noindent \textbf{\color{primary}MCQ} \hfill \textbf{\color{secondary}2 Marks}


\noindent For an infinitesimally small dipole in free space, the electric field $E_\theta$ in the far field
is proportional to $(e^{-jkr}/r) \sin \theta$, where $k = 2 \pi /\lambda$. A vertical infinitesimally small
electric dipole ($\delta l < < \lambda$) is placed at a distance $h(h > 0)$ above an infinite ideal
conducting plane, as shown in the figure. The minimum value of h, for which one
of the maxima in the far field radiation pattern occurs at $\theta =60^{\circ}$, is 

\begin{center}\includegraphics[max width=0.85\linewidth]{images/img\_26.jpg}\end{center}

\begin{center}\includegraphics[max width=0.85\linewidth]{images/img\_26.jpg}\end{center}


\begin{enumerate}[label=(\Alph*)]
    \item \textbf{(Correct)} $\lambda$
    \item 0.5$\lambda$
    \item 0.25$\lambda$
    \item 0.75$\lambda$
\end{enumerate}

\noindent \textbf{Correct Answer:} (A)

\begin{tcolorbox}[solbox, title=Detailed Solution -- Question 32]
\begin{center}\includegraphics[max width=0.85\linewidth]{images/img\_88.jpg}\end{center}

\begin{center}\includegraphics[max width=0.85\linewidth]{images/img\_88.jpg}\end{center}

\[
\begin{aligned}\left | \text{Total E }\right |&=\left | (E_{\text{single element}) \right | \left | (A.F.) \right | \\ \left | (A.F.) \right |&=\frac{\sin\left ( N\frac{\psi }{2} \right )}{\sin (\frac{\psi }{2})}=\frac{\sin \left ( 2\frac{\psi }{2} \right )}{\sin \left ( \frac{\psi }{2} \right )} \\ & =\frac{2\sin \left ( \frac{\psi }{2} \right )\cos (\frac{\psi }{2})}{\sin \frac{\psi }{2}} \\ & =2\cos (\frac{\psi }{2})\\  \left | A.F_{N} \right | &=\frac{A.F}{A.F_{max}}=\frac{2\cos (\frac{\psi }{2})}{2} \\ &=\left | \cos \left ( \frac{\psi }{2} \right ) \right | \\
\text{where, }\psi &=\beta d\cos \theta =\frac{2\pi }{\lambda }(2h)\cos \theta  \\ \left.\begin{matrix} \left | A.F_{N} \right |_{\theta =60^{\circ}} \end{matrix}\right|& =\left | \cos (\frac{2\pi }{\lambda }h\cos 60^{\circ}) \right |\\ &=\left | \cos \left ( \frac{\pi h}{\lambda } \right ) \right |\\ \cos \theta  \text{ is maximum ,where   } \theta &=n\pi \;\; n=0,1,2..\\ \frac{\pi h}{\lambda }&=n\pi \Rightarrow h=n\lambda \\ \Rightarrow \; \text{For } n=1,  \;\; h_{min}&=\lambda \end{aligned}}
\]
\end{tcolorbox}

\vspace{1.5em}
\hrule
\vspace{1.5em}

\section*{Question 33: Transmission Lines (GATE EC 2020)}
\addcontentsline{toc}{section}{Question 33: Transmission Lines (GATE EC 2020)}
\noindent \textbf{\color{primary}NAT} \hfill \textbf{\color{secondary}1 Mark}


\noindent A transmission line of length $3\lambda /4$ and having a characteristic impedance of 50$\Omega$ is
terminated with a load of 400$\Omega$. The impedance (rounded off to two decimal places)
seen at the input end of the transmission line is \underline{\hspace{1.5cm}} $\Omega$.


\begin{tcolorbox}[solbox, title=Detailed Solution -- Question 33]
\begin{center}\includegraphics[max width=0.85\linewidth]{images/img\_38.jpg}\end{center}

\begin{center}\includegraphics[max width=0.85\linewidth]{images/img\_38.jpg}\end{center}

\[
Z_{in}\,\, for \, (l=\lambda /4)=\frac{Z_{0}^{2}}{Z_{L}}=\frac{50^{2}}{400}=\frac{25}{4}=6.25\Omega
\]
\end{tcolorbox}

\vspace{1.5em}
\hrule
\vspace{1.5em}

\section*{Question 34: Transmission Lines (GATE EC 2020)}
\addcontentsline{toc}{section}{Question 34: Transmission Lines (GATE EC 2020)}
\noindent \textbf{\color{primary}MCQ} \hfill \textbf{\color{secondary}1 Mark}


\noindent The impedances Z=jX, for all X in the range ($-\infty ,\infty$), map to the Smith chart as


\begin{enumerate}[label=(\Alph*)]
    \item \textbf{(Correct)} a circle of radius 1 with centre at (0, 0).
    \item a point at the centre of the chart.
    \item a line passing through the centre of the chart
    \item a circle of radius 0.5 with centre at (0.5, 0).
\end{enumerate}

\noindent \textbf{Correct Answer:} (A)

\begin{tcolorbox}[solbox, title=Detailed Solution -- Question 34]
For given impedance Normalized impedance is

\[
\begin{aligned} \frac{Z}{Z_{0}}&=\frac{jX}{Z_{0}}\\ Z&=jX \\ \Rightarrow \; Z&=0+jX \end{aligned}
\]

Normalized Resistance  $=0 \; \Rightarrow \;  r=0$
$X=-\infty \text{ to } \infty$
$r=0$ and X from $-\infty \, to\, \infty$ is a unit circle (radius 1) and centre (0,0) on a complex reflection coeffient plane:

\begin{center}\includegraphics[max width=0.85\linewidth]{images/img\_40.jpg}\end{center}

\begin{center}\includegraphics[max width=0.85\linewidth]{images/img\_40.jpg}\end{center}
\end{tcolorbox}

\vspace{1.5em}
\hrule
\vspace{1.5em}

\section*{Question 35: Waveguides \& Optical Fibers (GATE EC 2019)}
\addcontentsline{toc}{section}{Question 35: Waveguides \& Optical Fibers (GATE EC 2019)}
\noindent \textbf{\color{primary}NAT} \hfill \textbf{\color{secondary}2 Marks}


\noindent A rectangular waveguide of width w and height h has cut-off frequencies for $TE_{10} \; and \;TE_{11}$ modes in the ratio 1:2. The aspect ratio w/h, rounded off to two decimal places, is \underline{\hspace{1.5cm}}


\begin{tcolorbox}[solbox, title=Detailed Solution -- Question 35]
\[
f_{c m n}=\frac{c}{2} \sqrt{\left(\frac{m}{a}\right)^{2}+\left(\frac{n}{b}\right)^{2}}
\]

 For TE mode
$f_{c10}=\frac{c}{2 w} \quad\ldots(i)$
 and For $TE_{11}$ mode.

\[
\begin{aligned} f_{c11} &=\frac{c}{2} \sqrt{\frac{1}{w}^{2}+\left(\frac{1}{h}\right)^{2}} &\ldots(i)\\ &=\frac{c}{2 w} \sqrt{1+\left(\frac{w}{h}\right)^{2}} &\ldots(ii)\\ \text{given}\quad\frac{f_{c10}}{f_{11}} &=\frac{1}{2} &\ldots(iii) \end{aligned}
\]

 put (i). (in) in (ii)

\[
\Rightarrow \frac{\frac{c}{2 w}}{\frac{c}{2 w} \sqrt{1+\left(\frac{w}{h}\right)^{2}}}=\frac{1}{2} \Rightarrow \sqrt{1+\left(\frac{w}{h}\right)^{2}}=2
\]

 On solving above equation, we get,
$\frac{w}{h}=\sqrt{3}=1.732$
\end{tcolorbox}

\vspace{1.5em}
\hrule
\vspace{1.5em}

\section*{Question 36: Waveguides \& Optical Fibers (GATE EC 2019)}
\addcontentsline{toc}{section}{Question 36: Waveguides \& Optical Fibers (GATE EC 2019)}
\noindent \textbf{\color{primary}MCQ} \hfill \textbf{\color{secondary}2 Marks}


\noindent The dispersion equation of a waveguide,which relates the wavenumber kto the frequency $\omega$, is 

$k(\omega) =(1/c)\sqrt{\omega ^2 -{\omega _0}^2}$

 where the speed of light $c=3 \times 10^8$ m/s, and $\omega_0$ is a constant.  If the group velocity is $2 \times 10^8$ m/s, then the phase velocity is


\begin{enumerate}[label=(\Alph*)]
    \item $1.5 \times 10^8$ m/s
    \item $2 \times 10^8$ m/s
    \item $3 \times 10^8$ m/s
    \item \textbf{(Correct)} $4.5 \times 10^8$ m/s
\end{enumerate}

\noindent \textbf{Correct Answer:} (D)

\begin{tcolorbox}[solbox, title=Detailed Solution -- Question 36]
By definition $v_{p}=\frac{\omega}{\beta}=\frac{\omega}{k}$
 where, $k(\omega)=\left(\frac{1}{c}\right) \sqrt{\omega^{2}-\omega_{0}^{2}} \quad$(given)
$\therefore v_{p}=\frac{c}{\sqrt{1-\left(\frac{\omega_{0}}{\omega}\right)^{2}}}$
 by definition,

\[
\begin{aligned} v_{g} &=\frac{d \omega}{d \beta}=\frac{d \omega}{d k} \\ &=\frac{d k}{d \omega}=\frac{1}{c} \frac{1}{2 \sqrt{\omega^{2}-\omega_{0}^{2}}} \times 2 \omega \\ \text{or}\quad v_{g} &=c \sqrt{1-\left(\frac{\omega_{0}}{\omega}\right)^{2}}\\ \because \quad v_{p} \cdot v_{g} &=c^{2} \\ \therefore \quad v_{p} &=\frac{c^{2}}{v_{g}}=\frac{\left(3 \times 10^{8}\right)^{2}}{2 \times 10^{8}} \\ &=4.5 \times 10^{8} \mathrm{m} / \mathrm{sec} \end{aligned}
\]
\end{tcolorbox}

\vspace{1.5em}
\hrule
\vspace{1.5em}

\section*{Question 37: Basics of Electromagnetics \& Maxwell's Eqns (GATE EC 2019)}
\addcontentsline{toc}{section}{Question 37: Basics of Electromagnetics \& Maxwell's Eqns (GATE EC 2019)}
\noindent \textbf{\color{primary}MCQ} \hfill \textbf{\color{secondary}2 Marks}


\noindent Two identical copper wires W1 and W2, placed in parallel as shown in the figure, carry currents I and 2I, respectively, in opposite directions. If the two wires are separated by a distance of 4r, then the magnitude of the magnetic field $\vec{B}$ between the wires at a distance r from W1 is  

\begin{center}\includegraphics[max width=0.85\linewidth]{images/img\_72.jpg}\end{center}

\begin{center}\includegraphics[max width=0.85\linewidth]{images/img\_72.jpg}\end{center}


\begin{enumerate}[label=(\Alph*)]
    \item $\frac{\mu _0 I}{6 \pi r}$
    \item $\frac{6 \mu _0 I}{5 \pi r}$
    \item \textbf{(Correct)} $\frac{5 \mu _0 I}{6 \pi r}$
    \item $\frac{{\mu _0}^2 I^2}{2 \pi r^2}$
\end{enumerate}

\noindent \textbf{Correct Answer:} (C)

\begin{tcolorbox}[solbox, title=Detailed Solution -- Question 37]
\begin{center}\includegraphics[max width=0.85\linewidth]{images/img\_32.jpg}\end{center}

\begin{center}\includegraphics[max width=0.85\linewidth]{images/img\_32.jpg}\end{center}

Magnetic flux density $(\vec{B})$ at r distance due to infinite line carrying current I is $|\vec{B}|=\frac{\mu_{0} I}{2 \pi \rho} .$
$\vec{B}$ at r distance due to $W_{1}$ wire
$=\left|\vec{B}_{1}\right|=\frac{\mu_{0} I}{2 \pi r}\qquad \ldots(i)$
$\vec{B}$ at 3r distance due to $W_{2}$ wire
$=\left|\vec{B}_{2}\right|=\frac{\mu_{0}(2 I)}{2 \pi(3 r)}\qquad \ldots(ii)$
 From right hand thumb rule, $\vec{B}$ due to both lines add in between conductors.

\[
\begin{aligned} {So,}\qquad |\vec{B}|&=\left|\vec{B}_{1}\right|+\left|\vec{B}_{2}\right|\\ \therefore \qquad |\vec{B}|&=\frac{\mu_{0} I}{2 \pi r}+\frac{2 \mu_{0} I}{6 \pi r}=\frac{5 \mu_{0} I}{6 \pi r} \end{aligned}
\]
\end{tcolorbox}

\vspace{1.5em}
\hrule
\vspace{1.5em}

\section*{Question 38: Antennas \& Radiating Systems (GATE EC 2019)}
\addcontentsline{toc}{section}{Question 38: Antennas \& Radiating Systems (GATE EC 2019)}
\noindent \textbf{\color{primary}NAT} \hfill \textbf{\color{secondary}1 Mark}


\noindent Radiation resistance of a small dipole current element of length l at a frequency of 3 GHz is 3 ohms. If the length is changed by 1%, then the percentage change in the radiation resistance, rounded off to two decimal places, is \underline{\hspace{1.5cm}} %.


\begin{tcolorbox}[solbox, title=Detailed Solution -- Question 38]
Radiation resistance of a small dipole current  element of length 'I' is 

\[
\begin{array}{l} R_{\mathrm{rad}}=80 \pi^{2}\left(\frac{l}{\lambda}\right)^{2} \Rightarrow R \propto l^{2} \\ \frac{R_{2}}{R_{1}}=\left(\frac{l_{2}}{l_{1}}\right)^{2} \end{array}
\]

 If length is changed by 1% then percentage change in the radiation resistance.
$\frac{R_{2}}{R_{1}}=\left(\frac{1.01 l}{l}\right)^{2}=1.0201$
$\text { Percentage change in radiation resistance }$

\[
\begin{array}{l} =\frac{R_{2}-R_{1}}{R_{1}} \times 100 \\ =0.0201 \times 100=2.01 \% \end{array}
\]
\end{tcolorbox}

\vspace{1.5em}
\hrule
\vspace{1.5em}

\section*{Question 39: Basics of Electromagnetics \& Maxwell's Eqns (GATE EC 2019)}
\addcontentsline{toc}{section}{Question 39: Basics of Electromagnetics \& Maxwell's Eqns (GATE EC 2019)}
\noindent \textbf{\color{primary}MCQ} \hfill \textbf{\color{secondary}1 Mark}


\noindent In the table shown, List I and List II, respectively, contain terms appearing on the left-hand side and the right-hand side of Maxwell's equations (in their standard form). Match the left-hand side with the corresponding right-hand side.  

\begin{center}\includegraphics[max width=0.85\linewidth]{images/img\_10.jpg}\end{center}

\begin{center}\includegraphics[max width=0.85\linewidth]{images/img\_10.jpg}\end{center}


\begin{enumerate}[label=(\Alph*)]
    \item 1-P, 2-R, 3-Q, 4-S
    \item \textbf{(Correct)} 1-Q, 2-R, 3-P, 4-S
    \item 1-Q, 2-S, 3-P, 4-R
    \item 1-R, 2-Q, 3-S, 4-P
\end{enumerate}

\noindent \textbf{Correct Answer:} (B)

\begin{tcolorbox}[solbox, title=Detailed Solution -- Question 39]
\[
\begin{aligned} \nabla \cdot \vec{D} &=\rho_{v} \\ \nabla \times \vec{E} &=-\frac{\partial \vec{B}}{\partial t} \\ \nabla \cdot \vec{B} &=0 \\ \nabla \times \vec{H} &=\vec{J}+\frac{\partial \vec{D}}{\partial t} \end{aligned}
\]
\end{tcolorbox}

\vspace{1.5em}
\hrule
\vspace{1.5em}

\section*{Question 40: Basics of Electromagnetics \& Maxwell's Eqns (GATE EC 2019)}
\addcontentsline{toc}{section}{Question 40: Basics of Electromagnetics \& Maxwell's Eqns (GATE EC 2019)}
\noindent \textbf{\color{primary}MCQ} \hfill \textbf{\color{secondary}1 Mark}


\noindent What is the electric flux ($\int \vec{E}\cdot d\hat{a}$) through a quarter-cylinder of height H (as shown in the figure) due to an infinitely long line charge along the axis of the cylinder with a charge density of Q?  

\begin{center}\includegraphics[max width=0.85\linewidth]{images/img\_28.jpg}\end{center}

\begin{center}\includegraphics[max width=0.85\linewidth]{images/img\_28.jpg}\end{center}


\begin{enumerate}[label=(\Alph*)]
    \item $\frac{HQ}{\varepsilon _0}$
    \item \textbf{(Correct)} $\frac{HQ}{4\varepsilon _0}$
    \item $\frac{H\varepsilon _0}{4Q}$
    \item $\frac{4H}{Q\varepsilon _0}$
\end{enumerate}

\noindent \textbf{Correct Answer:} (B)

\begin{tcolorbox}[solbox, title=Detailed Solution -- Question 40]
Electric field intensity $(\vec{E})$ at '$\rho$' distance due to infinite long line having line charge density Q is 

\[
\begin{aligned} \vec{E} &=\frac{Q}{2 \pi \varepsilon_{0} \rho} \hat{a}_{\rho} \\ \int \vec{E} \cdot \vec{da} &=\iint \frac{Q}{2 \pi \varepsilon_{0} \rho} \hat{a}_{p} \cdot \rho d \phi d z \hat{a}_{\rho} \\ &=\frac{Q}{2 \pi \varepsilon_{0}} \int_{<\pi / 2>} d \phi \int_{z=0}^{H} d z \\ &=\frac{Q}{2 \pi \varepsilon_{0}}\left(\frac{\pi}{2}\right) H=\frac{H Q}{4 \varepsilon_{0}} \end{aligned}
\]
\end{tcolorbox}

\vspace{1.5em}
\hrule
\vspace{1.5em}

\section*{Question 41: Uniform Plane Waves (GATE EC 2018)}
\addcontentsline{toc}{section}{Question 41: Uniform Plane Waves (GATE EC 2018)}
\noindent \textbf{\color{primary}NAT} \hfill \textbf{\color{secondary}2 Marks}


\noindent A uniform plane wave traveling in free space and having the electric field  
 
\[
\vec{E}=(\sqrt{2}\hat{a}_{x}-\hat{a}_{z})cos[6\sqrt{3}\pi \times 10^{8}t-2\pi (x+\sqrt{2}z)]V/m
\]
 
 is incident on a dielectric medium (relative permittivity $>$1, relative permeability = 1) as shown in the figure and there is no reflected wave.

\begin{center}\includegraphics[max width=0.85\linewidth]{images/img\_15.jpg}\end{center}

\begin{center}\includegraphics[max width=0.85\linewidth]{images/img\_15.jpg}\end{center}

The relative permittivity (correct to two decimal places) of the dielectric medium is\underline{\hspace{1.5cm}}.


\begin{tcolorbox}[solbox, title=Detailed Solution -- Question 41]
The wave has $E_{i}$ direction and propagation  direction both in same plane (ZX).
 The wave is plane of incidence (P) polarized.

\[
\begin{aligned} \beta_{x} &=1, \quad \beta_{2}=\sqrt{2} \\ \tan \theta_{i} &=\frac{\beta_{z}}{\beta_{x}}=\sqrt{2} \end{aligned}
\]

 No reflection in P polarized wave under Brewster's angle of incidence.

\[
\begin{aligned} \theta_{i} &=\theta_{B} \\ \tan \theta_{B} &=\sqrt{2}=\sqrt{\epsilon_{r}} \\ \epsilon_{r} &=2 \end{aligned}
\]
\end{tcolorbox}

\vspace{1.5em}
\hrule
\vspace{1.5em}

\section*{Question 42: Waveguides \& Optical Fibers (GATE EC 2018)}
\addcontentsline{toc}{section}{Question 42: Waveguides \& Optical Fibers (GATE EC 2018)}
\noindent \textbf{\color{primary}NAT} \hfill \textbf{\color{secondary}2 Marks}


\noindent The cutoff frequency of $TE_{01}$ mode of an air filled rectangular waveguide having inner
dimensions $a cm \times b cm ( a > b)$ is twice that of the dominant $TE_{10}$ mode. When the
waveguide is operated at a frequency which is 25% higher than the cutoff frequency of the
dominant mode, the guide wavelength is found to be 4 cm. The value of b (in cm, correct to
two decimal places) is \underline{\hspace{1.5cm}}.


\begin{tcolorbox}[solbox, title=Detailed Solution -- Question 42]
\[
\begin{array}{l} f_{c(01)}=2 f_{c(10)}=\frac{2 c}{2 a}=\frac{c}{a} \\ \frac{c}{2 b}=\frac{c}{a} \Rightarrow a=2 b \Rightarrow b=\frac{a}{2} \end{array}
\]

 Operating frequency,
$f=1.25 f_{c(10)}$
$f_{c(10)}< 1.25 f_{c(10)} < \left[f_{c(01)}=2 f_{c(10)}\right]$
 So, for the given frequency, the waveguide will work in $\mathrm{TE}_{10}$ mode.

\[
\begin{aligned} \text { So, } \lambda_{g} &=\frac{\lambda_{0}}{\sqrt{1-\left(\frac{f_{c(10)}}{f}\right)^{2}}} \\ &=\frac{c / f}{\sqrt{1-\left(\frac{1}{1.25}\right)^{2}}}=\frac{c / f}{0.6} \\ \frac{c}{(1.25) f_{c(10)}(0.6)} &=\lambda_{g}=4 \mathrm{cm} \\ \frac{c}{f_{c(10)}} &=3 \times 10^{-2}=2 a \\ a &=1.5 \mathrm{cm} \\ b &=\frac{a}{2}=0.75 \mathrm{cm} \end{aligned}
\]
\end{tcolorbox}

\vspace{1.5em}
\hrule
\vspace{1.5em}

\section*{Question 43: Uniform Plane Waves (GATE EC 2018)}
\addcontentsline{toc}{section}{Question 43: Uniform Plane Waves (GATE EC 2018)}
\noindent \textbf{\color{primary}MCQ} \hfill \textbf{\color{secondary}2 Marks}


\noindent The distance (in meters) a wave has to propagate in a medium having a skin depth of 0.1 m
so that the amplitude of the wave attenuates by 20 dB, is


\begin{enumerate}[label=(\Alph*)]
    \item 0.12
    \item \textbf{(Correct)} 0.23
    \item 0.46
    \item 2.3
\end{enumerate}

\noindent \textbf{Correct Answer:} (B)

\begin{tcolorbox}[solbox, title=Detailed Solution -- Question 43]
Attenuation constant,

\[
\begin{aligned} \alpha &=\frac{1}{\text { skin depth }}=10 \mathrm{Np} / \mathrm{m} \\ \text { 20log }_{10}\left(\frac{E_{0}}{E_{x}}\right) &=20 \mathrm{dB} \\ \frac{E_{o}}{E_{x}} &=10 \Rightarrow E_{x}=\frac{E_{0}}{10} \\ E_{x} &=E_{0} e^{-\alpha x}=E_{0} e^{-10 x}=\frac{E_{0}}{10} \\ e^{-10 x} &=\frac{1}{10} \\ x &=\frac{1}{10} \ln (10)=0.23 \mathrm{m} \end{aligned}
\]
\end{tcolorbox}

\vspace{1.5em}
\hrule
\vspace{1.5em}

\section*{Question 44: Transmission Lines (GATE EC 2018)}
\addcontentsline{toc}{section}{Question 44: Transmission Lines (GATE EC 2018)}
\noindent \textbf{\color{primary}NAT} \hfill \textbf{\color{secondary}1 Mark}


\noindent A lossy transmission line has resistance per unit length R = 0.05 $\Omega$ /m. The line is
distortionless and has characteristic impedance of 50 $\Omega$. The attenuation constant (in Np/m, correct to three decimal places) of the line is \underline{\hspace{1.5cm}}.


\begin{tcolorbox}[solbox, title=Detailed Solution -- Question 44]
For a distortionless transmission line,
$\frac{L}{R}=\frac{C}{G}$
 Propagation constant,

\[
\begin{aligned} \gamma &=\alpha+j \beta=\sqrt{(R+j \omega L)(G+j \omega C)} \\ &=\sqrt{R G}\left(1+j \omega \frac{L}{R}\right) \end{aligned}
\]

 Attenuation constant,
$\alpha=\sqrt{R G}$
 Characteristic impedance,

\[
\begin{aligned} Z_{0} &=\sqrt{\frac{(R+j \omega L)}{(G+j \omega C)}}=\sqrt{\frac{R}{G}} \\ \sqrt{G} &=\frac{\sqrt{R}}{Z_{0}} \\ \text{So,}\quad\alpha &=\sqrt{R} \cdot \frac{\sqrt{R}}{Z_{0}}=\frac{R}{Z_{0}} \\ &=\frac{0.05}{50}=\frac{0.01}{10}=0.001 \mathrm{Np} / \mathrm{m} \end{aligned}
\]
\end{tcolorbox}

\vspace{1.5em}
\hrule
\vspace{1.5em}

\section*{Question 45: Transmission Lines (GATE EC 2018)}
\addcontentsline{toc}{section}{Question 45: Transmission Lines (GATE EC 2018)}
\noindent \textbf{\color{primary}MCQ} \hfill \textbf{\color{secondary}1 Mark}


\noindent The points P, Q, and R shown on the Smith chart (normalized impedance chart) in the
following figure represent: 

\begin{center}\includegraphics[max width=0.85\linewidth]{images/img\_57.jpg}\end{center}

\begin{center}\includegraphics[max width=0.85\linewidth]{images/img\_57.jpg}\end{center}


\begin{enumerate}[label=(\Alph*)]
    \item P: Open Circuit, Q: Short Circuit, R: Matched Load
    \item P: Open Circuit, Q: Matched Load, R: Short Circuit
    \item \textbf{(Correct)} P: Short Circuit, Q: Matched Load, R: Open Circuit
    \item P: Short Circuit, Q: Open Circuit, R: Matched Load
\end{enumerate}

\noindent \textbf{Correct Answer:} (C)

\begin{tcolorbox}[solbox, title=Detailed Solution -- Question 45]
For Short circuit,
$r=x=0 \quad \Rightarrow \text { Point } " P^{\prime \prime}$
 For Open circuit,
$r=x=\infty \quad \Rightarrow \text { Point }^{\prime \prime} R^{\prime \prime}$
 For Matched load,
$r=1 \text { and } x=0 \Rightarrow \text { Point " } Q^{\prime \prime}$
 P: Short Circuit, Q: Matched Load R: Open circuit
\end{tcolorbox}

\vspace{1.5em}
\hrule
\vspace{1.5em}

\section*{Question 46: Waveguides \& Optical Fibers (GATE EC 2017-SET-2)}
\addcontentsline{toc}{section}{Question 46: Waveguides \& Optical Fibers (GATE EC 2017-SET-2)}
\noindent \textbf{\color{primary}MCQ} \hfill \textbf{\color{secondary}2 Marks}


\noindent Standard air-filled rectangular waveguides of dimensions a = 2.29 cm and b= 1.02 cm are designed for radar applications. It is desired that these waveguides operate only in the dominatnt $TE_{10}$ mode but not higher than 95% of the next higher cutoff frequency. The range of the allowable operating frequency f is.


\begin{enumerate}[label=(\Alph*)]
    \item $8.19 GHz \leq f \leq 13.1 GHz$
    \item \textbf{(Correct)} $8.19 GHz \leq f \leq 12.45 GHz$
    \item $6.55 GHz \leq f \leq 13.1 GHz$
    \item $1.64 GHz \leq f \leq 10.24 GHz$
\end{enumerate}

\noindent \textbf{Correct Answer:} (B)

\begin{tcolorbox}[solbox, title=Detailed Solution -- Question 46]
$a=2.29 \;\mathrm{cm}, b=1.02 \mathrm{cm}$
 Waveguide is operating in $(TE _{10})$ dominant mode 

\[
f_{c}\left(T E_{10}\right)=\frac{c}{2 a}=\frac{3 \times 10^{10}}{2 \times 2.29}=6.55 \mathrm{GHz}
\]

$\Rightarrow 25 \%$ above cut-off frequency
$=1.25 \times 6.55 \mathrm{GHz}=8.1875 \mathrm{GHz}$
 Next higher order mode is $\mathrm{TE}_{20}\left(\because a_{2} > 2b\right)$

\[
\begin{aligned} f_{c}\left(T E_{20}\right) &=\frac{c}{a}=\frac{3 \times 10^{10}}{2.29}=13.1 \mathrm{GH}_{2} \\ \Rightarrow 95 \% \text { of } f_{c}\left(\mathrm{TE}_{20}\right) &=0.95 \times 13.1 \\ &=12.445 \mathrm{GHz} \end{aligned}
\]

 Range of allowable operating frequency
$8.19 \mathrm{GHz} \leq f \leq 12.45 \mathrm{GHz}$
\end{tcolorbox}

\vspace{1.5em}
\hrule
\vspace{1.5em}

\section*{Question 47: Uniform Plane Waves (GATE EC 2017-SET-2)}
\addcontentsline{toc}{section}{Question 47: Uniform Plane Waves (GATE EC 2017-SET-2)}
\noindent \textbf{\color{primary}NAT} \hfill \textbf{\color{secondary}2 Marks}


\noindent The permittivity of water at optical frequencies is 1.75 $\varepsilon_{0}$. It is found that an isotropic light source at a distance d under water forms an illuminated circular area of radius 5m, as shown in the figure. The critical angle is $\theta_{c}$.

\begin{center}\includegraphics[max width=0.85\linewidth]{images/img\_33.jpg}\end{center}

\begin{center}\includegraphics[max width=0.85\linewidth]{images/img\_33.jpg}\end{center}

 The value of d (in meter) is \underline{\hspace{1.5cm}}


\begin{tcolorbox}[solbox, title=Detailed Solution -- Question 47]
Critical angle,

\[
\begin{aligned} \theta_{c} &=\sin ^{-1} \sqrt{\frac{\epsilon_{2}}{\epsilon_{1}}} \\ &=\sin ^{-1} \sqrt{\frac{\epsilon_{0}}{1.75 \epsilon_{0}}}=49.1^{\circ} \end{aligned}
\]

\begin{center}\includegraphics[max width=0.85\linewidth]{images/img\_83.jpg}\end{center}

\begin{center}\includegraphics[max width=0.85\linewidth]{images/img\_83.jpg}\end{center}

\[
\begin{aligned} \tan 49.1 &=1.547=\frac{5}{d} \\ \Rightarrow \quad d &=\frac{5}{1.1547}=4.33 \mathrm{m} \end{aligned}
\]
\end{tcolorbox}

\vspace{1.5em}
\hrule
\vspace{1.5em}

\section*{Question 48: Basics of Electromagnetics \& Maxwell's Eqns (GATE EC 2017-SET-2)}
\addcontentsline{toc}{section}{Question 48: Basics of Electromagnetics \& Maxwell's Eqns (GATE EC 2017-SET-2)}
\noindent \textbf{\color{primary}NAT} \hfill \textbf{\color{secondary}2 Marks}


\noindent An electron ($q_1$) is moving in free space with velocity $10^{5}$ m/s towards a stationary electron ($q_2$) far away. The closest distance that this moving electron gets to the stationary electron before the repulsive force diverts its path is \underline{\hspace{1.5cm}} $\times 10^{-8}$m. 
[Given, mass of electron $m=9.11 \times 10^{-31}$ kg, charge of electron $e=-1.6 \times 10^{-19}$C, and permittivity $\varepsilon_{0}=(1/36\pi)\times 10^{-9}F/m]$


\begin{tcolorbox}[solbox, title=Detailed Solution -- Question 48]
\begin{center}\includegraphics[max width=0.85\linewidth]{images/img\_11.jpg}\end{center}

\begin{center}\includegraphics[max width=0.85\linewidth]{images/img\_11.jpg}\end{center}

r is the distance at which kinetic energy of $q_{1}$ becomes zero [(because kinetic energy (KE) is converted into potential energy (PE)].
 When $q_{1}$ reaches 'r', it starts diverting. 
 Kinetic energy, $K E=\frac{1}{2} m v^{2}$ and work done in moving $q_{1}$ charge to distance 'r' is

\[
\begin{aligned} q_{1} v_{2}&=q_{1} \frac{q_{2}}{4 \pi \varepsilon_{0} r} \\ &\qquad\left(q_{1}=q_{2}=-1.6 \times 10^{-19} \mathrm{C}\right)\\ Now, \quad \frac{1}{2} m v^{2}&=\frac{q_{1} q_{2}}{4 \pi \varepsilon_{0} r} \\ \end{aligned}
\]

\[
\Rightarrow r=\frac{\left(2 \times-1.6 \times 10^{-19}\right) \times\left(-1.6 \times 10^{-19}\right)}{4 \pi \times \frac{10^{-9}}{36 \pi} \times 9.11 \times 10^{-31} \times\left(10^{5}\right)^{2}}
\]

$\simeq 5.06 \times 10^{-8} \mathrm{m}$
\end{tcolorbox}

\vspace{1.5em}
\hrule
\vspace{1.5em}

\section*{Question 49: Transmission Lines (GATE EC 2017-SET-2)}
\addcontentsline{toc}{section}{Question 49: Transmission Lines (GATE EC 2017-SET-2)}
\noindent \textbf{\color{primary}NAT} \hfill \textbf{\color{secondary}1 Mark}


\noindent A two - wire transmission line terminates in a television set. The VSWR measured on the line is 5.8. The percentage of power that is reflected from the television set is \underline{\hspace{1.5cm}}


\begin{tcolorbox}[solbox, title=Detailed Solution -- Question 49]
Given that, VSWR (or) $s=5.8$

\[
\begin{aligned} s &=\frac{1+|\Gamma|}{1-|\Gamma|}=5.8 \\ \Gamma &=\text { Reflection coefficient of } \\ &\quad \text { television set } \\ \Gamma &=\frac{s-1}{s+1}=\frac{4.8}{6.8}=0.70588 \\ \frac{P_{\text {reflected }}}{P_{\text {incident }}} &=\Gamma^{2}=0.4983 \text { or } 49.83 \% \end{aligned}
\]
\end{tcolorbox}

\vspace{1.5em}
\hrule
\vspace{1.5em}

\section*{Question 50: Basics of Electromagnetics \& Maxwell's Eqns (GATE EC 2017-SET-2)}
\addcontentsline{toc}{section}{Question 50: Basics of Electromagnetics \& Maxwell's Eqns (GATE EC 2017-SET-2)}
\noindent \textbf{\color{primary}MCQ} \hfill \textbf{\color{secondary}1 Mark}


\noindent Two conducting spheres S1 and S2 of radii a and b (b>a) respectively, are placed far apart and connected by a long, thin conducting wire, as shown in the figure.

\begin{center}\includegraphics[max width=0.85\linewidth]{images/img\_95.jpg}\end{center}

\begin{center}\includegraphics[max width=0.85\linewidth]{images/img\_95.jpg}\end{center}
 
 For some charge placed on this structure, the potential and surface electric field on S1 are $V_{a}$ and $E_{a}$, and that on S2 are $V_{b}$ and $E_{b}$, respectively, which of the following is CORRECT?


\begin{enumerate}[label=(\Alph*)]
    \item $V_{a}=V_{b} \; and \; E_{a}< E_{b}$
    \item $V_{a}> V_{b} \; and \; E_{a} > E_{b}$
    \item \textbf{(Correct)} $V_{a}=V_{b} \; and \;  E_{a} > E_{b}$
    \item $V_{a}> V_{b} \; and \;  E_{a}=E_{b}$
\end{enumerate}

\noindent \textbf{Correct Answer:} (C)

\begin{tcolorbox}[solbox, title=Detailed Solution -- Question 50]
\begin{center}\includegraphics[max width=0.85\linewidth]{images/img\_45.jpg}\end{center}

\begin{center}\includegraphics[max width=0.85\linewidth]{images/img\_45.jpg}\end{center}

When charge is placed on this structure equilibrium is established such that be  spheres are at same potential i.e.
$V_{a}=V_{b}$

\[
\begin{array}{c} V_{a}=V_{b} \\ \text { So, } \frac{Q_{a}}{4 \pi \epsilon_{o} a}=\frac{Q_{b}}{4 \pi \epsilon_{o} b} \end{array}
\]

$\frac{Q_{b}}{Q_{a}}=\frac{b}{a}$
 Now, surface electric fields.

\[
\frac{E_{a}}{E_{b}}=\left[\frac{Q_{a} / 4 \pi \varepsilon_{o} a^{2}}{Q_{b} / 4 \pi \varepsilon_{o} b^{2}}\right]=\frac{Q_{a} \times b^{2}}{Q_{b} \times a^{2}}=\frac{b}{a}>1
\]

$So, E_{a}>E_{b}$
\end{tcolorbox}

\vspace{1.5em}
\hrule
\vspace{1.5em}

\section*{Question 51: Antennas \& Radiating Systems (GATE EC 2017-SET-1)}
\addcontentsline{toc}{section}{Question 51: Antennas \& Radiating Systems (GATE EC 2017-SET-1)}
\noindent \textbf{\color{primary}MCQ} \hfill \textbf{\color{secondary}2 Marks}


\noindent A half wavelength dipole is kept in the x-y plane and oriented along $45^{\circ}$ from the x-axis. Determine the direction of null in the radiation pattern for $0\leq \phi \leq \pi$. Here the angle $\theta (0\leq \theta \leq \pi)$
is measured from the z-axis, and the angle $\phi(0\leq \phi \leq 2\pi)$ is measured from the x-axis in the x-y plane.


\begin{enumerate}[label=(\Alph*)]
    \item \textbf{(Correct)} $\theta =90^{\circ},\; \; \phi =45^{\circ}$
    \item $\theta =45^{\circ},\; \; \phi =90^{\circ}$
    \item $\theta =90^{\circ},\; \; \phi =135^{\circ}$
    \item $\theta =45^{\circ},\; \; \phi =135^{\circ}$
\end{enumerate}

\noindent \textbf{Correct Answer:} (A)

\begin{tcolorbox}[solbox, title=Detailed Solution -- Question 51]
x-y plane means Z=0 with $Z=r \cos \theta$ (spherical coordinates) $\theta=90^{\circ}$
 With null points along the axis of the dipole it is $\phi=45^{\circ}$ and $225^{\circ}$

\begin{center}\includegraphics[max width=0.85\linewidth]{images/img\_39.jpg}\end{center}

\begin{center}\includegraphics[max width=0.85\linewidth]{images/img\_39.jpg}\end{center}
\end{tcolorbox}

\vspace{1.5em}
\hrule
\vspace{1.5em}

\section*{Question 52: Uniform Plane Waves (GATE EC 2017-SET-1)}
\addcontentsline{toc}{section}{Question 52: Uniform Plane Waves (GATE EC 2017-SET-1)}
\noindent \textbf{\color{primary}MCQ} \hfill \textbf{\color{secondary}2 Marks}


\noindent The expression for an electric field in free space is $E=E_{0}=(\hat{x}+\hat{y}+j2\hat{z})e^{-j(\omega t-kx+ky)}$, where  x, y,
z represent the spatial coordinates, t represents time, and $\omega$, k are constants. This electric field


\begin{enumerate}[label=(\Alph*)]
    \item does not represent a plane wave
    \item represents a circular polarized plane wave propagating normal to the z-axis
    \item \textbf{(Correct)} represents an elliptically polarized plane wave propagating along x-y plane.
    \item represents a linearly polarized plane wave
\end{enumerate}

\noindent \textbf{Correct Answer:} (C)

\begin{tcolorbox}[solbox, title=Detailed Solution -- Question 52]
$\vec{E}$ field direction $\Rightarrow\left(\hat{a}_{x}+\hat{a}_{y}+j 2 \hat{a}_{z}\right)$
 Propagation direction $\Rightarrow k \hat{a}_{x}-k \hat{a}_{y}$
$\vec{E}$ is perpendicular to propagation
$\vec{E} \cdot \vec{P}=0$
 Component in $\hat{a}_{z}$ has magnitude of 2. 
 Component in X-Y plane has magnitude of $\sqrt{2}.$
 These two components are out of phase by $90^{\circ}$ and have unequal amplitudes. So, it is elliptically polarized wave.

\begin{center}\includegraphics[max width=0.85\linewidth]{images/img\_50.jpg}\end{center}

\begin{center}\includegraphics[max width=0.85\linewidth]{images/img\_50.jpg}\end{center}
\end{tcolorbox}

\vspace{1.5em}
\hrule
\vspace{1.5em}

\section*{Question 53: Miscellaneous Topics (GATE EC 2017-SET-1)}
\addcontentsline{toc}{section}{Question 53: Miscellaneous Topics (GATE EC 2017-SET-1)}
\noindent \textbf{\color{primary}NAT} \hfill \textbf{\color{secondary}2 Marks}


\noindent An optical fiber is kept along the $\hat{z}$ direction. The refractive indices for the electric fields along $\hat{x}$ and $\hat{y}$ directions in the fiber are $n_{x}$=1.5000 and $n_{y}$=1.5001, respectively ($n_{x} \neq n_{y}$ due to the imperfection in the fiber cross-section). The free space wavelength of a light wave propagating in the fiber is 1.5$\mu$m. If the light wave is circularly polarized at the input of the fiber, the minimum propagation distance after which it becomes linearly polarized, in centimeter, is \underline{\hspace{1.5cm}}.


\begin{tcolorbox}[solbox, title=Detailed Solution -- Question 53]
Initially the wave is circularly polarized. So, the
initial phase difference between field components
in $\hat{a}_{x}$ direction and $\hat{a}_{y}$ direction is $\frac{\pi}{2}$

To become linearly polarized, the wave must travel
a minimum distance, such that, the phase
difference at that point between the field
components in $\hat{a}_{x}$ direction and $\hat{a}_{y}$ direction is $\pi$

(i.e., the travel of this minimum distance should
provide an additional phase difference of $\pi / 2$
between $\hat{a}_{x}$ and $\hat{a}_{y}$ field components).
So,

\[
\begin{aligned} z_{min}k_z&\sim z_{min}k_y=\frac{\pi}{2} \\ z_{min}\left [ \frac{\omega }{V_{px}}\sim \frac{\omega }{V_{py}} \right ]&= \frac{\pi}{2} \\ 2 \pi z_{min}\left [ \frac{f}{c}\sqrt{\varepsilon _{rx}}\sim \frac{f}{c}\sqrt{\varepsilon _{ry}}\right ]&=\frac{\pi}{2} \\ \frac{4 z_{min}}{\lambda _0}[n_x\sim n_y]&=1 \\ z_{min}&=\frac{\lambda _0}{4(n_x\sim n_y)} \\ &=\frac{1.5}{4(1.5\sim 1.5001)}\mu m \\ &= \frac{1.5}{4(0.0001)}\mu m=\frac{1.5}{4}cm\\ z_{min}&=0.375cm \end{aligned}
\]
\end{tcolorbox}

\vspace{1.5em}
\hrule
\vspace{1.5em}

\section*{Question 54: Antennas \& Radiating Systems (GATE EC 2017-SET-1)}
\addcontentsline{toc}{section}{Question 54: Antennas \& Radiating Systems (GATE EC 2017-SET-1)}
\noindent \textbf{\color{primary}MCQ} \hfill \textbf{\color{secondary}1 Mark}


\noindent Consider a wireless communication link between a transmitter and a receiver located in free space, with finite and strictly positive capacity. If the effective areas of the transmitter and the receiver antennas, and the distance between them are all doubled, and everything else remains
unchanged, the maximum capacity of the wireless link


\begin{enumerate}[label=(\Alph*)]
    \item increases by a factor of 2
    \item decreases by a factor of 2
    \item \textbf{(Correct)} remains unchanged
    \item decreases by a factor of $\sqrt{2}$
\end{enumerate}

\noindent \textbf{Correct Answer:} (C)

\begin{tcolorbox}[solbox, title=Detailed Solution -- Question 54]
As per Friis free space propagation equation,

$W_{r}=\frac{W_{t} A_{e r} A_{e t}}{(\lambda d)^{2}}$

when $A_{e r}$ and $A_{e t}$ are doubled and $d$ also doubled, $W_{r}$ is same. Hence capacity is also same.
\end{tcolorbox}

\vspace{1.5em}
\hrule
\vspace{1.5em}

\section*{Question 55: Transmission Lines (GATE EC 2017-SET-1)}
\addcontentsline{toc}{section}{Question 55: Transmission Lines (GATE EC 2017-SET-1)}
\noindent \textbf{\color{primary}NAT} \hfill \textbf{\color{secondary}1 Mark}


\noindent The voltage of an electromagnetic wave propagating in a coaxial cable with uniform characteristic impedance is $V(\iota )=e^{-y \iota +j\omega t}$ volts, Where $\iota$ is the distance along the length of the cable in meters.  $\gamma =(0.1+j40)m^{-1}$ is the complex propagation constant, and $\omega = 2\pi \times 10^{9} rad/ s$  is the angular frequency. The absolute value of the attenuation in the cable in dB/meter is \underline{\hspace{1.5cm}}.


\begin{tcolorbox}[solbox, title=Detailed Solution -- Question 55]
\[
\begin{aligned} V(l) &=V_{O} e^{-\alpha l} e^{-j\beta l} e^{j\omega t} \\ \text { Attenuation } &=\frac{\mid \text { Input } \mid}{\mid \text { Output } \mid}=\frac{\left|V_{O}(0)\right|}{\left|V_{o}(l)\right|} \end{aligned}
\]

 Attenuation per meter $=\frac{\left|V_{0}\right|}{\left|V_{0}(1 \mathrm{m})\right|}=e^{\alpha}$
 Attenuation in 
\[
\mathrm{dB} / \mathrm{m}=\left(20 \log e^{\alpha}\right) \mathrm{dB} / \mathrm{m}
\]

$=20(0.1) \log e=0.868 \mathrm{dB} / \mathrm{m}$
\end{tcolorbox}

\vspace{1.5em}
\hrule
\vspace{1.5em}

\section*{Question 56: Basics of Electromagnetics \& Maxwell's Eqns (GATE EC 2016-SET-3)}
\addcontentsline{toc}{section}{Question 56: Basics of Electromagnetics \& Maxwell's Eqns (GATE EC 2016-SET-3)}
\noindent \textbf{\color{primary}MCQ} \hfill \textbf{\color{secondary}2 Marks}


\noindent Consider the charge profile shown in the figure. The resultant potential distribution is best described by 

\begin{center}\includegraphics[max width=0.85\linewidth]{images/img\_47.jpg}\end{center}

\begin{center}\includegraphics[max width=0.85\linewidth]{images/img\_47.jpg}\end{center}

\begin{center}\includegraphics[max width=0.85\linewidth]{images/img\_24.jpg}\end{center}

\begin{center}\includegraphics[max width=0.85\linewidth]{images/img\_24.jpg}\end{center}


\begin{enumerate}[label=(\Alph*)]
    \item A
    \item B
    \item C
    \item \textbf{(Correct)} D
\end{enumerate}

\noindent \textbf{Correct Answer:} (D)

\begin{tcolorbox}[solbox, title=Detailed Solution -- Question 56]
Applying Poisson's equations
$\nabla^{2} V=\frac{\partial^{2} V}{\partial x^{2}}=-\frac{\rho_{v}}{\epsilon}=K$
 Constant charge density

\[
\begin{aligned} \frac{\partial V}{\partial x} &=-K x+K^{\prime} \\ V &=\frac{-K x^{2}}{2}+K^{\prime} x+K^{\prime} \end{aligned}
\]

 Towards positive x or negative side.
 It is a second order parabolic increase.
 Due to symmetry of + and - charges K" = 0 is expected with V = O at centre and graph passing through origin.
 Beyond $x > 0 \text{ or } x <  b, \; E = 0$ due to capacitive nature of + and - charges
$V=-\int 0 \cdot \overrightarrow{d l}=\text { Constant }$
 This  constant is same V at x=a or x=b
\end{tcolorbox}

\vspace{1.5em}
\hrule
\vspace{1.5em}

\section*{Question 57: Miscellaneous Topics (GATE EC 2016-SET-3)}
\addcontentsline{toc}{section}{Question 57: Miscellaneous Topics (GATE EC 2016-SET-3)}
\noindent \textbf{\color{primary}NAT} \hfill \textbf{\color{secondary}2 Marks}


\noindent A radar operating at 5 GHz uses a common antenna for transmission and reception. The antenna has a gain of 150 and is aligned for maximum directional radiation and reception to a target 1 km away having radar cross-section of 3 $m^{2}$. If it transmits 100 kW, then the received power (in $\mu$W) is\underline{\hspace{1.5cm}}


\begin{tcolorbox}[solbox, title=Detailed Solution -- Question 57]
\[
\begin{aligned} f&=5 \mathrm{GHz}, G=150, R=1 \mathrm{km} \\ P_{t}&=100 \mathrm{kW}, \sigma=3 \mathrm{m}^{2} \\ P_{r}&= \frac{(G \lambda)^{2} \cdot \sigma}{(4 \pi)^{3} R^{4}} \times P_{t} \\ &=\frac{\left(150 \times \frac{3 \times 10^{8}}{5 \times 10^{9}}\right)^{2} \cdot 3}{(4 \pi)^{3} \cdot 10^{12}} \times 10^{5} \\ &=\frac{81 \times 3 \times 10^{5}}{(4 \pi)^{3} \times 10^{12}} \\ &= 1.224 \times 10^{-8}=0.012 \mu \mathrm{W} \end{aligned}
\]
\end{tcolorbox}

\vspace{1.5em}
\hrule
\vspace{1.5em}

\section*{Question 58: Waveguides \& Optical Fibers (GATE EC 2016-SET-3)}
\addcontentsline{toc}{section}{Question 58: Waveguides \& Optical Fibers (GATE EC 2016-SET-3)}
\noindent \textbf{\color{primary}MCQ} \hfill \textbf{\color{secondary}2 Marks}


\noindent Consider an air-filled rectangular waveguide with dimensions a = 2.286 cm and b = 1.016 cm. The increasing order of the cut-off frequencies for different modes is


\begin{enumerate}[label=(\Alph*)]
    \item $TE_{01} < TE_{10} < TE_{11} < TE_{20}$
    \item $TE_{20} < TE_{11} < TE_{10} < TE_{01}$
    \item \textbf{(Correct)} $TE_{10} < TE_{20} < TE_{01} < TE_{11}$
    \item $TE_{10} < TE_{11} < TE_{20} < TE_{01}$
\end{enumerate}

\noindent \textbf{Correct Answer:} (C)

\begin{tcolorbox}[solbox, title=Detailed Solution -- Question 58]
\[
\begin{aligned} a &=2.286 \mathrm{cm} \\ b &=1.016 \mathrm{cm} \\ f_{c_{T E_{10}}} &=\frac{c}{2 a} \\ f_{c_{T E_{01}}} &=\frac{c}{2 b} \\ f_{c_{T E_{20}}} &=\frac{c}{a} \\ f_{c_{T E_{11}}} &=\frac{c}{2} \sqrt{\left(\frac{1}{a}\right)^{2}+\left(\frac{1}{b}\right)^{2}} \\ \text{As }a > 2 b \Rightarrow f_{c_{T E_{10}}} & < f_{c_{T E_{20}}} < f_{c_{T E_{01}}} < f_{c_{T E_{11}}} \end{aligned}
\]
\end{tcolorbox}

\vspace{1.5em}
\hrule
\vspace{1.5em}

\section*{Question 59: Waveguides \& Optical Fibers (GATE EC 2016-SET-3)}
\addcontentsline{toc}{section}{Question 59: Waveguides \& Optical Fibers (GATE EC 2016-SET-3)}
\noindent \textbf{\color{primary}NAT} \hfill \textbf{\color{secondary}2 Marks}


\noindent Consider an air-filled rectangular waveguide with dimensions a = 2.286 cm and b = 1.016 cm. At 10 GHz operating frequency, the value of the propagation constant (per meter) of the corresponding propagating mode is \underline{\hspace{1.5cm}}


\begin{tcolorbox}[solbox, title=Detailed Solution -- Question 59]
Operation frequency $=10 \mathrm{GHz}$
 Cut off frequency for $\mathrm{TE}_{10}, \mathrm{f}_{c}=\frac{\mathrm{c}}{2 \mathrm{a}}$

\[
\begin{aligned} f_{c} &=\frac{3 \times 10^{8}}{2 \times 2.286 \times 10^{-2}} \\ &=6.57 \mathrm{GHz} \end{aligned}
\]

 Next mode $T E_{20}$ has $f_{c}=13 \mathrm{GHz}$
$10 \mathrm{GHz}$ operates only in dominant mode.

\[
\begin{aligned} r &=\sqrt{\left(\frac{m \pi}{a}\right)^{2}-\omega^{2} \mu \epsilon} \\ &=\sqrt{\left(\frac{\pi}{a}\right)^{2}-\frac{(2 \pi f)^{2}}{\left(3 \times 10^{8}\right)^{2}}} \\ &=\pi \sqrt{\left(\frac{1}{2.286 \times 10^{-2}}\right)^{2}-\left(\frac{2 \times 10 \times 10^{9}}{3 \times 10^{8}}\right)^{2}} \\ &=j 158\left(\mathrm{m}^{-1}\right) \end{aligned}
\]
\end{tcolorbox}

\vspace{1.5em}
\hrule
\vspace{1.5em}

\section*{Question 60: Uniform Plane Waves (GATE EC 2016-SET-3)}
\addcontentsline{toc}{section}{Question 60: Uniform Plane Waves (GATE EC 2016-SET-3)}
\noindent \textbf{\color{primary}MCQ} \hfill \textbf{\color{secondary}1 Mark}


\noindent Faraday's law of electromagnetic induction is mathematically described by which one of the following equations?


\begin{enumerate}[label=(\Alph*)]
    \item $\bigtriangledown \cdot \vec{B}=0$
    \item $\bigtriangledown \cdot \vec{D}=\rho _{v}$
    \item \textbf{(Correct)} $\bigtriangledown  \times \vec{E}=-\frac{\partial \vec{B}}{\partial t}$
    \item \[
\bigtriangledown  \times \vec{H}=\sigma \vec{E}+\frac{\partial \vec{D}}{\partial t}
\]
\end{enumerate}

\noindent \textbf{Correct Answer:} (C)

\begin{tcolorbox}[solbox, title=Detailed Solution -- Question 60]
Rate of change of magnetic field results in induced voltage, 
$\nabla \times \vec{E}=-\frac{\partial \vec{B}}{\partial t}$
\end{tcolorbox}

\vspace{1.5em}
\hrule
\vspace{1.5em}

\section*{Question 61: Uniform Plane Waves (GATE EC 2016-SET-3)}
\addcontentsline{toc}{section}{Question 61: Uniform Plane Waves (GATE EC 2016-SET-3)}
\noindent \textbf{\color{primary}MCQ} \hfill \textbf{\color{secondary}1 Mark}


\noindent If a right-handed circularly polarized wave is incident normally on a plane perfect conductor, then the reflected wave will be


\begin{enumerate}[label=(\Alph*)]
    \item right-handed circularly polarized
    \item \textbf{(Correct)} left-handed circularly polarized
    \item elliptically polarized with a tilt angle of $45^{\circ}$
    \item horizontally polarized
\end{enumerate}

\noindent \textbf{Correct Answer:} (B)

\begin{tcolorbox}[solbox, title=Detailed Solution -- Question 61]
Left circularly polarized, due to direction change
after the reflection from conductor.
\end{tcolorbox}

\vspace{1.5em}
\hrule
\vspace{1.5em}

\section*{Question 62: Basics of Electromagnetics \& Maxwell's Eqns (GATE EC 2016-SET-2)}
\addcontentsline{toc}{section}{Question 62: Basics of Electromagnetics \& Maxwell's Eqns (GATE EC 2016-SET-2)}
\noindent \textbf{\color{primary}MCQ} \hfill \textbf{\color{secondary}2 Marks}


\noindent A positive charge q is placed at x= 0 between two infinite metal plates placed at x =-d and at x = +d respectively. The metal plates lie in the yz plane.

\begin{center}\includegraphics[max width=0.85\linewidth]{images/img\_16.jpg}\end{center}

\begin{center}\includegraphics[max width=0.85\linewidth]{images/img\_16.jpg}\end{center}

 The charge is at rest at t = 0, when a voltage +V is applied to the plate at -d and voltage -V is
applied to the plate at x=+d . Assume that the quantity of the charge q is small enough that it does not perturb the field set up by the metal plates. The time that the charge q takes to reach the right plate is proportional to


\begin{enumerate}[label=(\Alph*)]
    \item d/V
    \item $\sqrt{d}/V$
    \item \textbf{(Correct)} $d/ \sqrt{V}$
    \item $\sqrt{d/V}$
\end{enumerate}

\noindent \textbf{Correct Answer:} (C)

\begin{tcolorbox}[solbox, title=Detailed Solution -- Question 62]
Velocity being free velocity,

\[
\begin{aligned} \frac{1}{2} m v^{2} &=q V \\ v &=\frac{d}{t}=\sqrt{\frac{2 q V}{m}} \\ \Rightarrow t \propto \frac{d}{\sqrt{V}} \end{aligned}
\]
\end{tcolorbox}

\vspace{1.5em}
\hrule
\vspace{1.5em}

\section*{Question 63: Transmission Lines (GATE EC 2016-SET-2)}
\addcontentsline{toc}{section}{Question 63: Transmission Lines (GATE EC 2016-SET-2)}
\noindent \textbf{\color{primary}NAT} \hfill \textbf{\color{secondary}2 Marks}


\noindent A microwave circuit consisting of lossless transmission lines $T_{1} \; and \;  T_{2}$ is shown in the figure. The plot shows the magnitude of the input reflection coefficient  $\tau$ as a function of frequency f. The phase velocity of the signal in the transmission lines is $2*10^{8} m/s$. 

\begin{center}\includegraphics[max width=0.85\linewidth]{images/img\_4.jpg}\end{center}

\begin{center}\includegraphics[max width=0.85\linewidth]{images/img\_4.jpg}\end{center}
 
 The length L (in meters) of $T_2$ is \underline{\hspace{1.5cm}}


\begin{tcolorbox}[solbox, title=Detailed Solution -- Question 63]
At frequency of $1 \mathrm{GHz}$ the line is matched and $| \Gamma |=0$ as seen in the graph.

\begin{center}\includegraphics[max width=0.85\linewidth]{images/img\_7.jpg}\end{center}

\begin{center}\includegraphics[max width=0.85\linewidth]{images/img\_7.jpg}\end{center}

 The line $T_{1}$ is having $50 \Omega$ load parallel to $T_{2}$ input impedance. 
 This load is $50\Omega$ when $T_{2}$ input impedance is infinite. 

\[
\begin{aligned} Z_{\text {in } 2}&=\infty=Z_{L 2} \text{ (open circuit) }\\ T_{2}\text{ has to be a }\frac{\lambda}{2}\text{ line}.&\\ f &=1 \mathrm{GHz} \\ \lambda &=\frac{2 \times 10^{10}}{1 \times 10^{9}} \mathrm{cm} \\ \frac{\lambda}{2} &=L=10 \mathrm{cm}=0.1 \mathrm{m} \end{aligned}
\]
\end{tcolorbox}

\vspace{1.5em}
\hrule
\vspace{1.5em}

\section*{Question 64: Transmission Lines (GATE EC 2016-SET-2)}
\addcontentsline{toc}{section}{Question 64: Transmission Lines (GATE EC 2016-SET-2)}
\noindent \textbf{\color{primary}MCQ} \hfill \textbf{\color{secondary}2 Marks}


\noindent A lossless microstrip transmission line consists of a trace of width w. It is drawn over a practically infinite ground plane and is separated by a dielectric slab of thickness t and relative permittivity $\varepsilon _{r} > 1$.  The inductance per unit length and the characteristic impedance of this line are L and $Z_{0}$, respectively.  

\begin{center}\includegraphics[max width=0.85\linewidth]{images/img\_94.jpg}\end{center}

\begin{center}\includegraphics[max width=0.85\linewidth]{images/img\_94.jpg}\end{center}

 Which one of the following inequalities is always satisfied?


\begin{enumerate}[label=(\Alph*)]
    \item $Z_{0} > \sqrt{\frac{Lt}{\varepsilon _{0}\varepsilon_{r}w }}$
    \item \textbf{(Correct)} $Z_{0} < \sqrt{\frac{Lt}{\varepsilon _{0}\varepsilon_{r} w}}$
    \item $Z_{0} >\sqrt{\frac{Lw}{\varepsilon _{0}\varepsilon_{r}t }}$
    \item $Z_{0} < \sqrt{\frac{Lw}{\varepsilon _{0}\varepsilon_{r}t }}$
\end{enumerate}

\noindent \textbf{Correct Answer:} (B)


\vspace{1.5em}
\hrule
\vspace{1.5em}

\section*{Question 65: Basics of Electromagnetics \& Maxwell's Eqns (GATE EC 2016-SET-2)}
\addcontentsline{toc}{section}{Question 65: Basics of Electromagnetics \& Maxwell's Eqns (GATE EC 2016-SET-2)}
\noindent \textbf{\color{primary}MCQ} \hfill \textbf{\color{secondary}2 Marks}


\noindent The parallel-plate capacitor shown in the figure has movable plates. The capacitor is charged so that the energy stored in it is E when the plate separation is d. The capacitor is then isolated electrically and the plates are moved such that the plate separation becomes 2d. 

\begin{center}\includegraphics[max width=0.85\linewidth]{images/img\_25.jpg}\end{center}

\begin{center}\includegraphics[max width=0.85\linewidth]{images/img\_25.jpg}\end{center}

 At this new plate separation, what is the energy stored in the capacitor, neglecting fringing effects?


\begin{enumerate}[label=(\Alph*)]
    \item \textbf{(Correct)} 2E
    \item $\sqrt{2}E$
    \item E
    \item E/2
\end{enumerate}

\noindent \textbf{Correct Answer:} (A)

\begin{tcolorbox}[solbox, title=Detailed Solution -- Question 65]
\[
\begin{array}{l} \text { Let, } E=E_{1} \text { , Energy } E_{1}=\frac{Q_{1}^{2}}{2 C_{1}} \\ \text { Electrically isolated } \Rightarrow Q_{2}=Q_{1} \\ \begin{aligned} d_{2}=2 d_{1} \Rightarrow \;& C_{2}=\frac{C_{1}}{2} \\ E_{2} &=\frac{Q_{2}^{2}}{2 C_{2}}=\frac{Q_{1}^{2}}{\frac{2 C_{1}}{2}}=2\left(\frac{Q_{1}^{2}}{2 C_{1}}\right) \\ &=2 E_{1}=2 E \end{aligned} \end{array}
\]
\end{tcolorbox}

\vspace{1.5em}
\hrule
\vspace{1.5em}

\section*{Question 66: Miscellaneous Topics (GATE EC 2016-SET-2)}
\addcontentsline{toc}{section}{Question 66: Miscellaneous Topics (GATE EC 2016-SET-2)}
\noindent \textbf{\color{primary}NAT} \hfill \textbf{\color{secondary}1 Mark}


\noindent Light from free space is incident at an angle $\theta _{i}$  to the normal of the facet of a step-index large core optical fibre. The core and cladding refractive indices are $n _{1}=1.5 \; and \; n_{2}=1.4$, respectively.

\begin{center}\includegraphics[max width=0.85\linewidth]{images/img\_66.jpg}\end{center}

\begin{center}\includegraphics[max width=0.85\linewidth]{images/img\_66.jpg}\end{center}
 
The maximum value of $\theta _{i}$ (in degrees) for which the incident light will be guided in the core of the fibre is \underline{\hspace{1.5cm}}


\begin{tcolorbox}[solbox, title=Detailed Solution -- Question 66]
\[
\begin{aligned} & \sin \alpha_{\max } &=\sqrt{n_{1}^{2}-n_{2}^{2}}=\sqrt{1.5^{2}-1.4^{2}} \\ \Rightarrow & \alpha_{\max } &=\sin ^{-1}(0.5385)=32.58^{\circ} \end{aligned}
\]
\end{tcolorbox}

\vspace{1.5em}
\hrule
\vspace{1.5em}

\section*{Question 67: Uniform Plane Waves (GATE EC 2016-SET-2)}
\addcontentsline{toc}{section}{Question 67: Uniform Plane Waves (GATE EC 2016-SET-2)}
\noindent \textbf{\color{primary}MCQ} \hfill \textbf{\color{secondary}1 Mark}


\noindent Let the electric field vector of a plane electromagnetic wave propagating in a homogenous medium be expressed as $E=\hat{x}E_{x}e^{-j(\omega t-\beta z)}$, where the propagation constant $\beta$ is a function of the angular frequency $\omega$. Assume that $\beta (\omega)$
and $E_{x}$ are known and are real. From the information available, which one of the following CANNOT be determined?


\begin{enumerate}[label=(\Alph*)]
    \item The type of polarization of the wave.
    \item The group velocity of the wave
    \item The phase velocity of the wave.
    \item \textbf{(Correct)} The power flux through the z = 0 plane.
\end{enumerate}

\noindent \textbf{Correct Answer:} (D)

\begin{tcolorbox}[solbox, title=Detailed Solution -- Question 67]
$\beta(\omega)$ is known so $V_{g}$ can be calculated. 
$v_{p}=\omega / \beta$ can be calculated.
 Polarization can be identified.
$\mu_{r}$ and $\epsilon_{r}$ cannot be found, due to which power flux cannot be calculated as power flux 

\[
P=\frac{1}{2} \frac{|E|^{2}}{\eta}, \text{ where }\eta=120 \pi \times \sqrt{\frac{\mu_{r}}{\epsilon_{r}}}
\]
\end{tcolorbox}

\vspace{1.5em}
\hrule
\vspace{1.5em}

\section*{Question 68: Basics of Electromagnetics \& Maxwell's Eqns (GATE EC 2016-SET-2)}
\addcontentsline{toc}{section}{Question 68: Basics of Electromagnetics \& Maxwell's Eqns (GATE EC 2016-SET-2)}
\noindent \textbf{\color{primary}MCQ} \hfill \textbf{\color{secondary}1 Mark}


\noindent A uniform and constant magnetic field $B=\hat{z}B$ exists in the $\hat{z}$ direction in vacuum. A particle of mass m with a small charge q is introduced into this region with an initial velocity $v=\hat{x}v_{x}+\hat{z}v_{z}$. Given that B, m, q, $v_{x}$ and $v_{z}$ are all non-zero, which one of the following describes the eventual trajectory of the particle?


\begin{enumerate}[label=(\Alph*)]
    \item \textbf{(Correct)} Helical motion in the $\hat{z}$ direction
    \item Circular motion in the xy plane
    \item Linear motion in the $\hat{z}$ direction
    \item Linear motion in the $\hat{x}$ direction
\end{enumerate}

\noindent \textbf{Correct Answer:} (A)

\begin{tcolorbox}[solbox, title=Detailed Solution -- Question 68]
$\mathrm{Ba}_{\mathrm{z}}$ magnetic field 
$v_{x} a_{x}+v_{z} a_{z}$ velocity 

\[
\begin{aligned} F&=Q(v \times B) \text{ by Lorentz's law}\\ &=Q\left(v_{x} a_{x}+v_{z} a_{z}\right) \times B a_{z} \\ F_{y} &=Q v_{x} \cdot B\left(-a_{y}\right) \end{aligned}
\]

 This results in a circular path in the XY plane with  $v_{z} a_{z}$ component causing a linear path. 
 Both result in a helical path along Z axis.
\end{tcolorbox}

\vspace{1.5em}
\hrule
\vspace{1.5em}

\section*{Question 69: Antennas \& Radiating Systems (GATE EC 2016-SET-1)}
\addcontentsline{toc}{section}{Question 69: Antennas \& Radiating Systems (GATE EC 2016-SET-1)}
\noindent \textbf{\color{primary}MCQ} \hfill \textbf{\color{secondary}2 Marks}


\noindent The far-zone power density radiated by a helical antenna is approximated as: 

\[
\vec{W}_{rad}=\vec{W}_{average}\approx \hat{a}_{r}C_{o}\frac{1}{r^{2}}cos^{4}\theta
\]

  The radiated power density is symmetrical with respect to $\phi$ and exists only in the upper hemisphere: $0\leq \theta \leq \frac{\pi }{2};0\leq \phi \leq 2\pi;C_{o}$ is a constant. The power radiated by the antenna (in watts) and the maximum directivity of the antenna, respectively, are


\begin{enumerate}[label=(\Alph*)]
    \item $1.5C_{o},10dB$
    \item \textbf{(Correct)} $1.256C_{o},10dB$
    \item $1.256C_{o},12dB$
    \item $1.5C_{o},12dB$
\end{enumerate}

\noindent \textbf{Correct Answer:} (B)

\begin{tcolorbox}[solbox, title=Detailed Solution -- Question 69]
\[
\begin{aligned} \text { Power radiated }&=\int P_{\mathrm{rad}} \cdot d s \\ =& \int_{\theta=0}^{\pi / 2} \int_{\theta=0}^{2 \pi} C_{0} \frac{1}{r^{2}} \cos ^{4} \theta \cdot r^{2} \sin \theta d \theta d \phi \\ =& 2 \pi \cdot \frac{C_{0}}{r^{2}} \cdot r^{2} \int_{\theta=0}^{\pi / 2} \cos ^{4} \theta \cdot d(-\cos \theta) \\ =&\left.2 \pi \cdot C_{0}\left(-\frac{\cos ^{5} \theta}{5}\right)\right|_{0} ^{\pi / 2} \\ =& \frac{2 \pi}{5} C_{0}=1.256 C_{0} \\ \text{Directivity}=& \frac{4 \pi \cdot U(\theta, \phi)}{\int(\theta, \phi) \cdot d \Omega} \\ =& \frac{4 \pi \cdot C_{0} \cdot \cos ^{4} \theta}{\int_{\theta=0}^{\pi / 2}\int_{\phi=0}^{2\pi}C_{0} \cdot \cos ^{4} \theta \cdot \sin \theta d \theta d \phi} \\ &=\frac{2 \cos ^{5} \theta}{1 / 5}=10 \cos ^{5} \theta \\ \max _{\operatorname{value}}=& 10 \\ \Rightarrow \max \text { value }(\text { in } \mathrm{dB}) &=10 \log _{10} 10=10 \mathrm{dB} \end{aligned}
\]
\end{tcolorbox}

\vspace{1.5em}
\hrule
\vspace{1.5em}

\section*{Question 70: Antennas \& Radiating Systems (GATE EC 2016-SET-1)}
\addcontentsline{toc}{section}{Question 70: Antennas \& Radiating Systems (GATE EC 2016-SET-1)}
\noindent \textbf{\color{primary}MSQ} \hfill \textbf{\color{secondary}2 Marks}


\noindent The electric field of a uniform plane wave travelling along the negative z direction is given by the following equation: 

$\vec{E}^{1}_{w}=(\hat{a}_{x}+j\hat{a}_{y})E_{0}e^{jkz}$ 

This wave is incident upon a receiving antenna placed at the origin and whose radiated electric field towards the incident wave is given by the following equation:  

$\vec{E}_{a}=(\hat{a}_{x}+2\hat{a}_{y})E_{I}\frac{1}{r}e^{-jkr}$    

The polarization of the incident wave, the polarization of the antenna and losses due to the polarization mismatch are, respectively,


\begin{enumerate}[label=(\Alph*)]
    \item Linear, Circular (clockwise), -5dB
    \item Circular (clockwise), Linear, -5dB
    \item \textbf{(Correct)} Circular (clockwise), Linear, -3dB
    \item \textbf{(Correct)} Circular (anti clockwise), Linear, -3dB
\end{enumerate}

\noindent \textbf{Correct Answer:} (C, D)

\begin{tcolorbox}[solbox, title=Detailed Solution -- Question 70]
$\vec{E}_{T}=\left(\hat{a}_{x}+j \hat{a}_{y}\right) E_{0} e^{j k z}$
$\Rightarrow$ Wave contains two orthogonal components and  Y component leads X component leads by $90^{\circ}$ and also wave is travelling in negative z-direction. 
$\vec{E}_{R}=\left(\hat{a}_{x}+2 \hat{a}_{y}\right) E_{1} \frac{1}{r} e^{-j k r}$
$\Rightarrow$ Wave contains two orthogonal components  with unequal amplitudes and both are in-phase.
$\Rightarrow$ Linear polarization.
 Polarization Loss Factor, $\mathrm{PLF}=\left|\hat{E}_{T} \cdot \hat{E}_{R}\right|^{2}$

\[
\begin{aligned} Where, E_{T}&=\frac{\hat{a}_{x}+j \hat{a}_{y}}{\sqrt{2}} \\ E_{R}&=\frac{\hat{a}_{x}+2 \hat{a}_{y}}{\sqrt{5}} \\ P L&= \left|\frac{1+j 2^{2}}{\sqrt{10} }\right|=\frac{5}{10}=\frac{1}{2} \\ \Rightarrow \quad P L F(d B)&=10 \log \frac{1}{2}=-3 \mathrm{dB} \end{aligned}
\]
\end{tcolorbox}

\vspace{1.5em}
\hrule
\vspace{1.5em}

\section*{Question 71: Antennas \& Radiating Systems (GATE EC 2016-SET-1)}
\addcontentsline{toc}{section}{Question 71: Antennas \& Radiating Systems (GATE EC 2016-SET-1)}
\noindent \textbf{\color{primary}NAT} \hfill \textbf{\color{secondary}2 Marks}


\noindent Two lossless X-band horn antennas are separated by a distance of 200$\lambda$. The amplitude reflection coefficients at the terminals of the transmitting and receiving antennas are 0.15 and 0.18, respectively. The maximum directivities of the transmitting and receiving antennas (over the isotropic antenna) are 18 dB and 22 dB, respectively. Assuming that the input power in the lossless transmission line connected to the antenna is 2 W, and that the antennas are perfectly aligned and
polarization matched, the power ( in mW) delivered to the load at the receiver is \underline{\hspace{1.5cm}}


\begin{tcolorbox}[solbox, title=Detailed Solution -- Question 71]
\begin{center}\includegraphics[max width=0.85\linewidth]{images/img\_59.jpg}\end{center}

\begin{center}\includegraphics[max width=0.85\linewidth]{images/img\_59.jpg}\end{center}

\[
\begin{aligned} G_{t}&=10^{1.8}, G_{r}=10^{2.2} \\ P_{r}&=\frac{\left(1-\left|\Gamma_{t}\right|^{2}\right)\left(1-\left|\Gamma_{r}\right|^{2}\right) G_{t} G_{r}}{\left(\frac{4 \pi d}{\lambda}\right)^{2}} \cdot P_{t}\\ &=\frac{\left(1-|0.15|^{2}\right)\left(1-|0.18|^{2}\right) 10^{1.8} \cdot 10^{2.2}}{\left(\frac{4 \pi 200 \lambda}{\lambda}\right)^{2}} \times 2 \\ &=2.99476 \times 10^{-3} \mathrm{W}=2.995 \mathrm{mW} \approx 3 \mathrm{mW} \end{aligned}
\]
\end{tcolorbox}

\vspace{1.5em}
\hrule
\vspace{1.5em}

\section*{Question 72: Basics of Electromagnetics \& Maxwell's Eqns (GATE EC 2016-SET-1)}
\addcontentsline{toc}{section}{Question 72: Basics of Electromagnetics \& Maxwell's Eqns (GATE EC 2016-SET-1)}
\noindent \textbf{\color{primary}MCQ} \hfill \textbf{\color{secondary}2 Marks}


\noindent The current density in a medium is given by 
 $\vec{j}=\frac{400 sin\theta }{2\pi (r^{2})+4}\hat{a}_{r}Am^{-2}$
 The total current and the average current density flowing through the portion of a spherical surface r = 0.8 m, $\frac{\pi }{12}\leq \theta \leq \frac{\pi }{4},0\leq \phi \leq 2\pi$ are given, respectively, by


\begin{enumerate}[label=(\Alph*)]
    \item \textbf{(Correct)} 15.09 A, 12.86 $Am^{-2}$
    \item 18.73 A, 13.65 $Am^{-2}$
    \item 12.86 A, 9.23 $Am^{-2}$
    \item 10.28 A, 7.56 $Am^{-2}$
\end{enumerate}

\noindent \textbf{Correct Answer:} (A)

\begin{tcolorbox}[solbox, title=Detailed Solution -- Question 72]
\[
\begin{aligned} I &=\int \overrightarrow{J} \cdot \overrightarrow{d s} \\ &=\int_{\theta=\frac{\pi}{12}}^{\pi / 4} \int_{\phi=0}^{2 \pi} \frac{400 \sin \theta}{2 \pi\left(r^{2}+4\right)} r^{2} \sin \theta d \theta d \phi \\ &=\left.\frac{400}{2 \pi\left(r^{2}+4\right)} \cdot r^{2} \cdot \phi\right|_{0} ^{2 \pi} \int_{\pi / 12}^{\pi / 4} \sin ^{2} \theta d \theta \\ &=\frac{400 r^{2}}{\left(r^{2}+4\right)} \int_{\pi / 12}^{\pi / 4}\left(\frac{1-\cos 2 \theta}{2}\right) d \theta \\ &=\frac{400 \cdot r^{2}}{\left(r^{2}+4\right)}\left(\frac{\left(\frac{\pi}{4}-\frac{\pi}{12}\right)}{2}-\left(\frac{\sin 2 \theta}{4}\right)_{\pi / 12}^{\pi / 4}\right) \\ &=\left.\frac{400 \cdot r^{2}}{\left(r^{2}+4\right)}\left(\frac{\pi}{12}-\left(\frac{1-1 / 2}{4}\right)\right)\right|_{r=0.8} \\ &=\frac{400 \times 0.8 \times 0.8}{4.64} \times 0.13=7.17 \mathrm{Amp} \end{aligned}
\]

\[
\begin{array}{l} \text { Total area }=\int d s=\iint r^{2} \sin \theta d \theta d \phi \\ =r^{2} \int_{0=\frac{\pi}{12}}^{\pi / 4} \sin \theta d \theta \cdot 2 \pi \\ =\left.r^{2} \cdot 2 \pi \cdot 0.259\right|_{r=0.8} \\ =0.8^{2} \times 0.5 \times 2 \pi \times \frac{1}{4} \\ =1.041 \mathrm{m}^{2} \\ \text { Average current }=\frac{7.17}{1.041}=6.88 \mathrm{A} / \mathrm{m}^{2} \end{array}
\]

 Note : GATE key mentioned (MTA - marks to all)
\end{tcolorbox}

\vspace{1.5em}
\hrule
\vspace{1.5em}

\section*{Question 73: Transmission Lines (GATE EC 2016-SET-1)}
\addcontentsline{toc}{section}{Question 73: Transmission Lines (GATE EC 2016-SET-1)}
\noindent \textbf{\color{primary}MCQ} \hfill \textbf{\color{secondary}1 Mark}


\noindent The propagation constant of a lossy transmission line is $(2 +j5) m^{-1}$
and its characteristic impedance is $(50 + j0) \Omega$ at $\omega = 10^{6}$ rad $s^{-1}$. The values of the line constants L, C,R, G are, respectively,


\begin{enumerate}[label=(\Alph*)]
    \item $L=200\mu H/m, C=0.1 \mu F/m,$ $R=50\Omega /m, G=0.02 S/m$
    \item \textbf{(Correct)} $L=250\mu H/m, C=0.1 \mu F/m,$ $R=100\Omega /m, G=0.04 S/m$
    \item $L=200\mu H/m, C=0.2 \mu F/m,$ $R=100\Omega /m, G=0.04 S/m$
    \item $L=250\mu H/m, C=0.2 \mu F/m,$ $R=50\Omega /m, G=0.04 S/m$
\end{enumerate}

\noindent \textbf{Correct Answer:} (B)

\begin{tcolorbox}[solbox, title=Detailed Solution -- Question 73]
\[
\begin{aligned} \gamma &=\sqrt{(R+j \omega L)(G+j \omega C)} \\ Z_{0} &=\sqrt{\frac{R+j \omega L}{G+j \omega C}} \\ \gamma \cdot Z_{0}=R+j \omega L =&(2+j 5)(50+j 0)=100+j 250 \\ R &=100 \Omega / m \\ L &=\frac{250}{\omega}=\frac{250}{10^{6}}=250 \mu \mathrm{H} / \mathrm{m} \\ \frac{\gamma}{Z_{0}}=\frac{2+j 5}{50}=G &+\dot{j \omega C}=0.04+j 0.1 \\ G &=0.04 \mathrm{S} / \mathrm{m} \\ C &=\frac{0.1}{\omega}=\frac{0.1}{10^{6}}=0.1 \mu \mathrm{F} / \mathrm{m} \end{aligned}
\]
\end{tcolorbox}

\vspace{1.5em}
\hrule
\vspace{1.5em}

\section*{Question 74: Basics of Electromagnetics \& Maxwell's Eqns (GATE EC 2016-SET-1)}
\addcontentsline{toc}{section}{Question 74: Basics of Electromagnetics \& Maxwell's Eqns (GATE EC 2016-SET-1)}
\noindent \textbf{\color{primary}NAT} \hfill \textbf{\color{secondary}1 Mark}


\noindent Concentric spherical shells of radii 2 m, 4 m, and 8 m carry uniform surface charge densities of 20 nC/$m^{2}$, -4 nC/$m^{2}$ and $\rho _{s}$ , respectively. The value of $\rho _{s}  (nC/m^{2})$ required to ensure that the
electric flux density $\overrightarrow{D}=\overrightarrow{0}$at radius 10 m is \underline{\hspace{1.5cm}}


\begin{tcolorbox}[solbox, title=Detailed Solution -- Question 74]
$\oint \vec{D} \cdot ds=Q \quad$ (charge enclosed) 
$Q_{1}+Q_{2}+Q_{3}=0$
 For D=0 

\[
\rho_{s 1} \cdot 4 \pi 2^{2}+\rho_{s 2} \cdot 4 \pi \cdot 4^{2}+\rho_{s 3} \cdot 4 \pi \cdot 8^{2}=Q=0
\]

$20 \cdot 4-4.4^{2}+\rho_{s 3} \cdot 8^{2}=0$
$80-64+\rho_{s 3} \cdot 8^{2}=0$
$\rho_{s 3}=\frac{-16}{64}=-0.25 \mathrm{nC} / \mathrm{m}^{2}$
\end{tcolorbox}

\vspace{1.5em}
\hrule
\vspace{1.5em}

\section*{Question 75: Transmission Lines (GATE EC 2015-SET-3)}
\addcontentsline{toc}{section}{Question 75: Transmission Lines (GATE EC 2015-SET-3)}
\noindent \textbf{\color{primary}NAT} \hfill \textbf{\color{secondary}2 Marks}


\noindent A coaxial capacitor of inner radius 1 mm and outer radius 5 mm has a capacitance per unit length of 172 pF/m. If the ratio of outer radius to inner radius is doubled, the capacitance per unit length (in pF/m) is \underline{\hspace{1.5cm}}.


\begin{tcolorbox}[solbox, title=Detailed Solution -- Question 75]
\[
\begin{aligned} C &=\frac{2 \pi \in}{ln \frac{b}{a}} \\ 172 \times 10^{-12} &=\frac{2 \pi \epsilon}{ln 5} \\ & \in \frac{172 \times 10^{-12} \times ln 5}{2 \pi} \end{aligned}
\]

 If ratio of $\frac{b}{a}$ is doubled then

\[
\begin{array}{l} C=\frac{2 \pi \epsilon}{\ln 10}=\frac{2 \pi \times \frac{172 \times 10^{-12} \times \ln 5}{2 \pi}}{\ln 10} \\ C=120.22 \mathrm{pF} \end{array}
\]
\end{tcolorbox}

\vspace{1.5em}
\hrule
\vspace{1.5em}

\section*{Question 76: Transmission Lines (GATE EC 2015-SET-3)}
\addcontentsline{toc}{section}{Question 76: Transmission Lines (GATE EC 2015-SET-3)}
\noindent \textbf{\color{primary}NAT} \hfill \textbf{\color{secondary}2 Marks}


\noindent A 200 m long transmission line having parameters shown in the figure is terminated into a load $R_{L}$. The line is connected to a 400 V source having source resistance $R_{S}$ through a switch, which is closed at t = 0. The transient response of the circuit at the input of the line (z = 0) is also drawn in the figure. The value of $R_{L}$ (in $\Omega$) is \underline{\hspace{1.5cm}}.

\begin{center}\includegraphics[max width=0.85\linewidth]{images/img\_41.jpg}\end{center}

\begin{center}\includegraphics[max width=0.85\linewidth]{images/img\_41.jpg}\end{center}


\begin{tcolorbox}[solbox, title=Detailed Solution -- Question 76]
\[
\begin{aligned} \text { At } t&=0 \\ &V_{\text {in }}=400 \times \frac{R_{0}}{R_{0}+R_{s}}=400 \times \frac{50}{200}=100 V \\ \text { At } t&=T_{0} \Gamma_{L} 100 \text { volts reflects }\\ \text { At } t&=2T_{0}\\ V_{\text {in }} &=100+\left(1+\Gamma_{\text {in }}\right) \Gamma_{L} 100=62.5 \\ \Gamma_{\text {in }} &=\frac{R_{s}-R_{o}}{R_{s}+R_{o}}=\frac{1}{2} \\ \text { Therefore, } \Gamma_{L} &=-\frac{1}{4}=\frac{R_{L}-50}{R_{L}+50} \\ R_{L} &=30 \Omega \end{aligned}
\]
\end{tcolorbox}

\vspace{1.5em}
\hrule
\vspace{1.5em}

\section*{Question 77: Transmission Lines (GATE EC 2015-SET-3)}
\addcontentsline{toc}{section}{Question 77: Transmission Lines (GATE EC 2015-SET-3)}
\noindent \textbf{\color{primary}MCQ} \hfill \textbf{\color{secondary}2 Marks}


\noindent Consider the 3 m long lossless air-filled transmission line shown in the figure. It has a characteristic impedance of 120$\pi \Omega$, is terminated by a short circuit, and is excited with a frequency of 37.5 MHz. What is the nature of the input impedance ($Z_{in}$)?

\begin{center}\includegraphics[max width=0.85\linewidth]{images/img\_42.jpg}\end{center}

\begin{center}\includegraphics[max width=0.85\linewidth]{images/img\_42.jpg}\end{center}


\begin{enumerate}[label=(\Alph*)]
    \item Open
    \item Short
    \item Inductive
    \item \textbf{(Correct)} Capacitive
\end{enumerate}

\noindent \textbf{Correct Answer:} (D)

\begin{tcolorbox}[solbox, title=Detailed Solution -- Question 77]
\[
\begin{aligned} l &=3 \mathrm{m} \\ f &=37.5 \mathrm{MHz} \\ \Rightarrow \quad \lambda &=\frac{3 \times 10^{8}}{37.5 \times 10^{6}}=8 \mathrm{m} \\ Z_{\text {in }}(S C) &=j Z_{0} \tan \beta l=j 120 \pi \tan \frac{2 \pi}{8} \cdot 3 \\ &=-j 120 \pi \Omega \end{aligned}
\]
\end{tcolorbox}

\vspace{1.5em}
\hrule
\vspace{1.5em}

\section*{Question 78: Basics of Electromagnetics \& Maxwell's Eqns (GATE EC 2015-SET-3)}
\addcontentsline{toc}{section}{Question 78: Basics of Electromagnetics \& Maxwell's Eqns (GATE EC 2015-SET-3)}
\noindent \textbf{\color{primary}NAT} \hfill \textbf{\color{secondary}2 Marks}


\noindent A vector field $D=2\rho ^{2} a_{\rho }+z a_{z}$ exists inside a cylindrical region enclosed by the surfaces $\rho$= 1, z = 0 and z = 5. Let S be the surface bounding this cylindrical region. The surface integral of this field on $S(\oint\oint_{S} D.ds)$  is \underline{\hspace{1.5cm}}.


\begin{tcolorbox}[solbox, title=Detailed Solution -- Question 78]
\[
\begin{array}{c} \vec{D}=2 \rho^{2} \hat{a}_{p}+z \hat{a}_{z} \\ \text { surface: } \rho=1, z=0 \text { and } z=5 \\ \oiint_{s} \vec{D} \cdot ds=\int_{\text {Top surface}} \vec{D} \cdot ds \\ +\int_{\text {Bottom surface}} \vec{D} \cdot ds \\ +\int_{\text {curved surface }} \vec{D} \cdot ds \\ =5 \pi+0+20 \pi=25 \pi=78.54 \end{array}
\]
\end{tcolorbox}

\vspace{1.5em}
\hrule
\vspace{1.5em}

\section*{Question 79: Transmission Lines (GATE EC 2015-SET-3)}
\addcontentsline{toc}{section}{Question 79: Transmission Lines (GATE EC 2015-SET-3)}
\noindent \textbf{\color{primary}MCQ} \hfill \textbf{\color{secondary}1 Mark}


\noindent A coaxial cable is made of two brass conductors. The spacing between the conductors is filled with Teflon ( $\varepsilon^{'}_{r}$= 2.1, tan $\delta$=0). Which one of the following circuits can represent the lumped element model of a small piece of this cable having length $\Delta$z? 

\begin{center}\includegraphics[max width=0.85\linewidth]{images/img\_31.jpg}\end{center}

\begin{center}\includegraphics[max width=0.85\linewidth]{images/img\_31.jpg}\end{center}


\begin{enumerate}[label=(\Alph*)]
    \item A
    \item \textbf{(Correct)} B
    \item C
    \item D
\end{enumerate}

\noindent \textbf{Correct Answer:} (B)

\begin{tcolorbox}[solbox, title=Detailed Solution -- Question 79]
The teflon dielectric between the lines has
$\tan \delta=\text{Loss tangent}=0$
 For the medium between the lines
$G=0$
\end{tcolorbox}

\vspace{1.5em}
\hrule
\vspace{1.5em}

\section*{Question 80: Antennas \& Radiating Systems (GATE EC 2015-SET-3)}
\addcontentsline{toc}{section}{Question 80: Antennas \& Radiating Systems (GATE EC 2015-SET-3)}
\noindent \textbf{\color{primary}MCQ} \hfill \textbf{\color{secondary}1 Mark}


\noindent The directivity of an antenna array can be increased by adding more antenna elements, as a larger number of elements


\begin{enumerate}[label=(\Alph*)]
    \item improves the radiation efficiency
    \item \textbf{(Correct)} increases the effective area of the antenna
    \item results in a better impedance matching
    \item allows more power to be transmitted by the antenna
\end{enumerate}

\noindent \textbf{Correct Answer:} (B)

\begin{tcolorbox}[solbox, title=Detailed Solution -- Question 80]
Directivity (0) is increased by adding more antenna elements in an antenna array. Effective area $(A_{e})$ and directivity (D) are related by 
$A_{e}=\frac{\lambda^{2}}{4 \pi} D$
$\Rightarrow D \uparrow A_{e} \uparrow$
\end{tcolorbox}

\vspace{1.5em}
\hrule
\vspace{1.5em}

\section*{Question 81: Uniform Plane Waves (GATE EC 2015-SET-2)}
\addcontentsline{toc}{section}{Question 81: Uniform Plane Waves (GATE EC 2015-SET-2)}
\noindent \textbf{\color{primary}MCQ} \hfill \textbf{\color{secondary}2 Marks}


\noindent The electric field of a plane wave propagating in a lossless non-magnetic medium is given by the following expression 
$E(z,t)=a_{x}5 cos(2\pi \times 10^{9}t+\beta z)+$ $a_{y} 3 cos(2\pi \times 10^{9}t+\beta z-\pi/2)$ 
 The type of the polarization is


\begin{enumerate}[label=(\Alph*)]
    \item Right Hand Circular.
    \item \textbf{(Correct)} Left Hand Elliptical
    \item Right Hand Elliptical
    \item Linear.
\end{enumerate}

\noindent \textbf{Correct Answer:} (B)

\begin{tcolorbox}[solbox, title=Detailed Solution -- Question 81]
$\vec{E}(z, t)=\hat{a}_{x} 5 \cos \left(2 \pi \times 10^{9} t+\beta z\right)$
$+\hat{a}_{y} 3 \cos \left(2 \pi \times 10^{9} t+\beta z-\frac{\pi}{2}\right)$
$\Rightarrow$ Wave is travelling in $-\hat{a}_{z}$ direction.
$\Rightarrow$ Wave has orthogonal components with unequal amplitudes and checking the time trace.
$\therefore$ Wave is left hand elliptically polarized.
\end{tcolorbox}

\vspace{1.5em}
\hrule
\vspace{1.5em}

\section*{Question 82: Antennas \& Radiating Systems (GATE EC 2015-SET-2)}
\addcontentsline{toc}{section}{Question 82: Antennas \& Radiating Systems (GATE EC 2015-SET-2)}
\noindent \textbf{\color{primary}NAT} \hfill \textbf{\color{secondary}2 Marks}


\noindent Two half-wave dipole antennas placed as shown in the figure are excited with sinusoidally varying currents of frequency 3 MHz and phase shift of $\pi/2$ between them(the element at the origin leads in phase). If the maximum radiated E-field at the point P in the x-y plane occurs at an azimuthal angle of $60^{\circ}$, the distance d (in meters) between the antennas is \underline{\hspace{1.5cm}}. 

\begin{center}\includegraphics[max width=0.85\linewidth]{images/img\_5.jpg}\end{center}

\begin{center}\includegraphics[max width=0.85\linewidth]{images/img\_5.jpg}\end{center}


\begin{tcolorbox}[solbox, title=Detailed Solution -- Question 82]
Both dipole antennas have isotropic pattern in $\phi$ view or H plane view.

\begin{center}\includegraphics[max width=0.85\linewidth]{images/img\_55.jpg}\end{center}

\begin{center}\includegraphics[max width=0.85\linewidth]{images/img\_55.jpg}\end{center}

$\psi=\beta d \cos \theta+\alpha$
 where, 
\[
\alpha=-90^{\circ}, \beta=\frac{2 \pi}{\lambda}=\frac{2 \pi}{\frac{3 \times 10^{8}}{3 \times 10^{6}}}=\frac{2 \pi}{100}
\]

 Since maximum radiation $(\psi=0)$ is at $\theta=60^{\circ}$

\[
\begin{array}{l} \Rightarrow \quad 0=\frac{2 \pi}{100} d \cos 60^{\circ}-90^{\circ} \\ \Rightarrow \quad 90=\frac{360}{100} d \frac{1}{2} \\ \Rightarrow \quad d=50 \mathrm{m} \end{array}
\]
\end{tcolorbox}

\vspace{1.5em}
\hrule
\vspace{1.5em}

\section*{Question 83: Waveguides \& Optical Fibers (GATE EC 2015-SET-2)}
\addcontentsline{toc}{section}{Question 83: Waveguides \& Optical Fibers (GATE EC 2015-SET-2)}
\noindent \textbf{\color{primary}NAT} \hfill \textbf{\color{secondary}2 Marks}


\noindent An air-filled rectangular waveguide of internal dimensions $a \; cm \times b \;  cm \; (a > b)$ has a cutoff frequency of 6 GHz for the dominant $TE_{10}$ mode. For the same waveguide, if the cutoff frequencyof the $TM_{11}$ mode is 15 GHz, the cutoff frequency of the $TE_{01}$ mode in GHz is \underline{\hspace{1.5cm}}.


\begin{tcolorbox}[solbox, title=Detailed Solution -- Question 83]
Given, $f_{c} T E_{10}$ mode is $6 \mathrm{GHz}$

\[
\Rightarrow \quad \frac{c}{2 a}=\frac{3 \times 10^{10}}{2 \times a}=6 \times 10^{9}
\]

$\Rightarrow a=2.5 \mathrm{cm}$
 Given $f_{c}$ for $T M_{11}$ mode is $15 \mathrm{GHz}$

\[
\begin{array}{l} \Rightarrow \frac{c}{2} \sqrt{\frac{1}{a^{2}}+\frac{1}{b^{2}}}=15 \times 10^{9} \\ \begin{array}{l} \frac{3 \times 10^{10}}{2} \sqrt{\frac{1}{6.25}+\frac{1}{b^{2}}}=15 \times 10^{9} \\ \Rightarrow \quad b=\frac{2.5}{\sqrt{5.25}} \end{array} \\ \begin{aligned} \Rightarrow \quad & \quad 2 \end{aligned} \end{array}
\]

 Now $f_{c}$ for $T E_{01}$ mode is,

\[
\begin{aligned} f_{c} &=\frac{c}{2 b}=\frac{3 \times 10^{10}}{2.5} \sqrt{5.25} \\ &=13.75 \mathrm{GHz} \end{aligned}
\]
\end{tcolorbox}

\vspace{1.5em}
\hrule
\vspace{1.5em}

\section*{Question 84: Uniform Plane Waves (GATE EC 2015-SET-2)}
\addcontentsline{toc}{section}{Question 84: Uniform Plane Waves (GATE EC 2015-SET-2)}
\noindent \textbf{\color{primary}MCQ} \hfill \textbf{\color{secondary}1 Mark}


\noindent The electric field of a uniform plane electromagnetic wave is 
 $\vec{E}=(\vec{a}_{x}+j4\vec{a}_{y})exp[j(2\pi \times 10^{7}t-0.2z)]$  
 The polarization of the wave is


\begin{enumerate}[label=(\Alph*)]
    \item right handed circular
    \item left handed circular
    \item right handed elliptical
    \item \textbf{(Correct)} left handed elliptical
\end{enumerate}

\noindent \textbf{Correct Answer:} (D)

\begin{tcolorbox}[solbox, title=Detailed Solution -- Question 84]
\[
\begin{aligned} \vec{E}(z, t)=& \hat{a}_{x} 5 \cos \left(2 \pi \times 10^{9} t+\beta z\right) \\ &+\hat{a}_{y} 3 \cos \left(2 \pi \times 10^{9} t+\beta z-\frac{\pi}{2}\right) \end{aligned}
\]

$\Rightarrow$ Wave is travelling in $-\hat{a}_{z}$ direction.
$\Rightarrow$ Wave has orthogonal components with unequal  amplitudes and checking the time trace.
$\therefore$ Wave is left hand elliptically polarized.
\end{tcolorbox}

\vspace{1.5em}
\hrule
\vspace{1.5em}

\section*{Question 85: Basics of Electromagnetics \& Maxwell's Eqns (GATE EC 2015-SET-2)}
\addcontentsline{toc}{section}{Question 85: Basics of Electromagnetics \& Maxwell's Eqns (GATE EC 2015-SET-2)}
\noindent \textbf{\color{primary}NAT} \hfill \textbf{\color{secondary}1 Mark}


\noindent In a source free region in vacuum, if the electrostatic potential $\varphi = 2x^{2}+y^{2}+cz^{2}$, the value of constant c must be \underline{\hspace{1.5cm}}.


\begin{tcolorbox}[solbox, title=Detailed Solution -- Question 85]
$\phi=2 x^{2}+y^{2}+c z^{2}$
 In source-free region, 

\[
\begin{aligned} \nabla \cdot \vec{D}&=0 \\ \Rightarrow \quad \in(\nabla \cdot \vec{E}) &=0 \\ \Rightarrow \quad \nabla \cdot \vec{E} &=0\\ where,\quad\vec{E}&=-\nabla V=-\nabla \phi \\ Now, \quad \nabla V&=\frac{\partial \phi}{\partial x} \hat{a}_{x}+\frac{\partial \phi}{\partial y} \hat{a}_{y}+\frac{\partial \phi}{\partial z} \hat{a}_{z} \\ \Rightarrow \quad -\vec{E}&=4 x \hat{a}_{x}+2 y \hat{a}_{y}+2 c z \hat{a}_{z} \\ \text{Again }\nabla \cdot \vec{E}&=0 \\ \Rightarrow \frac{\partial E_{x}}{\partial x}+\frac{\partial E_{y}}{\partial y}&+\frac{\partial E_{t}}{\partial z}=4+2+2 c=0 \\ c&=-3 \end{aligned}
\]
\end{tcolorbox}

\vspace{1.5em}
\hrule
\vspace{1.5em}

\section*{Question 86: Uniform Plane Waves (GATE EC 2015-SET-1)}
\addcontentsline{toc}{section}{Question 86: Uniform Plane Waves (GATE EC 2015-SET-1)}
\noindent \textbf{\color{primary}NAT} \hfill \textbf{\color{secondary}2 Marks}


\noindent Consider a uniform plane wave with amplitude $(E_{0})$ of 10 V/m and 1.1 GHz frequency travelling in air, and incident normally on a dielectric medium with complex relative permittivity ($\varepsilon_{r}$) and permeability ($\mu_{r}$) as shown in the figure.

\begin{center}\includegraphics[max width=0.85\linewidth]{images/img\_54.jpg}\end{center}

\begin{center}\includegraphics[max width=0.85\linewidth]{images/img\_54.jpg}\end{center}

The magnitude of the transmitted electric field component (in V/m) after it has travelled a distance of 10 cm inside the dielectric region is\underline{\hspace{1.5cm}}.


\begin{tcolorbox}[solbox, title=Detailed Solution -- Question 86]
Material has $\mu_{R}$ and $\epsilon_{R}$ complex.

\[
\eta \text { of material }=\sqrt{\frac{j \omega \mu}{\sigma+j \omega \epsilon}}=\sqrt{\frac{\mu_{o} \mu_{R}}{\epsilon_{o} \epsilon_{R}}}=120 \pi
\]

 where, $\sigma=0 \quad$ (dielectric) 
 No reflections at the boundary.

\[
\begin{aligned} \gamma &=\text { Propagation constant } \\ &=\sqrt{j \omega \mu_{o} \mu_{R}\left(0+j \omega \epsilon_{o} \epsilon_{R}\right)} \\ &=\sqrt{j \omega \mu_{o} \epsilon_{o}(1-j 2)} \\ &=j \omega \sqrt{\mu_{o} \epsilon_{o}}+j 2 \omega \sqrt{\mu_{o} \epsilon_{o}} \\ \alpha &=2 \omega \sqrt{\mu_{o} \epsilon_{o}}=4.6 \frac{1}{\mathrm{m}} \\ |E| &=E_{o} e^{-\alpha z} \text { with } Z=0.1 \mathrm{m} \\ |E| &=0.1 \mathrm{V} / \mathrm{m} \end{aligned}
\]
\end{tcolorbox}

\vspace{1.5em}
\hrule
\vspace{1.5em}

\section*{Question 87: Uniform Plane Waves (GATE EC 2015-SET-1)}
\addcontentsline{toc}{section}{Question 87: Uniform Plane Waves (GATE EC 2015-SET-1)}
\noindent \textbf{\color{primary}NAT} \hfill \textbf{\color{secondary}2 Marks}


\noindent The electric field intensity of a plane wave traveling in free space is given by the following
expression 
 $E(x,t)=a_{y}24\pi cos(\omega t-k_{0}x)(V/m)$  
 In this field, consider a square area 10 cm x 10 cm on a plane x + y = 1. The total time-averaged power (in mW) passing through the square area is \underline{\hspace{1.5cm}}.


\begin{tcolorbox}[solbox, title=Detailed Solution -- Question 87]
$\vec{E}(x, t)=\hat{a}_{y} 24 \pi \cos \left(\omega t-K_{o} x\right) V / m$
 Power density

\[
\vec{P}=\frac{1}{2} \frac{|E|^{2}}{\eta}=\frac{1}{2} \frac{(24 \pi)^{2}}{120 \pi} \hat{a}_{x}
\]

 Time-averaged power

\[
\begin{aligned} P_{\mathrm{avg}} &=\vec{P} \cdot \overrightarrow{d s} \\ &=\left[\frac{1}{2} \frac{(24 \pi)^{2}}{120 \pi} \hat{a}_{x}\right] \cdot\left[\frac{\hat{a}_{x}+\hat{a}_{y}}{\sqrt{2}} \times 10^{-2}\right] \\ P_{\mathrm{avg}} &=\frac{24 \pi}{\sqrt{2}}=53.29 \mathrm{mW} \end{aligned}
\]
\end{tcolorbox}

\vspace{1.5em}
\hrule
\vspace{1.5em}

\section*{Question 88: Waveguides \& Optical Fibers (GATE EC 2015-SET-1)}
\addcontentsline{toc}{section}{Question 88: Waveguides \& Optical Fibers (GATE EC 2015-SET-1)}
\noindent \textbf{\color{primary}MCQ} \hfill \textbf{\color{secondary}2 Marks}


\noindent The longitudinal component of the magnetic field inside an air-filled rectangular waveguide made of a perfect electric conductor is given by the following expression 
 $H_{z}(x,y,z,t)=0.1 cos(25\pi x) cos(30.3 \pi y)$ $cos(12 \pi *10^{9}t-\beta z) (A/m)$ 
 The cross-sectional dimensions of the waveguide are given as a = 0.08 m and b = 0.033 m. The
mode of propagation inside the waveguide is


\begin{enumerate}[label=(\Alph*)]
    \item $TM_{12}$
    \item $TM_{21}$
    \item \textbf{(Correct)} $TE_{21}$
    \item $TE_{12}$
\end{enumerate}

\noindent \textbf{Correct Answer:} (C)

\begin{tcolorbox}[solbox, title=Detailed Solution -- Question 88]
\[
\begin{aligned} H_{z}(x, y, z, t) &=0.1 \cos (25 \pi x) \cos (30.3 \pi y) \\ & \cos \left(12 \pi \times 10^{9} t-\beta z\right)^{-1 / m}\\ a &=0.08 m, b=0.033 m \end{aligned}
\]

 Axial component is $\mathrm{H} \Rightarrow$ the propagating mode is 
$T E_{\text {mn}}$, m, n can be found by

\[
\begin{array}{l} \quad \frac{m \pi}{a} x=25 \pi x \\ \Rightarrow \quad \frac{m}{0.08}=25 \\ \Rightarrow \quad m=2 \\ \Rightarrow \quad \frac{n\pi}{b} y=30.3 \pi y \\ \Rightarrow \quad \frac{n}{0.033}=30.3 \\ \Rightarrow \quad n=1 \end{array}
\]

$\therefore$ The mode of propagation is $T E_{21}$
\end{tcolorbox}

\vspace{1.5em}
\hrule
\vspace{1.5em}

\section*{Question 89: Uniform Plane Waves (GATE EC 2015-SET-1)}
\addcontentsline{toc}{section}{Question 89: Uniform Plane Waves (GATE EC 2015-SET-1)}
\noindent \textbf{\color{primary}NAT} \hfill \textbf{\color{secondary}1 Mark}


\noindent The electric field component of a plane wave traveling in a lossless dielectric medium is given by $\vec{E}(z,t)=\hat{a_{y}}2cos(10^{8}t-\frac{z}{\sqrt{2}})V/m$. The wavelength (in m) for the wave is \underline{\hspace{1.5cm}}.


\begin{tcolorbox}[solbox, title=Detailed Solution -- Question 89]
\[
\vec{E}(z, t)=\hat{a}_{y} 2 \cos \left(10^{8} t-\frac{z}{\sqrt{2}}\right) \mathrm{V} / \mathrm{m}
\]

 General form of 
$\vec{E}(z, t)=\hat{a}_{y} E_{o} \cos (\omega t-\beta z) \vee / m$
 Comparing both, we get $\beta=\frac{1}{\sqrt{2}}$
 since, $\beta=\frac{2 \pi}{\lambda}$

\[
\Rightarrow \quad \lambda=\frac{2 \pi}{\beta}=\frac{2 \pi}{\frac{1}{\sqrt{2}}}=2 \sqrt{2} \pi m
\]

$\Rightarrow \quad \lambda=8.88 \mathrm{m}$
\end{tcolorbox}

\vspace{1.5em}
\hrule
\vspace{1.5em}

\section*{Question 90: Basics of Electromagnetics \& Maxwell's Eqns (GATE EC 2015-SET-1)}
\addcontentsline{toc}{section}{Question 90: Basics of Electromagnetics \& Maxwell's Eqns (GATE EC 2015-SET-1)}
\noindent \textbf{\color{primary}MCQ} \hfill \textbf{\color{secondary}1 Mark}


\noindent Consider a straight, infinitely long, current carrying conductor lying on the z-axis. Which one of the following plots (in linear scale) qualitatively represents the dependence of $H_{\phi}$ on r, where $H_{\phi}$ is the magnitude of the azimuthal component of magnetic field outside the conductor and r is the radial distance from the conductor?

\begin{center}\includegraphics[max width=0.85\linewidth]{images/img\_49.jpg}\end{center}

\begin{center}\includegraphics[max width=0.85\linewidth]{images/img\_49.jpg}\end{center}


\begin{enumerate}[label=(\Alph*)]
    \item A
    \item B
    \item \textbf{(Correct)} C
    \item D
\end{enumerate}

\noindent \textbf{Correct Answer:} (C)

\begin{tcolorbox}[solbox, title=Detailed Solution -- Question 90]
$H_{\phi}=\frac{I}{2 \pi r} \hat{a}_{\phi}$
$\Rightarrow \quad H_{\phi} \propto \frac{1}{r}$option(c) is satisfied.
\end{tcolorbox}

\vspace{1.5em}
\hrule
\vspace{1.5em}

\section*{Question 91: Basics of Electromagnetics \& Maxwell's Eqns (GATE EC 2014-SET-4)}
\addcontentsline{toc}{section}{Question 91: Basics of Electromagnetics \& Maxwell's Eqns (GATE EC 2014-SET-4)}
\noindent \textbf{\color{primary}MCQ} \hfill \textbf{\color{secondary}2 Marks}


\noindent If $\vec{E}=-9(2y^{3}-3yz^{2})\hat{x}-(6xy^{2}-3xz^{2})\hat{y}+(6xyz)\hat{z}$ is the electric field in a source free region, a valid expression for the electrostatic potential is


\begin{enumerate}[label=(\Alph*)]
    \item $xy^{3}-yz^{2}$
    \item $2xy^{3}-xyz^{2}$
    \item $y^{3}+xyz^{2}$
    \item \textbf{(Correct)} $2xy^{3}-3xyz^{2}$
\end{enumerate}

\noindent \textbf{Correct Answer:} (D)

\begin{tcolorbox}[solbox, title=Detailed Solution -- Question 91]
\[
\begin{aligned} E &=-\frac{\partial V}{\partial x} \hat{a}_{x}-\frac{\partial V}{\partial y} \hat{a}_{y}-\frac{\partial V}{\partial Z} \hat{a}_{z} \\ \frac{\partial V}{\partial x} &=2 y^{3}-3 y z^{2} &V=2 x y^{3}-3 x y z^{2}+K_{1} \\ \frac{\partial V}{\partial y} &=6 x y^{2}-3 x z^{2} &V=2 x y^{3}-3 x y z^{2}+K_{2} \\ \frac{\partial V}{\partial z} &=-6 x y z &V=-3 x y z^{2}+K_{3} \end{aligned}
\]

 Comparing all three for uniqueness solution of potential

\[
\begin{array}{c} K_{1}=K_{2}=0, \quad K_{3}=2 x y^{3} \\ \vec{E}=-\left(2 y^{3}-3 y z^{2}\right) \hat{x}-\left(6 x y^{2}-3 x z^{2}\right) \hat{y}+(6 x y z) \hat{z} \end{array}
\]

 as we know that $E=-\nabla V$

\[
E=-\left(\frac{\partial}{\partial x} \hat{x}+\frac{\partial}{\partial y} \hat{y}+\frac{\partial}{\partial z} \hat{z}\right) V \quad\ldots(i)
\]

 Clearly option (D) satisfies the given relation.
\end{tcolorbox}

\vspace{1.5em}
\hrule
\vspace{1.5em}

\section*{Question 92: Basics of Electromagnetics \& Maxwell's Eqns (GATE EC 2014-SET-4)}
\addcontentsline{toc}{section}{Question 92: Basics of Electromagnetics \& Maxwell's Eqns (GATE EC 2014-SET-4)}
\noindent \textbf{\color{primary}NAT} \hfill \textbf{\color{secondary}2 Marks}


\noindent The electric field (assumed to be one-dimensional) between two points A and B is shown. Let $\psi _{A}$ and $\psi _{B}$ be the electrostatic potentials at A and B, respectively. The value of $\psi _{B}-\psi _{A}$ in Volts is 

\begin{center}\includegraphics[max width=0.85\linewidth]{images/img\_64.jpg}\end{center}

\begin{center}\includegraphics[max width=0.85\linewidth]{images/img\_64.jpg}\end{center}


\begin{tcolorbox}[solbox, title=Detailed Solution -- Question 92]
\begin{center}\includegraphics[max width=0.85\linewidth]{images/img\_52.jpg}\end{center}

\begin{center}\includegraphics[max width=0.85\linewidth]{images/img\_52.jpg}\end{center}

E is linear from A to B and E=mx+C 

\[
\begin{aligned} V&=-\int_{B}^{A}(m x+c) d x \\ \text { At point } A: x&= 0, E=c=2 \times 10^{6} \vee / m =x 10^{-6}\\ \text { At point } B: x &=5 \times 10^{-6} \\ E &=4 \times 10^{6}=m \times 5 \times 10^{-6}+2 \times 10^{6} \\ m &=4 \times 10^{11}\\ -\int^{\Psi_{A}}_{\Psi_{B}} E \cdot d l&=-\int_{5\mu m}^{0}\left[20 \times 10^{5} \times 5 \times 10^{-6}+\frac{4 \times 10^{11}}{2}\left[\left(5 \times 10^{-6}\right)-0^{2}\right]\right] \\ \Psi_{A}-\Psi_{B} &=-\left[\left(20 \times 10^{5} \times 5 \times 10^{-6}\right)+\frac{4 \times 10^{11}}{2}\left[\left(5 \times 10^{-6}\right)^{2}-0^{2}\right]\right] \\ &=-\left[10+\left(2 \times 10^{11} \times 25 \times 10^{-12}\right)\right] \\ &=-15 \mathrm{V} \end{aligned}
\]
\end{tcolorbox}

\vspace{1.5em}
\hrule
\vspace{1.5em}

\section*{Question 93: Antennas \& Radiating Systems (GATE EC 2014-SET-4)}
\addcontentsline{toc}{section}{Question 93: Antennas \& Radiating Systems (GATE EC 2014-SET-4)}
\noindent \textbf{\color{primary}MCQ} \hfill \textbf{\color{secondary}1 Mark}


\noindent Match column A with column B. 

\begin{center}\includegraphics[max width=0.85\linewidth]{images/img\_22.jpg}\end{center}

\begin{center}\includegraphics[max width=0.85\linewidth]{images/img\_22.jpg}\end{center}


\begin{enumerate}[label=(\Alph*)]
    \item $1\rightarrow P  \;  2\rightarrow Q  \; 3\rightarrow R$
    \item \textbf{(Correct)} $1\rightarrow R  \; 2\rightarrow P  \; 3\rightarrow Q$
    \item $1\rightarrow Q  \;  2\rightarrow P  \; 3\rightarrow R$
    \item $1\rightarrow R  \; 2\rightarrow Q  \; 3\rightarrow P$
\end{enumerate}

\noindent \textbf{Correct Answer:} (B)

\begin{tcolorbox}[solbox, title=Detailed Solution -- Question 93]
Point electromagnetic source radiates in all
direction
Dish antenna radiates any electromagnetic
energy in any particular direction with narrow
beam-width and high directivity
Yagi-uda antenna is high gain antenna used
for T.V. reception. Its radiation is along the axis
of the antenna.

\begin{center}\includegraphics[max width=0.85\linewidth]{images/img\_1.jpg}\end{center}

\begin{center}\includegraphics[max width=0.85\linewidth]{images/img\_1.jpg}\end{center}
\end{tcolorbox}

\vspace{1.5em}
\hrule
\vspace{1.5em}

\section*{Question 94: Antennas \& Radiating Systems (GATE EC 2014-SET-4)}
\addcontentsline{toc}{section}{Question 94: Antennas \& Radiating Systems (GATE EC 2014-SET-4)}
\noindent \textbf{\color{primary}NAT} \hfill \textbf{\color{secondary}1 Mark}


\noindent For an antenna radiating in free space, the electric field at a distance of 1 km is found to be 12mV/m. Given that intrinsic impedance of the free space is $120\pi \Omega$, the magnitude of average power density due to this antenna at a distance of 2 km from the antenna (in nW/$m^{2}$) is \underline{\hspace{1.5cm}}.


\begin{tcolorbox}[solbox, title=Detailed Solution -- Question 94]
Given electric field, E at a distance of
$1 \mathrm{km}=12 \mathrm{mV} / \mathrm{m}$
 Also we know that $E \propto \frac{1}{r}$, where r is the distance 
 where electric field is to be measured. 
 So, electric field, E at a distance of $2 \mathrm{km}$
$=\frac{12 \mathrm{mV} / \mathrm{m}}{2}=6 \mathrm{mV} / \mathrm{m}$
 Also, power density due to the antenna is given as,

\[
\begin{aligned} P_{\text {avg }} &=\frac{1}{2} \frac{E^{2}}{\eta}=\frac{1}{2} \times \frac{6 \times 6 \times 10^{-6}}{120 \pi} \\ &=47.7 \mathrm{nW} / \mathrm{m}^{2} \end{aligned}
\]
\end{tcolorbox}

\vspace{1.5em}
\hrule
\vspace{1.5em}

\section*{Question 95: Uniform Plane Waves (GATE EC 2014-SET-3)}
\addcontentsline{toc}{section}{Question 95: Uniform Plane Waves (GATE EC 2014-SET-3)}
\noindent \textbf{\color{primary}NAT} \hfill \textbf{\color{secondary}2 Marks}


\noindent Assume that a plane wave in air with an electric field $\vec{E}=10cos(\omega t-3x-\sqrt{3}z)\hat{a}_{y} V/m$ is incident on a non-magnetic dielectric
slab of relative permittivity 3 which covers the region $z > 0$. The angle of
transmission in the dielectric slab is \underline{\hspace{1.5cm}}degrees.


\begin{tcolorbox}[solbox, title=Detailed Solution -- Question 95]
\begin{center}\includegraphics[max width=0.85\linewidth]{images/img\_75.jpg}\end{center}

\begin{center}\includegraphics[max width=0.85\linewidth]{images/img\_75.jpg}\end{center}

\[
\begin{aligned} \beta_{x}&=3, \beta_{z}=\sqrt{3}\text{ from }E_{i} \text{ of the wave}\\ \theta_{i} &=60^{\circ} \\ \frac{\sin 60^{\circ}}{\sin \theta_{t}} &=\sqrt{3} \\ \theta_{t} &=30^{\circ} \end{aligned}
\]
\end{tcolorbox}

\vspace{1.5em}
\hrule
\vspace{1.5em}

\section*{Question 96: Basics of Electromagnetics \& Maxwell's Eqns (GATE EC 2014-SET-3)}
\addcontentsline{toc}{section}{Question 96: Basics of Electromagnetics \& Maxwell's Eqns (GATE EC 2014-SET-3)}
\noindent \textbf{\color{primary}MCQ} \hfill \textbf{\color{secondary}2 Marks}


\noindent A region shown below contains a perfect conducting half-space and air. The
surface current $\vec{K}_{s}$ on the surface of the perfect conductor is $\hat{K}_{s}=\hat{x}_{2}$ amperes per meter. The tangential $\vec{H}$ field in the air just above the perfect conductor is

\begin{center}\includegraphics[max width=0.85\linewidth]{images/img\_48.jpg}\end{center}

\begin{center}\includegraphics[max width=0.85\linewidth]{images/img\_48.jpg}\end{center}


\begin{enumerate}[label=(\Alph*)]
    \item $(\hat{x}+\hat{z})2$ amperes per meter
    \item $\hat{x}2$ amperes per meter
    \item $-\hat{z}2$ amperes per meter
    \item \textbf{(Correct)} $\hat{z}2$ amperes per meter
\end{enumerate}

\noindent \textbf{Correct Answer:} (D)

\begin{tcolorbox}[solbox, title=Detailed Solution -- Question 96]
since, we know from the property of magnetic boundary conditions i.e.
$\vec{H}_{t 2}-\vec{H}_{t 1}=\hat{K}_{s} \times a_{N}$
 In the given question the tangential field $\vec{H}$ in the air just above the perfect conductor is asked with H=0 inside the conductor

\[
\begin{array}{l} \vec{H}_{t 1}=0 \\ \vec{H}_{t 2}=\hat{x} 2 \times \hat{y} \\ \vec{H}_{t 2}=\hat{z} 2 \mathrm{A} / \mathrm{m} \end{array}
\]
\end{tcolorbox}

\vspace{1.5em}
\hrule
\vspace{1.5em}

\section*{Question 97: Transmission Lines (GATE EC 2014-SET-3)}
\addcontentsline{toc}{section}{Question 97: Transmission Lines (GATE EC 2014-SET-3)}
\noindent \textbf{\color{primary}MCQ} \hfill \textbf{\color{secondary}1 Mark}


\noindent In the following figure, the transmitter $T_{X}$
sends a wideband modulated RF signal via a coaxial cable to the receiver $R_{X}$. The output impedance $Z_{T}$ of $T_{X}$,
the characteristic impedance $Z_{0}$ of the cable and the input impedance $Z_{R}$ of $R_{X}$ are all real. 

\begin{center}\includegraphics[max width=0.85\linewidth]{images/img\_53.jpg}\end{center}

\begin{center}\includegraphics[max width=0.85\linewidth]{images/img\_53.jpg}\end{center}


\begin{enumerate}[label=(\Alph*)]
    \item The signal gets distorted if $Z_{R}\neq Z_{0}$, irrespective of the value of $Z_{T}$
    \item The signal gets distorted if $Z_{T}\neq Z_{0}$, irrespective of the value of $Z_{R}$
    \item \textbf{(Correct)} Signal distortion implies impedance mismatch at both ends :$Z_{T}\neq Z_{0}$ and $Z_{R}\neq Z_{0}$
    \item Impedance mismatches do NOT result in signal distortion but reduce power
transfer efficiency
\end{enumerate}

\noindent \textbf{Correct Answer:} (C)

\begin{tcolorbox}[solbox, title=Detailed Solution -- Question 97]
The signal wave gets distorted if there is any
impedance mismatch between the characteristic
impedance $Z_{0}$ and the transmitting impedance $Z_{T}$
and $Z_{0}$ and receiving impedance $Z_{R}$ 

i.e. $Z_{T} \neq Z_{0}$ and $Z_{R} \neq Z_{0}$

then only distortion occurs. 

As standing waveforms the signal has multi
harmonics and hence distortion occurs.
\end{tcolorbox}

\vspace{1.5em}
\hrule
\vspace{1.5em}

\section*{Question 98: Waveguides \& Optical Fibers (GATE EC 2014-SET-3)}
\addcontentsline{toc}{section}{Question 98: Waveguides \& Optical Fibers (GATE EC 2014-SET-3)}
\noindent \textbf{\color{primary}NAT} \hfill \textbf{\color{secondary}1 Mark}


\noindent Consider an air filled rectangular waveguide with a cross-section of 5cm x 3cm. For this waveguide, the cut-off frequency (in MHz) of $TE_{21}$ mode is \underline{\hspace{1.5cm}}.


\begin{tcolorbox}[solbox, title=Detailed Solution -- Question 98]
Given $a=5 \mathrm{cm}=5,10^{-2} \mathrm{m}$

\[
\begin{array}{c} b=3 \mathrm{cm}=3 \times 10^{-2} \mathrm{m} \\ f_{c}=\frac{C}{2 \pi} \sqrt{\left(\frac{n m}{a}\right)^{2}+\left(\frac{r m}{b}\right)^{2}} \\ f_{C T E_{21}}=\frac{C}{2} \sqrt{\left(\frac{2}{5 \mathrm{cm}}\right)^{2}+\left(\frac{1}{3 \mathrm{cm}}\right)^{2}} \\ =\frac{3 \times 10^{8}}{2} \sqrt{\left(\frac{2 \times 100}{5}\right)^{2}+\left(\frac{100}{3}\right)^{2}} \\ f_{c}=7810 \mathrm{MHz} \end{array}
\]
\end{tcolorbox}

\vspace{1.5em}
\hrule
\vspace{1.5em}

\section*{Question 99: Waveguides \& Optical Fibers (GATE EC 2014-SET-2)}
\addcontentsline{toc}{section}{Question 99: Waveguides \& Optical Fibers (GATE EC 2014-SET-2)}
\noindent \textbf{\color{primary}NAT} \hfill \textbf{\color{secondary}2 Marks}


\noindent For a rectangular waveguide of internal dimensions $a \times b(a > b)$, the cur-off
frequency for the $TE_{11}$ mode is the arithmetic mean of the cut-off frequencies for $TE_{10}$ mode and $TE_{20}$ mode. If a = $\sqrt{5}$ cm, the value of b (in cm) is\underline{\hspace{1.5cm}}.


\begin{tcolorbox}[solbox, title=Detailed Solution -- Question 99]
Cut-off frequency of a rectangular waveguide is 

\[
\begin{aligned} f_{c}&=\frac{c}{2 \pi} \sqrt{\left(\frac{m \pi}{a}\right)^{2}+\left(\frac{n\pi}{b}\right)^{2}} \\ f_{c} \text{ for }\quad T E_{10}&=\frac{c}{2 \pi} \sqrt{\left(\frac{1 \pi}{a}\right)^{2}}=\frac{c}{2 a} \\ f_{c} \text{ for }\quad T E_{20}&=\frac{c}{2 \pi} \sqrt{\left(\frac{2 \pi}{a}\right)^{2}}=\frac{c}{2 a} \\ f_{c} \text { for } \quad T E_{11} &=\frac{c}{2 \pi} \sqrt{\left(\frac{\pi}{a}\right)^{2}}+\left(\frac{\pi}{b}\right)^{2} \\ &=\frac{c}{2} \sqrt{\left(\frac{1}{a}\right)^{2}+\left(\frac{1}{b}\right)^{2}} \\ \text { Given, } f_{C T_{11}} &=\frac{f_{C_{I E 10}}+f_{C I E_{20}}}{2} \end{aligned}
\]

\[
\begin{aligned} \frac{c}{2} \sqrt{\left(\frac{1}{a}\right)^{2}+\left(\frac{1}{b}\right)^{2}}&=\frac{\frac{c}{2 a}+\frac{c}{a}}{2} \\ \sqrt{\left(\frac{1}{a}\right)^{2}+\left(\frac{1}{b}\right)^{2}}&=\frac{1}{2 a}+\frac{1}{a}=\frac{3}{2 a} \\ \left(\frac{1}{a}\right)^{2}+\left(\frac{1}{b}\right)^{2}&=\frac{9}{4 a^{2}} \\ \Rightarrow \quad\left(\frac{1}{b}\right)^{2} &=\frac{5}{4 a^{2}} \\ b &=\frac{2 a}{\sqrt{5}}=\frac{2 \times \sqrt{5}}{\sqrt{5}} \\ &=2 \mathrm{cm} \end{aligned}
\]
\end{tcolorbox}

\vspace{1.5em}
\hrule
\vspace{1.5em}

\section*{Question 100: Transmission Lines (GATE EC 2014-SET-2)}
\addcontentsline{toc}{section}{Question 100: Transmission Lines (GATE EC 2014-SET-2)}
\noindent \textbf{\color{primary}NAT} \hfill \textbf{\color{secondary}2 Marks}


\noindent In the transmission line shown, the impedance $Z_{in}$ (in ohms) between node A and
the ground is \underline{\hspace{1.5cm}}.

\begin{center}\includegraphics[max width=0.85\linewidth]{images/img\_14.jpg}\end{center}

\begin{center}\includegraphics[max width=0.85\linewidth]{images/img\_14.jpg}\end{center}


\begin{tcolorbox}[solbox, title=Detailed Solution -- Question 100]
\begin{center}\includegraphics[max width=0.85\linewidth]{images/img\_43.jpg}\end{center}

\begin{center}\includegraphics[max width=0.85\linewidth]{images/img\_43.jpg}\end{center}

\[
\begin{aligned} Z_{\mathrm{in}_{1}} &=Z_{0}\left(\frac{Z_{L}+j Z_{0} \tan \beta l}{Z_{0}+j Z_{L} \tan \beta l}\right) \\ \beta l &=\frac{2 \pi}{\lambda} \times \frac{\lambda}{2}=\pi \\ Z_{\mathrm{in}} &=Z_{0}\left(\frac{Z_{L}+j 0}{Z_{0}+j 0}\right) \\ Z_{\mathrm{in}_{1}} &=Z_{L}=50 \Omega \\ \therefore \quad Z_{\mathrm{in}} &=\left(Z_{\mathrm{in}_{1}} \| 100 \Omega\right)=\frac{100 \times 50}{150}=33.3 \Omega \end{aligned}
\]
\end{tcolorbox}

\vspace{1.5em}
\hrule
\vspace{1.5em}

\section*{Question 101: Uniform Plane Waves (GATE EC 2014-SET-2)}
\addcontentsline{toc}{section}{Question 101: Uniform Plane Waves (GATE EC 2014-SET-2)}
\noindent \textbf{\color{primary}MCQ} \hfill \textbf{\color{secondary}2 Marks}


\noindent If the electric field of a plane wave is 
 $\vec{E}(z,t)=\hat{x}3 cos(\omega t-kz+30^{\circ})$- $\hat{y}4sin(\omega t-kz+45^{\circ})$ (mV/m), 
 the polarization state of the plane wave is


\begin{enumerate}[label=(\Alph*)]
    \item \textbf{(Correct)} left elliptical
    \item left circular
    \item right elliptical
    \item right circular
\end{enumerate}

\noindent \textbf{Correct Answer:} (A)

\begin{tcolorbox}[solbox, title=Detailed Solution -- Question 101]
\[
\vec{E}(z, t)=3 \hat{x} \cos \left(\omega t-k z+30^{\circ}\right)-4 \hat{y} \sin \left(\omega t-k z+45^{\circ}\right)
\]

 The plane wave is travelling in +z direction

\[
\begin{array}{l} \hat{E}_{x}=3 \cos \left(\omega t-k z+30^{\circ}\right) \\ \hat{E}_{y}=-4 \sin \left(\omega t-k z+45^{\circ}\right) \end{array}
\]

 At z=0, tracing E locus for various times

\[
\begin{array}{c} \text { At } t=0, E_{x}=3 \cos 30^{\circ}=3 \cdot \frac{1}{2}=\frac{3}{2}=1.5 \\ E_{y}=-4 \sin 45^{\circ}=-2 \sqrt{2}=-2.83 \end{array}
\]

 At $t=\frac{T}{6}$
$E_{x}=3 \cos \left(\frac{2 \pi}{T} \times \frac{T}{6}+30^{\circ}\right)=0$
$E_{y}=-4 \sin \left(60^{\circ}+45^{\circ}\right)=-3.86$

\begin{center}\includegraphics[max width=0.85\linewidth]{images/img\_67.jpg}\end{center}

\begin{center}\includegraphics[max width=0.85\linewidth]{images/img\_67.jpg}\end{center}

This is left elliptical polarized plane-wave.
\end{tcolorbox}

\vspace{1.5em}
\hrule
\vspace{1.5em}

\section*{Question 102: Uniform Plane Waves (GATE EC 2014-SET-2)}
\addcontentsline{toc}{section}{Question 102: Uniform Plane Waves (GATE EC 2014-SET-2)}
\noindent \textbf{\color{primary}MCQ} \hfill \textbf{\color{secondary}1 Mark}


\noindent Which one of the following field patterns represents a TEM wave travelling in the
positive x direction ?


\begin{enumerate}[label=(\Alph*)]
    \item $E=+8\hat{y}, H=-4\hat{z}$
    \item \textbf{(Correct)} $E=-2\hat{y}, H=-3\hat{z}$
    \item $E=+2\hat{z}, H=+2\hat{y}$
    \item $E=-3\hat{y}, H=+4\hat{z}$
\end{enumerate}

\noindent \textbf{Correct Answer:} (B)

\begin{tcolorbox}[solbox, title=Detailed Solution -- Question 102]
Since, the TEM wave or plane wave is travelling in  positive x-direction, then conditions should be

\[
\begin{aligned} P &=\vec{E} \times \vec{H} \\ \hat{a}_{x} &=\hat{a}_{y} \times \hat{a}_{z}\\ \text{or, }\quad\hat{a}_{x}&=\left(\hat{a}_{y}\right) \times\left(-\hat{a}_{z}\right) \end{aligned}
\]

 So, the option (B) satisfies the condition.
\end{tcolorbox}

\vspace{1.5em}
\hrule
\vspace{1.5em}

\section*{Question 103: Transmission Lines (GATE EC 2014-SET-2)}
\addcontentsline{toc}{section}{Question 103: Transmission Lines (GATE EC 2014-SET-2)}
\noindent \textbf{\color{primary}NAT} \hfill \textbf{\color{secondary}1 Mark}


\noindent To maximize power transfer, a lossless transmission line is to be matched to a resistive load impedance via a  $\lambda/4$   transformer as shown. 

\begin{center}\includegraphics[max width=0.85\linewidth]{images/img\_74.jpg}\end{center}

\begin{center}\includegraphics[max width=0.85\linewidth]{images/img\_74.jpg}\end{center}

  The characteristic impedance (in $\Omega$) of the $\lambda/4$  transformer is \underline{\hspace{1.5cm}}


\begin{tcolorbox}[solbox, title=Detailed Solution -- Question 103]
Input impedance for quarter wave transformer is given as,

\[
\begin{aligned} Z_{\text {in }}&=\frac{Z_{0}^{2}}{Z_{L}} \\ \text { Given: } Z_{\text {in }}&=50 \Omega, \quad Z_{L}=100 \Omega\\ \therefore \quad Z_{0}^{2}&=50 \times 100 \\ Z_{0}&=\sqrt{50 \times 100} \\ Z_{0}&=70.72 \Omega \end{aligned}
\]
\end{tcolorbox}

\vspace{1.5em}
\hrule
\vspace{1.5em}

\section*{Question 104: Transmission Lines (GATE EC 2014-SET-1)}
\addcontentsline{toc}{section}{Question 104: Transmission Lines (GATE EC 2014-SET-1)}
\noindent \textbf{\color{primary}NAT} \hfill \textbf{\color{secondary}2 Marks}


\noindent The input impedance of a $\frac{\lambda }{8}$ section of a lossless transmission line of characteristic impedance 50 $\Omega$ is found to be real when the other end is terminated by a load $Z_{L}=(R+jX)\Omega$. If X is  30 $\Omega$, the value of R(in $\Omega$) is\underline{\hspace{1.5cm}}.


\begin{tcolorbox}[solbox, title=Detailed Solution -- Question 104]
\[
\begin{aligned} Z_{\text {in }}&=Z_{0}\left(\frac{Z_{L}+j Z_{o} \tan \beta l}{Z_{o}+j Z_{L} \tan \beta l}\right)\\ \text{at }\;l&=\frac{\lambda}{8} \\ \Rightarrow \quad \beta l&=\frac{2 \pi}{\lambda} \times \frac{\lambda}{8}=\frac{\pi}{4} \\ \therefore \quad\tan \frac{\pi}{4}&=1 \\ Z_{\text {in }}&=Z_{o}\left[\frac{Z_{L}+j 50}{Z_{o}+j Z_{L}}\right] \end{aligned}
\]

\[
\begin{aligned} \text{Given},\quad Z_{L}&=R+j X=R+j 30 \text{ and }Z_{0}=50 \Omega \\ 30 R &=80(50+R) \\ Z_{\text {in }} &=50\left[\frac{R+j 30+j 50}{50+j(R+j 30)}\right] \\ &=50\left[\frac{R+j 80}{(50-30)+j R}\right]=50\left[\frac{R+j 80}{20+j R}\right] \\ &=50\left[\frac{(R+j 80)(20-j R)}{20^{2}+R^{2}}\right] \end{aligned}
\]

 Since, only real part of $Z_{\text {in }}$ exists, so imaginary 
 part of $Z_{\text {in }}=0$

\[
\begin{aligned} \therefore \quad Z_{\text {in }} &=50\left[\frac{20 R+1600 j-j R^{2}+50 R}{20^{2}+R^{2}}\right] \\ \text { Zimaginary } &=0=\frac{1600-R^{2}}{20^{2}+R^{2}}\\ \therefore \quad R^{2} &=1600 \\ R &=40 \Omega \end{aligned}
\]
\end{tcolorbox}

\vspace{1.5em}
\hrule
\vspace{1.5em}

\section*{Question 105: Transmission Lines (GATE EC 2014-SET-1)}
\addcontentsline{toc}{section}{Question 105: Transmission Lines (GATE EC 2014-SET-1)}
\noindent \textbf{\color{primary}MCQ} \hfill \textbf{\color{secondary}2 Marks}


\noindent For a parallel plate transmission line, let v be the speed of propagation and Z
be the characteristic impedance. Neglecting fringe effects, a reduction of the
spacing between the plates by a factor of two results in


\begin{enumerate}[label=(\Alph*)]
    \item halving of v and no change in Z
    \item \textbf{(Correct)} no changes in v and halving of Z
    \item no change in both v and Z
    \item halving of both v and Z
\end{enumerate}

\noindent \textbf{Correct Answer:} (B)

\begin{tcolorbox}[solbox, title=Detailed Solution -- Question 105]
\[
\begin{aligned} Z_{0}&=\sqrt{\frac{L}{C}} \text{ for parallel plate line}\\ \frac{C}{l}&=\frac{\epsilon W}{d} \text{ and } \frac{L}{l}=\frac{\mu d}{W} \\ \text{where}, \quad W&= \text{Width of the line}\\ d &=\text { Separation between the lines } \\ Z_{0} & \propto d \quad (Z_{0}\text{ change by half})\\ \text { Velocity } &=\frac{1}{\sqrt{L C}}=\frac{1}{\sqrt{\mu_{0} \epsilon_{0}}}=3 \times 10^{8} \mathrm{m} / \mathrm{s}\\ &\qquad\text{(No change in V)} \end{aligned}\\
\]
\end{tcolorbox}

\vspace{1.5em}
\hrule
\vspace{1.5em}

\section*{Question 106: Antennas \& Radiating Systems (GATE EC 2014-SET-1)}
\addcontentsline{toc}{section}{Question 106: Antennas \& Radiating Systems (GATE EC 2014-SET-1)}
\noindent \textbf{\color{primary}NAT} \hfill \textbf{\color{secondary}2 Marks}


\noindent In spherical coordinates, let $\hat{a}_{\theta },\hat{a}_{\phi }$
denote unit vectors along the $\theta, \phi$  directions. 
$E=\frac{100}{r}sin\theta cos(\omega t-\beta r)\hat{a}_\theta V/m$ 
 and 
  $E=\frac{0.265}{r}sin\theta cos(\omega t-\beta r)\hat{a}_\phi A/m$  
  represent the electric and magnetic field components of the EM wave at large
distances $r$  from a dipole antenna, in free space. The average power (W) crossing
the hemispherical shell located at r = 1 km, $0\leq \theta \leq \pi /2$ is \underline{\hspace{1.5cm}}.


\begin{tcolorbox}[solbox, title=Detailed Solution -- Question 106]
\[
\begin{aligned} \vec{E} &=\operatorname{Re}\left[\frac{100}{r} \sin \theta e^{-\beta r} e^{j \omega t}\right] \hat{a}_{\theta} V / m \\ &=\operatorname{Re}\left[\frac{0.265}{r} \sin \theta e^{-\beta \beta r} e^{j \omega t}\right] \hat{a}_{\phi} V / m \end{aligned}
\]

 Poynting vector,

\[
\vec{P}=\frac{1}{2} \operatorname{Re}\left[\vec{E} \times \vec{H}^{*}\right]=\frac{26.5}{2 r^{2}} \sin ^{2} \theta \hat{a}_{r} W / m^{2}
\]

 The average power crossing the given 
 hemispherical shell is,

\[
\begin{aligned} \text { Power } &=\int_{0}^{\pi / 2} \int_{\phi=0}^{2 \pi} \vec{P} \cdot r^{2} \sin \theta d \theta d \phi \hat{a}_{r} \\ &=\frac{26.5}{2} \times 2 \pi \int_{\theta=0}^{\pi / 2} \sin ^{3} \theta d \theta \\ &=\frac{26.5}{2} \times 2 \pi \frac{2}{3}=55.5 \mathrm{W} \end{aligned}
\]
\end{tcolorbox}

\vspace{1.5em}
\hrule
\vspace{1.5em}

\section*{Question 107: Basics of Electromagnetics \& Maxwell's Eqns (GATE EC 2014-SET-1)}
\addcontentsline{toc}{section}{Question 107: Basics of Electromagnetics \& Maxwell's Eqns (GATE EC 2014-SET-1)}
\noindent \textbf{\color{primary}MCQ} \hfill \textbf{\color{secondary}1 Mark}


\noindent The force on a point charge $+q$ kept at a distance d from the surface of an
infinite grounded metal plate in a medium of permittivity $\varepsilon$
is


\begin{enumerate}[label=(\Alph*)]
    \item 0
    \item $\frac{q^{2}}{16\pi \varepsilon d^{2}}$ away from the plate
    \item \textbf{(Correct)} $\frac{q^{2}}{16\pi \varepsilon d^{2}}$ towards the plate
    \item $\frac{q^{2}}{4\pi \varepsilon d^{2}}$ towards the plate
\end{enumerate}

\noindent \textbf{Correct Answer:} (C)


\vspace{1.5em}
\hrule
\vspace{1.5em}

\section*{Question 108: Miscellaneous Topics (GATE EC 2014-SET-1)}
\addcontentsline{toc}{section}{Question 108: Miscellaneous Topics (GATE EC 2014-SET-1)}
\noindent \textbf{\color{primary}MCQ} \hfill \textbf{\color{secondary}1 Mark}


\noindent A two-port network has scattering parameters given by 
\[
[S]=\begin{bmatrix} S_{11} &S_{12} \\ S_{21} & S_{22} \end{bmatrix}
\]
.  If the
port-2 of the two port is short circuited, the $S_{11}$
parameter for the resultant one port network is


\begin{enumerate}[label=(\Alph*)]
    \item $\frac{S_{11}-S_{11}S_{22}+S_{12}S_{21}}{1+S_{22}}$
    \item \textbf{(Correct)} $\frac{S_{11}+S_{11}S_{22}-S_{12}S_{21}}{1+S_{22}}$
    \item $\frac{S_{11}+S_{11}S_{22}+S_{12}S_{21}}{1-S_{22}}$
    \item $\frac{S_{11}-S_{11}S_{22}+S_{12}S_{21}}{1-S_{22}}$
\end{enumerate}

\noindent \textbf{Correct Answer:} (B)

\begin{tcolorbox}[solbox, title=Detailed Solution -- Question 108]
\[
\begin{array}{l} b_{1}=s_{11} \cdot a_{1}+s_{12} \cdot a_{2} \qquad\ldots(i)\\ b_{2}=s_{21} \cdot a_{1}+s_{22} a_{2}\qquad\ldots(ii) \end{array}
\]

 When output is short, $a_{2}=-b_{2}$

\[
\begin{aligned} -a_{2} &=s_{21} \cdot a_{1}+s_{22} a_{2} \\ \Rightarrow \quad a_{2} &=\frac{-s_{21} \cdot a_{1}}{1+s_{22}} \end{aligned}
\]

 Substituting in first equation,

\[
\begin{aligned} b_{1} &=s_{11} \cdot a_{1}+s_{12}\left(\frac{-s_{21} \cdot a_{1}}{1+s_{22}}\right) \\ &=\frac{s_{11} \cdot a_{1}+s_{11} \cdot s_{22} a_{1}-s_{12} \cdot s_{21} a_{1}}{1+s_{22}} \\ \frac{b_{1}}{a_{1}} &=s_{11}=\frac{s_{11}+s_{11} \cdot s_{22}-s_{12} \cdot s_{21}}{1+s_{22}} \end{aligned}
\]
\end{tcolorbox}

\vspace{1.5em}
\hrule
\vspace{1.5em}

\section*{Question 109: Uniform Plane Waves (GATE EC 2013)}
\addcontentsline{toc}{section}{Question 109: Uniform Plane Waves (GATE EC 2013)}
\noindent \textbf{\color{primary}MCQ} \hfill \textbf{\color{secondary}1 Mark}


\noindent A monochromatic plane wave of wavelength $\lambda =600\mu m$ is propagating in the direction as shown in the figure below.$\overrightarrow{E_{i}}$ , $\overrightarrow{E_{r}}$ , and $\overrightarrow{E_{t}}$ denote incident, reflected, and transmitted electric field vectors associated with the wave 

\begin{center}\includegraphics[max width=0.85\linewidth]{images/img\_81.jpg}\end{center}

\begin{center}\includegraphics[max width=0.85\linewidth]{images/img\_81.jpg}\end{center}

 The expression for $\overrightarrow{E_{r}}$ is


\begin{enumerate}[label=(\Alph*)]
    \item \textbf{(Correct)} \[
0.23\frac{E_{0}}{\sqrt{2}}\left ( \hat{a_{x}}+\hat{a_{z}} \right )e^{-j\frac{\pi \times 10^{4}\left ( x-z \right )}{3\sqrt{2}}}V/m
\]
    \item \[
-\frac{E_{0}}{\sqrt{2}}\left ( \hat{a_{x}}+\hat{a_{z}} \right )e^{j\frac{\pi \times 10^{4}z}{3}}V/m
\]
    \item \[
0.44\frac{E_{0}}{\sqrt{2}}\left ( \hat{a_{x}}+\hat{a_{z}} \right )e^{-j\frac{\pi \times 10^{4}\left ( x-z \right )}{3\sqrt{2}}}V/m
\]
    \item \[
\frac{E_{0}}{\sqrt{2}}\left ( \hat{a_{x}}+\hat{a_{z}} \right )e^{-j\frac{\pi \times 10^{4}\left ( x+z \right )}{3}}V/m
\]
\end{enumerate}

\noindent \textbf{Correct Answer:} (A)

\begin{tcolorbox}[solbox, title=Detailed Solution -- Question 109]
For reflected wave. 
$E_{r}$ propagation direction changes without change in angle $45^{\circ}$
$E_{r} \text { propagation direction }=\frac{\hat{a}_{x}-\hat{a}_{z}}{\sqrt{2}}$
$\Gamma$reflection coefficient
$\Gamma=-0.23$
$E_{r}$ field direction suitably changes satisfying
$E \cdot P=0$
\end{tcolorbox}

\vspace{1.5em}
\hrule
\vspace{1.5em}

\section*{Question 110: Uniform Plane Waves (GATE EC 2013)}
\addcontentsline{toc}{section}{Question 110: Uniform Plane Waves (GATE EC 2013)}
\noindent \textbf{\color{primary}MCQ} \hfill \textbf{\color{secondary}1 Mark}


\noindent A monochromatic plane wave of wavelength $\lambda =600\mu m$ is propagating in the direction as shown in the figure below.$\overrightarrow{E_{i}}$ , $\overrightarrow{E_{r}}$ , and $\overrightarrow{E_{t}}$ denote incident, reflected, and transmitted electric field vectors associated with the wave 

\begin{center}\includegraphics[max width=0.85\linewidth]{images/img\_81.jpg}\end{center}

\begin{center}\includegraphics[max width=0.85\linewidth]{images/img\_81.jpg}\end{center}
 
The angle of incidence $\theta _{i}$ and the expression for $\overrightarrow{E_{i}}$ are


\begin{enumerate}[label=(\Alph*)]
    \item $60^{\circ}$ and 
\[
\frac{E_{0}}{\sqrt{2}}\left ( \hat{a_{x}}-\hat{a_{z}} \right )e^{-j\frac{\pi \times 10^{4}\left ( x+z \right )}{3\sqrt{2}}}V/m
\]
    \item $45^{\circ}$ and 
\[
\frac{E_{0}}{\sqrt{2}}\left ( \hat{a_{x}}+\hat{a_{z}} \right )e^{-j\frac{\pi \times 10^{4}z}{3}}V/m
\]
    \item \textbf{(Correct)} $45^{\circ}$ and 
\[
\frac{E_{0}}{\sqrt{2}}\left ( \hat{a_{x}}-\hat{a_{z}} \right )e^{-j\frac{\pi \times 10^{4}\left ( x+z \right )}{3\sqrt{2}}}V/m
\]
    \item $60^{\circ}$ and 
\[
\frac{E_{0}}{\sqrt{2}}\left ( \hat{a_{x}}-\hat{a_{z}} \right )e^{-j\frac{\pi \times 10^{4}z}{3}}V/m
\]
\end{enumerate}

\noindent \textbf{Correct Answer:} (C)

\begin{tcolorbox}[solbox, title=Detailed Solution -- Question 110]
$E_{i}$ propagates in ZXplane. 
$E_{i}$ directed in ZX-plane. 
 It is a p-plane of incidence polarized wave.
 Using Snell's law,

\[
\begin{aligned} \frac{\sin \theta_{i}}{\sin 19.2^{\circ}} &=\sqrt{\frac{4.5 \epsilon_{o}}{\epsilon_{o}}} &\theta_{i}=45^{\circ} \\ \beta_{x} &=\beta_{z}=\beta \cos 45^{\circ}=\frac{\beta}{\sqrt{2}} \end{aligned}
\]

 (for propagation direction)
 For field direction $=\frac{\hat{a}_{x}-\hat{a}_{z}}{\sqrt{2}}\left(45^{\circ}\right)$
\end{tcolorbox}

\vspace{1.5em}
\hrule
\vspace{1.5em}

\section*{Question 111: Transmission Lines (GATE EC 2013)}
\addcontentsline{toc}{section}{Question 111: Transmission Lines (GATE EC 2013)}
\noindent \textbf{\color{primary}MCQ} \hfill \textbf{\color{secondary}1 Mark}


\noindent The return loss of a device is found to be 20 dB. The voltage standing wave ratio (VSWR) and magnitude of reflection coefficient are respectively


\begin{enumerate}[label=(\Alph*)]
    \item \textbf{(Correct)} 1.22 and 0.1
    \item 0.81 and 0.1
    \item - 1.22 and 0.1
    \item 2.44 and 0.2
\end{enumerate}

\noindent \textbf{Correct Answer:} (A)

\begin{tcolorbox}[solbox, title=Detailed Solution -- Question 111]
\[
\begin{aligned} \text { Return loss }&=-20 \log |\Gamma| \\ -20 \log |\Gamma|&=20 \\ \Rightarrow \quad|\Gamma|&=10^{-1}=0.1 \\ \text{VSWR}&=\frac{1+|\Gamma|}{1-|\Gamma|}=\frac{1+0.1}{1-0.1}=1.22 \end{aligned}
\]
\end{tcolorbox}

\vspace{1.5em}
\hrule
\vspace{1.5em}

\section*{Question 112: Basics of Electromagnetics \& Maxwell's Eqns (GATE EC 2012)}
\addcontentsline{toc}{section}{Question 112: Basics of Electromagnetics \& Maxwell's Eqns (GATE EC 2012)}
\noindent \textbf{\color{primary}MCQ} \hfill \textbf{\color{secondary}1 Mark}


\noindent An infinitely long uniform solid wire of radius a carries a uniform dc current of density $\overrightarrow{j}$. 

A hole of radius $b (b  <  a)$ is now drilled along the length of the wire at a distance d from the center
of the wire as shown below. 

\begin{center}\includegraphics[max width=0.85\linewidth]{images/img\_77.jpg}\end{center}

\begin{center}\includegraphics[max width=0.85\linewidth]{images/img\_77.jpg}\end{center}
 
 The magnetic field inside the hole is


\begin{enumerate}[label=(\Alph*)]
    \item \textbf{(Correct)} uniform and depends only on d
    \item uniform and depends only on b
    \item uniform and depends on both b and d
    \item non uniform
\end{enumerate}

\noindent \textbf{Correct Answer:} (A)

\begin{tcolorbox}[solbox, title=Detailed Solution -- Question 112]
\begin{center}\includegraphics[max width=0.85\linewidth]{images/img\_91.jpg}\end{center}

\begin{center}\includegraphics[max width=0.85\linewidth]{images/img\_91.jpg}\end{center}

The H field due to entire cylinder without hole.
$H\left(r_{1}\right)=\frac{J r_{1}}{2}$
 The Hfield due to hole only.
$H\left(r_{2}\right)=\frac{J r_{2}}{2}$
 Current inside the hole is zero means
$H(r)=\frac{J r_{1}}{2}-\frac{J r_{2}}{2} \simeq \frac{J d}{2}$
\end{tcolorbox}

\vspace{1.5em}
\hrule
\vspace{1.5em}

\section*{Question 113: Basics of Electromagnetics \& Maxwell's Eqns (GATE EC 2012)}
\addcontentsline{toc}{section}{Question 113: Basics of Electromagnetics \& Maxwell's Eqns (GATE EC 2012)}
\noindent \textbf{\color{primary}MCQ} \hfill \textbf{\color{secondary}1 Mark}


\noindent An infinitely long uniform solid wire of radius a carries a uniform dc current of density $\overrightarrow{j}$. 

The magnetic field at a distance $r$  from the center of the wire is proportional to


\begin{enumerate}[label=(\Alph*)]
    \item $r \; for \; r < a \; and \; 1/r^{2}\; for \;r > a$
    \item $0 \; for \; r  <  a \; and \; 1 /r \; for \; r  >  a$
    \item \textbf{(Correct)} $r \; for \; r  <  a \; and \; 1 /r \; for \; r  >  a$
    \item $0 \; for\; r < a \; and \; 1/r^{2} \;for \; r > a$
\end{enumerate}

\noindent \textbf{Correct Answer:} (C)

\begin{tcolorbox}[solbox, title=Detailed Solution -- Question 113]
We know that magnetic flux density at a distance r from the wire is, 

\[
\begin{aligned} |B| &=\frac{\mu_{0} I}{2 \pi r} \\ \text { For } r < a \rightarrow I&=J \cdot \pi r^{2} \\ \text { so } \quad |B|&=\frac{\mu_{0} J r}{2}\\ |B| \propto r \text { for } r < a \\ \text { for } r > a,\quad I&=J \times \pi a^{2} \\ \text { So, } |B|&=\frac{\mu_{0} J \pi a^{2}}{2 \pi r} \\ |B| \propto \frac{1}{r} \text { for } r > a \end{aligned}
\]
\end{tcolorbox}

\vspace{1.5em}
\hrule
\vspace{1.5em}

\section*{Question 114: Waveguides \& Optical Fibers (GATE EC 2012)}
\addcontentsline{toc}{section}{Question 114: Waveguides \& Optical Fibers (GATE EC 2012)}
\noindent \textbf{\color{primary}MCQ} \hfill \textbf{\color{secondary}1 Mark}


\noindent The magnetic field along the propagation direction inside a rectangular waveguide with the crosssection shown in the figure is $H_{z}=3 cos(2.094 \times  10^{2}x)cos(2.618 \times 10^{2}y)$ $cos(6.283 \times 10^{10}t-\beta z)$

\begin{center}\includegraphics[max width=0.85\linewidth]{images/img\_58.jpg}\end{center}

\begin{center}\includegraphics[max width=0.85\linewidth]{images/img\_58.jpg}\end{center}
 
 The phase velocity $v_{p}$ of the wave inside the waveguide satisfies


\begin{enumerate}[label=(\Alph*)]
    \item \textbf{(Correct)} $v_{p}> c$
    \item $v_{p}=c$
    \item $0 < v_{p}< c$
    \item $v_{p}=0$
\end{enumerate}

\noindent \textbf{Correct Answer:} (A)

\begin{tcolorbox}[solbox, title=Detailed Solution -- Question 114]
From the given data

\[
\begin{aligned} \omega_{c} &=\left(\sqrt{\left(\frac{m \pi}{a}\right)^{2}+\left(\frac{n\pi}{b}\right)^{2}}\right) c \\ f_{c} &=(\sqrt{(2.09)^{2}+(2.618)^{2}}) \times \frac{3 \times 10^{10}}{2 \times 3.14} \\ &=16 \mathrm{GHz} \\ f &=\frac{6.283 \times 10^{10}}{2 \times 3.14}=10 \mathrm{GHz} \\ \text { Hence, } \quad f& < f_{c^{\prime}} \gamma=\alpha+j 0\\ \beta=0, \quad \quad v_{p}&=\frac{\omega}{\beta}=\infty \end{aligned}
\]
\end{tcolorbox}

\vspace{1.5em}
\hrule
\vspace{1.5em}

\section*{Question 115: Transmission Lines (GATE EC 2012)}
\addcontentsline{toc}{section}{Question 115: Transmission Lines (GATE EC 2012)}
\noindent \textbf{\color{primary}MCQ} \hfill \textbf{\color{secondary}1 Mark}


\noindent A transmission line with a characteristic impedance of 100 $\Omega$ is used to match a 50 $\Omega$ section to a 200 $\Omega$ section. If the matching is to be done both at 429 MHz and 1 GHz, the length of the transmission line can be approximately


\begin{enumerate}[label=(\Alph*)]
    \item 82.5 cm
    \item 1.05 cm
    \item \textbf{(Correct)} 1.58 cm
    \item 1.75 cm
\end{enumerate}

\noindent \textbf{Correct Answer:} (C)

\begin{tcolorbox}[solbox, title=Detailed Solution -- Question 115]
The $100 \Omega$ device used to match $50 \Omega$ and $200 \Omega$ devices is a quarter wave transformer. 
Length of matching device

$=(2 n+1) \frac{\lambda}{4}=\text { odd multiple }$

Where $n$ -depends on frequency of $\lambda / 4$

$\lambda_{1} \text { of frequency } 429 \mathrm{MHz}=0.7 \mathrm{m}$

$\lambda_{2}$ of frequency $1 \mathrm{GHz}=0.3 \mathrm{m}$

Common multiple $=2.1 \mathrm{m}$

Dividing by $4=0.525 \mathrm{m}$ (not available in options) 
Another odd multiple of $0.525 \times 3=1.58 \mathrm{m}$
\end{tcolorbox}

\vspace{1.5em}
\hrule
\vspace{1.5em}

\section*{Question 116: Antennas \& Radiating Systems (GATE EC 2012)}
\addcontentsline{toc}{section}{Question 116: Antennas \& Radiating Systems (GATE EC 2012)}
\noindent \textbf{\color{primary}MCQ} \hfill \textbf{\color{secondary}1 Mark}


\noindent The radiation pattern of an antenna in spherical co-ordinates is given by 
 $F(\theta )=cos^{4}\theta ;0\leq \theta\leq \pi /2$ 
The directivity of the antenna is


\begin{enumerate}[label=(\Alph*)]
    \item \textbf{(Correct)} 10 dB
    \item 12.6 dB
    \item 11.5 dB
    \item 18 dB
\end{enumerate}

\noindent \textbf{Correct Answer:} (A)

\begin{tcolorbox}[solbox, title=Detailed Solution -- Question 116]
$F(\theta)=\cos ^{4} \theta, 0 \leq \theta \leq \pi / 2$
 we know that directivity D is
$D=\frac{4 \pi u_{\max }}{\pi_{\mathrm{rad}}} \qquad\ldots(1)$
$\mathrm{F}(\theta)$ is nothing but radiation intensity $u(\theta, \phi)$ and
$W_{\mathrm{rad}}$ is radiated power.

\[
W_{\mathrm{rad}}=\int_{0}^{2 \pi \pi / 2} \int_{0}^{2} \cos ^{4} \theta \sin \theta d \phi\qquad\ldots(2)
\]

 above equation is written from the formula
$W_{\mathrm{rad}}=\iint_{\theta \phi} u(\theta, \phi) d \Omega$
 where $\alpha \Omega$ is solid angle and $d \Omega=\sin \theta d \theta d \phi$
 So from (2)

\[
\begin{aligned} W_{\mathrm{rad}} &=\int_{0}^{2 \pi \pi / 2} \int_{0}^{2} \cos ^{4} \theta \sin d \theta d \phi \\ W_{\mathrm{rad}} &=2 \pi \int_{0}^{\pi / 2} \cos ^{4} \theta \sin \theta d \theta \\W_{\mathrm{rad}}&=-2 \pi\left[\frac{\cos ^{5} \theta}{5}\right]_{0}^{\pi / 2}\\
W_{\mathrm{rad}}&=\frac{2 \pi}{5}&\ldots(3)\\
\text { So, } \quad D&=\frac{4 \pi u_{\max }}{\left(\frac{2 \pi}{5}\right)}=10 u_{\max }=10\\
&\text { In } \mathrm{dB} \text { directivity }=10 \mathrm{log}_{10} \mathrm{D}=10 \mathrm{dB}\end{aligned}
\]
\end{tcolorbox}

\vspace{1.5em}
\hrule
\vspace{1.5em}

\section*{Question 117: Transmission Lines (GATE EC 2012)}
\addcontentsline{toc}{section}{Question 117: Transmission Lines (GATE EC 2012)}
\noindent \textbf{\color{primary}MCQ} \hfill \textbf{\color{secondary}1 Mark}


\noindent A coaxial cable with an inner diameter of 1 mm and outer diameter of 2.4 mm is filled with a
dielectric of relative permittivity 10.89. Given $\mu _{0}=4 \pi \times 10^{-7} H/m,$ $\varepsilon _{0}=\frac{10^{-9}}{36 \pi} F/m$, the characteristic impedance of the cable is


\begin{enumerate}[label=(\Alph*)]
    \item 330 $\Omega$
    \item 100 $\Omega$
    \item 143.3 $\Omega$
    \item \textbf{(Correct)} 15.3 $\Omega$
\end{enumerate}

\noindent \textbf{Correct Answer:} (D)

\begin{tcolorbox}[solbox, title=Detailed Solution -- Question 117]
Characteristics impedance of the co-axial cable is given by 
$Z_{o}=\frac{138}{\sqrt{\epsilon_{r}}} \log \frac{D}{d}$
$=\frac{138}{\sqrt{10.89}} \log \left(\frac{2.4}{1}\right)=15.89 \Omega$
\end{tcolorbox}

\vspace{1.5em}
\hrule
\vspace{1.5em}

\section*{Question 118: Uniform Plane Waves (GATE EC 2012)}
\addcontentsline{toc}{section}{Question 118: Uniform Plane Waves (GATE EC 2012)}
\noindent \textbf{\color{primary}MCQ} \hfill \textbf{\color{secondary}1 Mark}


\noindent The electric field of a uniform plane electromagnetic wave in free space, along the positive
$x$ direction, is given by $\vec{E}=10(\hat{a}_{y}+j\hat{a}_{z})e^{-j25x}$. The frequency and polarization of the wave, respectively, are


\begin{enumerate}[label=(\Alph*)]
    \item \textbf{(Correct)} 1.2 GHz and left circular
    \item 4 Hz and left circular
    \item 1.2 GHz and right circular
    \item 4 Hz and right circular
\end{enumerate}

\noindent \textbf{Correct Answer:} (A)

\begin{tcolorbox}[solbox, title=Detailed Solution -- Question 118]
Out of phase by $90^{\circ}$and equal amplitude the wave is circularly polarized.
 Check the trace of E with the time which gives left  circularly polarized. 

\[
\begin{aligned} a_{y}&=\sin \omega t \\ a_{z}&=\cos \omega t \end{aligned}
\]
\end{tcolorbox}

\vspace{1.5em}
\hrule
\vspace{1.5em}

\section*{Question 119: Uniform Plane Waves (GATE EC 2012)}
\addcontentsline{toc}{section}{Question 119: Uniform Plane Waves (GATE EC 2012)}
\noindent \textbf{\color{primary}MCQ} \hfill \textbf{\color{secondary}1 Mark}


\noindent A plane wave propagating in air with $\vec{E}=(8\hat{a}_{x}+6\hat{a}_{y}+5\hat{a}_{z})e^{j(\omega t+3x-4y)}$ V/m is incident on a perfectly conducting slab positioned at 0 $\leq$ x. The E field of the reflected wave is


\begin{enumerate}[label=(\Alph*)]
    \item $(-8\hat{a}_{x}-6\hat{a}_{y}-5\hat{a}_{z})e^{j(\omega t+3x+4y)} V/m$
    \item $(-8\hat{a}_{x}+6\hat{a}_{y}-5\hat{a}_{z})e^{j(\omega t+3x+4y)} V/m$
    \item \textbf{(Correct)} $(-8\hat{a}_{x}-6\hat{a}_{y}-5\hat{a}_{z})e^{j(\omega t-3x-4y)} V/m$
    \item $(-8\hat{a}_{x}+6\hat{a}_{y}-5\hat{a}_{z})e^{j(\omega t-3x-4y)} V/m$
\end{enumerate}

\noindent \textbf{Correct Answer:} (C)

\begin{tcolorbox}[solbox, title=Detailed Solution -- Question 119]
\begin{center}\includegraphics[max width=0.85\linewidth]{images/img\_78.jpg}\end{center}

\begin{center}\includegraphics[max width=0.85\linewidth]{images/img\_78.jpg}\end{center}

Incident wave propagation $\rightarrow-3 \hat{a}_{x}+4 \hat{a}_{y}$
 Reflected wave propagation $\rightarrow 3 \hat{a}_{x}+4 \hat{a}_{y}$
$E_{i}$ incident direction $\rightarrow 8 \hat{a}_{x}+6 \hat{a}_{y}+5 \hat{a}_{z}$

\[
\left.\begin{array}{l}E_{i} \text { tangential } \rightarrow 6 \hat{a}_{y}+5 \hat{a}_{z} \\ E_{r} \text { tangential } \rightarrow-6 \hat{a}_{y}-5 \hat{a}_{z}\end{array}\right.} E_{\text {tang }}=0
\]

\[
\left.\begin{array}{rl}E_{i} \text { normal } & =8 \hat{a}_{x} \\ E_{r} \text { normal } & =8 \hat{a}_{x}\end{array}\right.} E_{\text {normal }}=E_{\max
\]

$E_{r}$ reflected direction $=8 \hat{a}_{x}-6 \hat{a}_{y}-5 \hat{a}_{z}$
\end{tcolorbox}

\vspace{1.5em}
\hrule
\vspace{1.5em}

\section*{Question 120: Uniform Plane Waves (GATE EC 2011)}
\addcontentsline{toc}{section}{Question 120: Uniform Plane Waves (GATE EC 2011)}
\noindent \textbf{\color{primary}MCQ} \hfill \textbf{\color{secondary}1 Mark}


\noindent The electric and magnetic fields for a TEM wave of frequency 14 GHz in a
homogeneous medium of relative permittivity er and relative permeability $\mu _{r}$=1 are given by
 $\vec{E}=E_{p}e^{j(\omega t-280\pi t)}\hat{u}_{z} V/m$ 
$\vec{H}=3e^{j(\omega t-280\pi y)}\hat{u}_{x} V/m$ 

Assuming the speed of light in free space to be $3 \times 10^{8}$ m/s, the intrinsic impedance of free space to be 120$\pi$, the relative permittivity $\varepsilon_{r}$ of the medium and the electric field amplitude $E_{p}$ are


\begin{enumerate}[label=(\Alph*)]
    \item $\varepsilon _{r}=3,E_{p}=120 \pi$
    \item $\varepsilon _{r}=3,E_{p}=360 \pi$
    \item $\varepsilon _{r}=9,E_{p}=360 \pi$
    \item \textbf{(Correct)} $\varepsilon _{r}=9,E_{p}=120 \pi$
\end{enumerate}

\noindent \textbf{Correct Answer:} (D)

\begin{tcolorbox}[solbox, title=Detailed Solution -- Question 120]
\[
\begin{aligned} \text{Given:}\quad\vec{E}&=E_{p} e^{i(\omega t-280 \pi y)} \hat{u}_{z} V / m\\ \vec{H}&=3 e^{i(\omega t-280 \pi y)} \hat{u}_{x} A / m \\ v&=3 \times 10^{8} \mathrm{m} / \mathrm{s} \end{aligned}
\]

 From given expression we conclude that,

\[
\begin{aligned} \beta &=280 \pi=\frac{2 \pi}{\lambda} \\ \therefore \quad & \lambda=\frac{1}{140} \text { meter } \\ v &=f \lambda=14 \times 10^{9} \times \frac{1}{140} \mathrm{m} / \mathrm{se} \\ v &=1 \times 10^{8} \mathrm{m} / \mathrm{sec} \\ \text { but, } \quad v &=\frac{c}{\sqrt{\epsilon_{r} \mu_{r}}} \\ \therefore \quad 1 \times 10^{8} &=\frac{3 \times 10^{8}}{\sqrt{1 \times \epsilon_{r}}} \\ \therefore \quad \epsilon_{r}&=9\\ \frac{E_{P}}{H_{P}}&=\eta=\sqrt{\frac{\mu}{\epsilon}}=\sqrt{\frac{\mu_{0} \times 1}{\epsilon_{0} \times 9}}=\frac{1}{3} \sqrt{\frac{\mu_{0}}{\epsilon_{0}}} \\ \Rightarrow \frac{E_{P}}{3}&=\frac{1}{3} \times 120 \pi \\ \therefore \quad E_{P}&=120 \pi \end{aligned}
\]
\end{tcolorbox}

\vspace{1.5em}
\hrule
\vspace{1.5em}

\section*{Question 121: Transmission Lines (GATE EC 2011)}
\addcontentsline{toc}{section}{Question 121: Transmission Lines (GATE EC 2011)}
\noindent \textbf{\color{primary}MCQ} \hfill \textbf{\color{secondary}1 Mark}


\noindent A transmission line of characteristic impedance 50$\Omega$ is terminated in a load
impedance $Z_{L}$. The VSWR of the line is measured as 5 and the first of the
voltage maxima in the line is observed at a distance of $\frac{\lambda }{4}$ from the load. The value of $Z_{L}$ is


\begin{enumerate}[label=(\Alph*)]
    \item \textbf{(Correct)} 10 $\Omega$
    \item 250 $\Omega$
    \item (19.23 + j46.15) $\Omega$
    \item (19.23 - j 46.15) $\Omega$
\end{enumerate}

\noindent \textbf{Correct Answer:} (A)

\begin{tcolorbox}[solbox, title=Detailed Solution -- Question 121]
$\text { VSWR }=5=\frac{1+|\Gamma|}{1-|\Gamma|}$
$\therefore |\Gamma|=\frac{2}{3} \Rightarrow \Gamma=\pm \frac{2}{3}$
 If maximum is at $\lambda / 4$ from load then minimum will be at load itself.
 For the phase of $\Gamma$
$2 \beta x_{\max }=2 n \pi+\theta$
 where $x_{\max }$ is distance of $V_{\max }$ from load.

\[
\begin{aligned} 2 \cdot \frac{2 \pi}{\lambda} \cdot \frac{\lambda}{4}&=\theta \\ \Gamma&=-\frac{2}{3}\\ \text{Now},-\frac{2}{3}&=\frac{Z_{L}-Z_{0}}{Z_{L}+Z_{0}}=\frac{Z_{L}-50}{Z_{L}+50}\\ \Rightarrow Z_{L}&=10 \Omega \end{aligned}
\]
\end{tcolorbox}

\vspace{1.5em}
\hrule
\vspace{1.5em}

\section*{Question 122: Basics of Electromagnetics \& Maxwell's Eqns (GATE EC 2011)}
\addcontentsline{toc}{section}{Question 122: Basics of Electromagnetics \& Maxwell's Eqns (GATE EC 2011)}
\noindent \textbf{\color{primary}MCQ} \hfill \textbf{\color{secondary}1 Mark}


\noindent A current sheet $\bar{J}$ = 10$\hat{u}_{y}$ A/m lies on the dielectric interface $x=0$ between two dielectric media with $\varepsilon _{r1}$=5, $\mu _{r1}$=1 in Region-1 $(x < 0)$ and $\varepsilon _{r2}$=5, $\mu _{r2}$=2 in Region-2 $(x > 0)$. If the magnetic field in Region-1 at $x = 0^{-}$ is $\vec{H}_{1} = 3\hat{u}_{x}$ + 30$\hat{u}_{y}$ A/m. The magnetic field in Region-2 at $x = 0^{+}$ is  

\begin{center}\includegraphics[max width=0.85\linewidth]{images/img\_71.jpg}\end{center}

\begin{center}\includegraphics[max width=0.85\linewidth]{images/img\_71.jpg}\end{center}


\begin{enumerate}[label=(\Alph*)]
    \item \textbf{(Correct)} $\vec{H_{2}}=1.5\hat{u}_{x}+30\hat{u}_{y}-10\hat{u}_{z} A/m$
    \item $\vec{H}_{2}=3\hat{u}_{x}+30\hat{u}_{y}-10\hat{u}_{z} A/m$
    \item $\vec{H}_{2}=1.5\hat{u}_{x}+40\hat{u}_{y} A/m$
    \item $\vec{H}_{2}=3\hat{u}_{x}+30\hat{u}_{y}+10\hat{u}_{z} A/m$
\end{enumerate}

\noindent \textbf{Correct Answer:} (A)

\begin{tcolorbox}[solbox, title=Detailed Solution -- Question 122]
For magnetic field boundary relations are 

\[
\begin{aligned} \vec{B}_{n_{1}} &=\vec{B}_{n_{2}} \\ \text { and } \vec{H}_{t_{1}}-\vec{H}_{t_{2}} &=-\vec{J}_{s} \times \vec{a}_{n} \\ \Rightarrow \quad B_{x_{1}}&=B_{x_{2}}\qquad\text{(x is normal component)}\\ \mu_{1} H_{x_{1}} &=\mu_{2} H_{x 2} \\ \Rightarrow \quad 1 \times 3 &=2 \times H_{x 2} \\ H_{x_{2}} &=1.5 \\ \Rightarrow \quad \vec{H}_{11}-\vec{H}_{t_{2}} &=-10 \hat{u}_{y} \times \hat{u}_{x}=+10 \hat{u}_{z} \\ \vec{H}_{t_{2}} &=H_{t_{1}}-10 \hat{u}_{z} \\ \vec{H}_{t_{2}} &=30 \hat{u}_{y}-10 \hat{u}_{z} \\ \Rightarrow \quad \vec{H}_{2} &=1.5 \hat{u}_{x}+30 \hat{u}_{y}-10 \hat{u}_{z} \mathrm{A} / \mathrm{m} \end{aligned}
\]
\end{tcolorbox}

\vspace{1.5em}
\hrule
\vspace{1.5em}

\section*{Question 123: Waveguides \& Optical Fibers (GATE EC 2011)}
\addcontentsline{toc}{section}{Question 123: Waveguides \& Optical Fibers (GATE EC 2011)}
\noindent \textbf{\color{primary}MCQ} \hfill \textbf{\color{secondary}1 Mark}


\noindent The modes in a rectangular waveguide are denoted by $\frac{TE_{mn}}{TM_{mn}}$  where
m and n are the eigen numbers along the larger and smaller dimensions of the
waveguide respectively. Which one of the following statements is TRUE?


\begin{enumerate}[label=(\Alph*)]
    \item \textbf{(Correct)} The $TM_{10}$ mode of the wave does not exist
    \item The $TE_{10}$ mode of the wave does not exist
    \item The $TM_{10}$ and $TE_{10}$ the modes both exist and have the same cut-off frequencies
    \item The $TM_{10}$ and $TM_{01}$ modes both exist and have the same cut-off frequencies
\end{enumerate}

\noindent \textbf{Correct Answer:} (A)

\begin{tcolorbox}[solbox, title=Detailed Solution -- Question 123]
In case of rectangular wave guide  $T E_{\text {rm }}$  exists for  all values of  m  and  n except  m=0  and  n=0 . For   $T M_{m n}$  Io exist both values of mand  n  must be non
zero
\end{tcolorbox}

\vspace{1.5em}
\hrule
\vspace{1.5em}

\section*{Question 124: Transmission Lines (GATE EC 2011)}
\addcontentsline{toc}{section}{Question 124: Transmission Lines (GATE EC 2011)}
\noindent \textbf{\color{primary}MCQ} \hfill \textbf{\color{secondary}1 Mark}


\noindent A transmission line of characteristic impedance 50 W is terminated by a 50 $\Omega$ load. When excited by a sinusoidal voltage source at 10 GHz, the phase
difference between two points spaced 2 mm apart on the line is found to $\frac{\pi }{4}$ radians. The phase velocity of the wave along the line is


\begin{enumerate}[label=(\Alph*)]
    \item $0.8 \times 10^{8} m/s$
    \item $1.2 \times 10^{8} m/s$
    \item \textbf{(Correct)} $1.6 \times 10^{8} m/s$
    \item $3 \times 10^{8} m/s$
\end{enumerate}

\noindent \textbf{Correct Answer:} (C)

\begin{tcolorbox}[solbox, title=Detailed Solution -- Question 124]
\[
\begin{aligned} \text { Phase difference }&=\frac{2 \pi}{\lambda} \text { (path difference) } \\ \Rightarrow \frac{\pi}{4}&=\frac{2 \pi}{\lambda}\left(2 \times 10^{-3}\right) \\ \therefore \quad &=8 \times 2 \times 10^{-3} \\ &=16 \times 10^{-3} \mathrm{m} \\ f&= 10 \mathrm{GHz}=10 \times 10^{9} \mathrm{Hz} \end{aligned}
\]

 Hence the phase velocity of wave along the line is,

\[
\begin{aligned} & v=f \lambda=10 \times 10^{9} \times 16 \times 10^{-3} \mathrm{m} / \mathrm{sec} \\ \therefore \quad v &=1.6 \times 10^{8} \mathrm{m} / \mathrm{s} \end{aligned}
\]
\end{tcolorbox}

\vspace{1.5em}
\hrule
\vspace{1.5em}

\section*{Question 125: Antennas \& Radiating Systems (GATE EC 2011)}
\addcontentsline{toc}{section}{Question 125: Antennas \& Radiating Systems (GATE EC 2011)}
\noindent \textbf{\color{primary}MCQ} \hfill \textbf{\color{secondary}1 Mark}


\noindent Consider the following statements regarding the complex Poynting vector $\vec{P}$ for the power radiated by a point source in an infinite homogeneous and lossless medium. Re($\vec{P}$) denotes the real part of $\vec{P}$, S denotes a spherical surface whose centre is at the point source, and $\hat{n}$
denotes the unit surface normal on S. Which of the following statements is TRUE?


\begin{enumerate}[label=(\Alph*)]
    \item Re($\vec{P}$) remains constant at any radial distance from the source
    \item Re($\vec{P}$) increases with increasing radial distance from the source
    \item \textbf{(Correct)} $\oint \oint_{S}Re(\vec{P}).\hat{n} dS$ remains constant at any radial distance from the source
    \item $\oint \oint_{S}Re(\vec{P}).\hat{n} dS$ decreases with increasing radial distance form the source
\end{enumerate}

\noindent \textbf{Correct Answer:} (C)

\begin{tcolorbox}[solbox, title=Detailed Solution -- Question 125]
Power density of any point source decreases with
distance i.e. the density decreases and area of
cross-over increases with the product being
constant.
\end{tcolorbox}

\vspace{1.5em}
\hrule
\vspace{1.5em}

\section*{Question 126: Transmission Lines (GATE EC 2010)}
\addcontentsline{toc}{section}{Question 126: Transmission Lines (GATE EC 2010)}
\noindent \textbf{\color{primary}MCQ} \hfill \textbf{\color{secondary}1 Mark}


\noindent In the circuit shown, all the transmission line sections are lossless. The Voltage
Standing Wave Ration(VSWR) on the 60 $\Omega$ line is  

\begin{center}\includegraphics[max width=0.85\linewidth]{images/img\_8.jpg}\end{center}

\begin{center}\includegraphics[max width=0.85\linewidth]{images/img\_8.jpg}\end{center}


\begin{enumerate}[label=(\Alph*)]
    \item 1.00
    \item \textbf{(Correct)} 1.64
    \item 2.50
    \item 3.00
\end{enumerate}

\noindent \textbf{Correct Answer:} (B)

\begin{tcolorbox}[solbox, title=Detailed Solution -- Question 126]
\[
Z_{\text {in }}=Z_{0}\left(\frac{Z_{L}+j Z_{0} \tan \beta l}{Z_{0}+j Z_{L} \tan \beta l}\right)
\]

\begin{center}\includegraphics[max width=0.85\linewidth]{images/img\_23.jpg}\end{center}

\begin{center}\includegraphics[max width=0.85\linewidth]{images/img\_23.jpg}\end{center}

 Input impedance,

\[
\begin{aligned} Z_{1} &=30\left[\frac{0+j 30 \tan \left(\frac{2 \pi}{\lambda}\right)\left(\frac{\lambda}{8}\right)}{30+0}\right] \\ \Rightarrow Z_{1} &=j 30 \\ Z_{2} &=30 \sqrt{2}\left[\frac{30+j 30 \sqrt{2} \tan \left(\frac{2 \pi}{\lambda}\right)\left(\frac{\lambda}{4}\right)}{30 \sqrt{2}+j 30 \tan \left(\frac{2 \pi}{\lambda}\right)\left(\frac{\lambda}{4}\right)}\right] \\ &=30 \sqrt{2}\left[\frac{\frac{30}{\tan \pi / 2}+j 30 \sqrt{2}}{\frac{30 \sqrt{2}}{\tan \pi / 2}+j 30}\right] \\ &=60 \Omega \\ \text { or, } Z_{2}&=\frac{(30 \sqrt{2})^{2}}{30}=60 \Omega & \end{aligned}
\]

 Load impedance,
$Z_{L}=Z_{1}+Z_{2}=j 30+6$
 Magnitude of reflection coefficient,

\[
\begin{aligned} |\Gamma|&=\left|\frac{Z_{L}-Z_{0}}{Z_{L}+Z_{0}}\right|=\left|\frac{60+j 30-60}{60+j 30+60}\right|\\ &=\left|\frac{j 30}{120+j 30}\right|=\mid \frac{j 1}{4+j \mid} \\ &=\frac{1}{\sqrt{16+1}=\frac{1}{\sqrt{17}}} \end{aligned}
\]

 VSWR on 60 $\Omega$ line,

\[
\text { VSWR }=\frac{1+|\Gamma|}{1-|\Gamma|}=\frac{1+\frac{1}{\sqrt{17}}}{1-\frac{1}{\sqrt{17}}}=1.64
\]
\end{tcolorbox}

\vspace{1.5em}
\hrule
\vspace{1.5em}

\section*{Question 127: Uniform Plane Waves (GATE EC 2010)}
\addcontentsline{toc}{section}{Question 127: Uniform Plane Waves (GATE EC 2010)}
\noindent \textbf{\color{primary}MCQ} \hfill \textbf{\color{secondary}1 Mark}


\noindent A plane wave having the electric field components $\vec{E}_{i}=24 \cos (3 \times 10^{8}t-\beta y)\hat{a}_{z} V/m$ and traveling in free space is incident normally on a lossless medium with $\mu =\mu_{0}$ and $\varepsilon =9\varepsilon_{0}$ which occupies the region $y\geq 0$. The reflected magnetic field component is given by


\begin{enumerate}[label=(\Alph*)]
    \item \textbf{(Correct)} $\frac{1}{10\pi }cos(3 \times 10^{8}t+y)\hat{a}_{x} \; A/m$
    \item $\frac{1}{20\pi }cos(3 \times 10^{8}t+y)\hat{a}_{x}\; A/m$
    \item $-\frac{1}{20\pi }cos(3 \times 10^{8}t+y)\hat{a}_{x} \; A/m$
    \item $-\frac{1}{10\pi }cos(3 \times 10^{8}t+y)\hat{a}_{x} \; A/m$
\end{enumerate}

\noindent \textbf{Correct Answer:} (A)

\begin{tcolorbox}[solbox, title=Detailed Solution -- Question 127]
$\vec{E}_{i}=24 \cos \left(3 \times 10^{8} t-\beta y\right) \hat{a}_{z} \vee / m$
 The wave is travelling in +y direction.

\[
\begin{aligned} H_{i}=& \frac{1}{\eta}\left(\hat{a}_{y} \times \vec{E}_{i}\right) \\ =& \frac{24 \cos }{120 \pi}\left(3 \times 10^{8} t-\beta y\right) \hat{a}_{x} \\ =& \frac{1}{5 \pi} \cos \left(3 \times 10^{8} t-\beta y\right) \hat{a}_{x} \\ \frac{H_{r}}{H_{i}}&=\frac{\eta_{1}-\eta_{2}}{\eta_{1}+\eta_{2}}=\frac{\sqrt{\frac{\mu_{0}}{\epsilon_{1}}}-\sqrt{\frac{\mu_{0}}{\epsilon_{2}}}}{\sqrt{\frac{\mu_{0}}{\epsilon_{1}}}+\sqrt{\frac{\mu_{0}}{\epsilon_{2}}}} \\ =\frac{\sqrt{\epsilon_{2}}-\sqrt{\epsilon_{1}}}{\sqrt{\epsilon_{2}}+\sqrt{\epsilon_{1}}}&=\frac{\sqrt{9}-\sqrt{1}}{\sqrt{9}+\sqrt{1}}=\frac{3-1}{3+1}=\frac{2}{4}=\frac{1}{2} \end{aligned}
\]

Therefore, reflected magnetic field component,

\[
\vec{H}_{r}=\frac{1}{10 \pi} \cos \left(3 \times 10^{8} t+\beta y\right) \hat{a}_{x}
\]

Due to change in propagation direction to -y
\end{tcolorbox}

\vspace{1.5em}
\hrule
\vspace{1.5em}

\section*{Question 128: Uniform Plane Waves (GATE EC 2010)}
\addcontentsline{toc}{section}{Question 128: Uniform Plane Waves (GATE EC 2010)}
\noindent \textbf{\color{primary}MCQ} \hfill \textbf{\color{secondary}1 Mark}


\noindent The electric field component of a time harmonic plane EM wave traveling in
a nonmagnetic lossless dielectric medium has an amplitude of 1 V/m. If the
relative permittivity of the medium is 4, the magnitude of the time-average
power density vector (in W/$m^{2}$) is


\begin{enumerate}[label=(\Alph*)]
    \item $\frac{1}{30\pi }$
    \item $\frac{1}{60\pi }$
    \item \textbf{(Correct)} $\frac{1}{120\pi }$
    \item $\frac{1}{240\pi }$
\end{enumerate}

\noindent \textbf{Correct Answer:} (C)

\begin{tcolorbox}[solbox, title=Detailed Solution -- Question 128]
Magnitude of the time-average power density
 vector

\[
\begin{aligned} P_{a v} &=\frac{E^{2}}{2 \eta} \\ \text { where, } \quad\eta &=\sqrt{\mu / \epsilon} \\ \eta &=\sqrt{\frac{\mu_{0}}{4 \epsilon_{0}}}=\frac{\eta_{0}}{2}=\frac{120 \pi}{2} \\ P_{a v} &=\frac{1^{2}}{2\left(\frac{120 \pi}{2}\right)}=\frac{1}{120 \pi} \mathrm{W} / \mathrm{m}^{2} \end{aligned}
\]
\end{tcolorbox}

\vspace{1.5em}
\hrule
\vspace{1.5em}

\section*{Question 129: Transmission Lines (GATE EC 2010)}
\addcontentsline{toc}{section}{Question 129: Transmission Lines (GATE EC 2010)}
\noindent \textbf{\color{primary}MCQ} \hfill \textbf{\color{secondary}1 Mark}


\noindent A transmission line has a characteristic impedance of 50 $\Omega$ and a resistance of 0.1$\Omega$/m . If the line is distortion less, the attenuation constant(in Np/m) is


\begin{enumerate}[label=(\Alph*)]
    \item 500
    \item 5
    \item 0.014
    \item \textbf{(Correct)} 0.002
\end{enumerate}

\noindent \textbf{Correct Answer:} (D)

\begin{tcolorbox}[solbox, title=Detailed Solution -- Question 129]
For distortionless transmission line,

\[
\begin{aligned} L G&=R C \\ \Rightarrow \frac{L}{C}&=\frac{R}{G} \end{aligned}
\]

 Characteristic impedance,
$Z_{0}=\sqrt{\frac{L}{C}}=\sqrt{\frac{R}{G}}$
 Attenuation constant,

\[
\begin{aligned} \alpha &=\sqrt{R G}=\sqrt{R} \cdot \frac{\sqrt{R}}{Z_{0}}=\frac{R}{Z_{0}}=\frac{0.1}{50} \\ &=0.002 \mathrm{Np} / \mathrm{m} \end{aligned}
\]
\end{tcolorbox}

\vspace{1.5em}
\hrule
\vspace{1.5em}

\section*{Question 130: Miscellaneous Topics (GATE EC 2010)}
\addcontentsline{toc}{section}{Question 130: Miscellaneous Topics (GATE EC 2010)}
\noindent \textbf{\color{primary}MCQ} \hfill \textbf{\color{secondary}1 Mark}


\noindent If the scattering matrix [S] of a two port network is 

\[
|S|=\begin{bmatrix} 0.2\angle 0^{\circ}&0.9\angle 90^{\circ} \\ 0.9\angle 90^{\circ}& 0.1\angle 90^{\circ} \end{bmatrix}
\]

 then the network is


\begin{enumerate}[label=(\Alph*)]
    \item lossless and reciprocal
    \item lossless but not reciprocal
    \item \textbf{(Correct)} not lossless but reciprocal
    \item neither lossless nor reciprocal
\end{enumerate}

\noindent \textbf{Correct Answer:} (C)

\begin{tcolorbox}[solbox, title=Detailed Solution -- Question 130]
For reciprocal networks, $S_{12}=S_{21}$
 For symmetrical networks, $S_{11}=S_{22}$
 For antimetrical networks, $S_{11}=-S_{22}$
 For lossless reciprocal networks, 
$\left|S_{11}\right|=\left|S_{22}\right|$

and $\left|S_{11}\right|^{2}+\left|S_{12}\right|^{2}=1$
\end{tcolorbox}

\vspace{1.5em}
\hrule
\vspace{1.5em}

\section*{Question 131: Transmission Lines (GATE EC 2009)}
\addcontentsline{toc}{section}{Question 131: Transmission Lines (GATE EC 2009)}
\noindent \textbf{\color{primary}MCQ} \hfill \textbf{\color{secondary}1 Mark}


\noindent A transmission line terminates in two branches, each of length $\frac{\lambda }{4}$, as shown. The branches are terminated by 50$\Omega$ loads. The lines are lossless and have the characteristic impedances shown. Determine the impedance $Z_{i}$ as seen by the source 

\begin{center}\includegraphics[max width=0.85\linewidth]{images/img\_6.jpg}\end{center}

\begin{center}\includegraphics[max width=0.85\linewidth]{images/img\_6.jpg}\end{center}


\begin{enumerate}[label=(\Alph*)]
    \item $200\Omega$
    \item $100\Omega$
    \item $50\Omega$
    \item \textbf{(Correct)} $25\Omega$
\end{enumerate}

\noindent \textbf{Correct Answer:} (D)

\begin{tcolorbox}[solbox, title=Detailed Solution -- Question 131]
\[
\begin{aligned} z_{1}&=\frac{Z_{0}^{2}}{Z_{4}}=\frac{(100)^{2}}{50}=200 \Omega \\ Z_{2}&=\frac{Z_{0}^{2}}{Z_{2}}=\frac{(100)^{2}}{50}=200 \Omega \\ Z_{L}&=Z_{1} \| I_{2}=100 \Omega \\ Z_{i}&=\frac{Z_{0}^{2}}{Z_{L}}=\frac{(50)^{2}}{100}=25 \Omega \end{aligned}
\]
\end{tcolorbox}

\vspace{1.5em}
\hrule
\vspace{1.5em}

\section*{Question 132: Basics of Electromagnetics \& Maxwell's Eqns (GATE EC 2009)}
\addcontentsline{toc}{section}{Question 132: Basics of Electromagnetics \& Maxwell's Eqns (GATE EC 2009)}
\noindent \textbf{\color{primary}MCQ} \hfill \textbf{\color{secondary}1 Mark}


\noindent A magnetic field in air is measured to be 
 $\vec{B}=B_{0}(\frac{x}{x^{2}+y^{2}}\hat{y}-\frac{y}{x^{2}+y^{2}}\hat{x})$ 
What current distribution leads to this field ?
[Hint : The algebra is trivial in cylindrical coordinates.]


\begin{enumerate}[label=(\Alph*)]
    \item $\vec{J}=\frac{B_{0}\hat{z}}{\mu _{0}}(\frac{1}{x^{2}+y^{2}}),r\neq 0$
    \item $\vec{J}=\frac{B_{0}\hat{z}}{\mu _{0}}(\frac{2}{x^{2}+y^{2}}),r\neq 0$
    \item \textbf{(Correct)} $\vec{J}=0,r\neq 0$
    \item $\vec{J}=\frac{B_{0}\hat{z}}{\mu _{0}}(\frac{1}{x^{2}+y^{2}}),r\neq 0$
\end{enumerate}

\noindent \textbf{Correct Answer:} (C)

\begin{tcolorbox}[solbox, title=Detailed Solution -- Question 132]
Applying $\nabla \cdot H=J$

\[
\begin{aligned} \nabla \times B&=\left|\begin{array}{ccc} \hat{a}_{x} & \hat{a}_{y} & \hat{a}_{z} \\ \frac{\partial}{\partial x} & \frac{\partial}{\partial y} & \frac{\partial}{\partial z} \\ \frac{-y}{x^{2}+y^{2}} & \frac{x}{x^{2}+y^{2}} & 0 \end{array}\right|=0 \\ J&=0 \end{aligned}
\]
\end{tcolorbox}

\vspace{1.5em}
\hrule
\vspace{1.5em}

\section*{Question 133: Basics of Electromagnetics \& Maxwell's Eqns (GATE EC 2009)}
\addcontentsline{toc}{section}{Question 133: Basics of Electromagnetics \& Maxwell's Eqns (GATE EC 2009)}
\noindent \textbf{\color{primary}MCQ} \hfill \textbf{\color{secondary}1 Mark}


\noindent If a vector field $\vec{V}$ is related to another vector field $\vec{A}$ through $\vec{V}=\bigtriangledown *\vec{A}$, which of the following is true? (Note : C and $S_C$ refer to any closed contour and any surface whose boundary is C.)


\begin{enumerate}[label=(\Alph*)]
    \item \[
\oint_{C}^{}\vec{V} \cdot \vec{dl}=\int_{S_{C}}^{} \int_{}^{}\vec{A}\cdot \vec{dS}
\]
    \item \textbf{(Correct)} \[
\oint_{C}^{}\vec{A} \cdot \vec{dl}=\int_{S_{C}}^{} \int_{}^{}\vec{V}\cdot \vec{dS}
\]
    \item \[
\oint_{C}^{}\triangledown \times \vec{V} \cdot \vec{dl}=\int_{S_{C}}^{} \int_{}^{}\triangledown \times \vec{A}\cdot \vec{dS}
\]
    \item \[
\oint_{C}^{}\triangledown \times \vec{A} \cdot \vec{dl}=\int_{S_{C}}^{} \int_{}^{}\vec{V}\cdot \vec{dS}
\]
\end{enumerate}

\noindent \textbf{Correct Answer:} (B)


\vspace{1.5em}
\hrule
\vspace{1.5em}

\section*{Question 134: Basics of Electromagnetics \& Maxwell's Eqns (GATE EC 2009)}
\addcontentsline{toc}{section}{Question 134: Basics of Electromagnetics \& Maxwell's Eqns (GATE EC 2009)}
\noindent \textbf{\color{primary}MCQ} \hfill \textbf{\color{secondary}1 Mark}


\noindent Two infinitely long wires carrying current are as shown in the figure below. One wire is in the y-z plane and parallel to the y-axis. The other wire is in the x-y plane and parallel to the x-axis. Which components of the resulting magnetic field are non-zero at the origin ? 

\begin{center}\includegraphics[max width=0.85\linewidth]{images/img\_73.jpg}\end{center}

\begin{center}\includegraphics[max width=0.85\linewidth]{images/img\_73.jpg}\end{center}


\begin{enumerate}[label=(\Alph*)]
    \item x,y, z components
    \item x,y components
    \item y,z components
    \item \textbf{(Correct)} x,z components
\end{enumerate}

\noindent \textbf{Correct Answer:} (D)

\begin{tcolorbox}[solbox, title=Detailed Solution -- Question 134]
$H_{dicection}$ due to any current wire
 =$I$ flow direction $\times$ (cross) radial vector from current to point of consideration.
\end{tcolorbox}

\vspace{1.5em}
\hrule
\vspace{1.5em}

\section*{Question 135: Waveguides \& Optical Fibers (GATE EC 2009)}
\addcontentsline{toc}{section}{Question 135: Waveguides \& Optical Fibers (GATE EC 2009)}
\noindent \textbf{\color{primary}MCQ} \hfill \textbf{\color{secondary}1 Mark}


\noindent Which of the following statements is true regarding the fundamental mode of
the metallic waveguides shown ?

\begin{center}\includegraphics[max width=0.85\linewidth]{images/img\_65.jpg}\end{center}

\begin{center}\includegraphics[max width=0.85\linewidth]{images/img\_65.jpg}\end{center}


\begin{enumerate}[label=(\Alph*)]
    \item \textbf{(Correct)} Only P has no cutoff-frequency
    \item Only Q has no cutoff-frequency
    \item Only R has no cutoff-frequency
    \item All three have cutoff-frequencies
\end{enumerate}

\noindent \textbf{Correct Answer:} (A)

\begin{tcolorbox}[solbox, title=Detailed Solution -- Question 135]
P is coaxial line and support TEM wave
$\therefore$  P has no cutoff frequency

Q and R are wave-guides and cutoff frequency
each depends upon their dimensions.
\end{tcolorbox}

\vspace{1.5em}
\hrule
\vspace{1.5em}

\section*{Question 136: Antennas \& Radiating Systems (GATE EC 2008)}
\addcontentsline{toc}{section}{Question 136: Antennas \& Radiating Systems (GATE EC 2008)}
\noindent \textbf{\color{primary}MCQ} \hfill \textbf{\color{secondary}1 Mark}


\noindent At 20 GHz, the gain of a parabolic dish antenna of 1 meter and 70% efficiency is


\begin{enumerate}[label=(\Alph*)]
    \item 15 dB
    \item 25 dB
    \item 35 dB
    \item \textbf{(Correct)} 45 dB
\end{enumerate}

\noindent \textbf{Correct Answer:} (D)

\begin{tcolorbox}[solbox, title=Detailed Solution -- Question 136]
\[
\text { Gain }=\frac{4 \pi}{\lambda^{2}} A_{e}=\frac{4 \pi}{\lambda^{2}} \pi R^{2}
\]

 Along with illumination or efficiency

\[
\begin{aligned} \text { Gain } &=\eta \pi^{2} \frac{(2 R)^{2}}{\lambda^{2}}=7\left(\frac{D}{\lambda}\right)^{2} \\ &=0.7 \pi^{2}\left(\frac{1 \times 20 \times 10^{9}}{3 \times 10^{8}}\right)^{2} \\ &=30705.4359 \\ G(\mathrm{dB}) &=10 \log _{10} G \\ &=10 \log _{10}(30705.4359)=45 \mathrm{dB} \end{aligned}
\]
\end{tcolorbox}

\vspace{1.5em}
\hrule
\vspace{1.5em}

\section*{Question 137: miscellaneous Topics (GATE EC 2008)}
\addcontentsline{toc}{section}{Question 137: miscellaneous Topics (GATE EC 2008)}
\noindent \textbf{\color{primary}MCQ} \hfill \textbf{\color{secondary}1 Mark}


\noindent In the design of a single mode step index optical fibre close to upper cut-off, the
single-mode operation is not preserved if


\begin{enumerate}[label=(\Alph*)]
    \item radius as well as operating wavelength are halved
    \item radius as well as operating wavelength are doubled
    \item \textbf{(Correct)} radius is halved and operating wavelength is doubled
    \item radius is doubled and operating wavelength is halved
\end{enumerate}

\noindent \textbf{Correct Answer:} (C)


\vspace{1.5em}
\hrule
\vspace{1.5em}

\section*{Question 138: Uniform Plane Waves (GATE EC 2008)}
\addcontentsline{toc}{section}{Question 138: Uniform Plane Waves (GATE EC 2008)}
\noindent \textbf{\color{primary}MCQ} \hfill \textbf{\color{secondary}1 Mark}


\noindent A uniform plane wave in the free space is normally incident on an infinitely
thick dielectric slab (dielectric constant $\varepsilon_{r}$=9). The magnitude of the reflection coefficient is


\begin{enumerate}[label=(\Alph*)]
    \item 0
    \item 0.3
    \item \textbf{(Correct)} 0.5
    \item 0.8
\end{enumerate}

\noindent \textbf{Correct Answer:} (C)

\begin{tcolorbox}[solbox, title=Detailed Solution -- Question 138]
Dielectric constant of free space, $\epsilon_{r_{1}}=1$
 Dielectric constant of dielectric slab, $\epsilon_{r_{2}}=9$
 Magnitude of reflection coefficient 

\[
|\Gamma|=\left|\frac{\sqrt{\epsilon_{r_{1}}}-\sqrt{\epsilon_{r_{2}}}}{\sqrt{\epsilon_{r_{1}}}+\sqrt{\epsilon_{r_{2}}}}\right|=\left|\frac{1-3}{1+3}\right|=\frac{2}{4}=0.5
\]
\end{tcolorbox}

\vspace{1.5em}
\hrule
\vspace{1.5em}

\section*{Question 139: Transmission Lines (GATE EC 2008)}
\addcontentsline{toc}{section}{Question 139: Transmission Lines (GATE EC 2008)}
\noindent \textbf{\color{primary}MCQ} \hfill \textbf{\color{secondary}1 Mark}


\noindent One end of a loss-less transmission line having the characteristic impedance of
75$\Omega$ and length of 1 cm is short-circuited. At 3 GHz, the input impedance at the other end of transmission line is


\begin{enumerate}[label=(\Alph*)]
    \item 0
    \item Resistive
    \item Capacitive
    \item \textbf{(Correct)} Inductive
\end{enumerate}

\noindent \textbf{Correct Answer:} (D)

\begin{tcolorbox}[solbox, title=Detailed Solution -- Question 139]
It is a short circuit line with $l =1 cm.$
 As $\quad \lambda=\frac{3 \times 10^{8}}{3 \times 10^{9}}=0.3=30 \mathrm{cm}$
$l < \frac{\lambda}{4}$
 The line has inductive nature.
 As $Z_{S C}=j Z_{o} \tan \beta l$
\end{tcolorbox}

\vspace{1.5em}
\hrule
\vspace{1.5em}

\section*{Question 140: Waveguides \& Optical Fibers (GATE EC 2008)}
\addcontentsline{toc}{section}{Question 140: Waveguides \& Optical Fibers (GATE EC 2008)}
\noindent \textbf{\color{primary}MCQ} \hfill \textbf{\color{secondary}1 Mark}


\noindent A rectangular waveguide of internal dimensions (a=4cm and b=3cm) is to
be operated in $TE_{11}$ mode. The minimum operating frequency is


\begin{enumerate}[label=(\Alph*)]
    \item \textbf{(Correct)} 6.25 GHz
    \item 6.0 GHz
    \item 5.0 GHz
    \item 3.75 GHz
\end{enumerate}

\noindent \textbf{Correct Answer:} (A)

\begin{tcolorbox}[solbox, title=Detailed Solution -- Question 140]
Minimum operating frequency,

\[
\begin{aligned} f_{0}&=\frac{c}{2 \pi} \sqrt{\left(\frac{m \pi}{a}\right)^{2}+\left(\frac{n \pi}{b}\right)^{2}}\\ \text{For }\;T E_{11}&\;\text{ mode}\\ f_{0} &=\frac{3 \times 10^{8}}{2} \sqrt{\left(\frac{1}{0.04}\right)^{2}+\left(\frac{1}{0.03}\right)^{2}} \\ &=\frac{3 \times 10^{8}}{2} \times \frac{0.05}{0.04 \times 0.03}=6.25 \mathrm{GHz} \end{aligned}
\]
\end{tcolorbox}

\vspace{1.5em}
\hrule
\vspace{1.5em}

\section*{Question 141: Basics of Electromagnetics \& Maxwell's Eqns (GATE EC 2008)}
\addcontentsline{toc}{section}{Question 141: Basics of Electromagnetics \& Maxwell's Eqns (GATE EC 2008)}
\noindent \textbf{\color{primary}MCQ} \hfill \textbf{\color{secondary}1 Mark}


\noindent For static electric and magnetic fields in an inhomogeneous source-free medium, which of the following represents the correct form of Maxwell's equations ?


\begin{enumerate}[label=(\Alph*)]
    \item $\bigtriangledown \cdot E=0,\bigtriangledown \times B=0$
    \item $\bigtriangledown \cdot E=0,\bigtriangledown \cdot B=0$
    \item $\bigtriangledown  \times  E=0,\bigtriangledown  \times B=0$
    \item \textbf{(Correct)} $\bigtriangledown  \times  E=0,\bigtriangledown \cdot B=0$
\end{enumerate}

\noindent \textbf{Correct Answer:} (D)

\begin{tcolorbox}[solbox, title=Detailed Solution -- Question 141]
Intensity - curl - line integral 
 Density divergence - surface integral
 This is order of understanding Maxwell's equations. 
 For any material

\[
\begin{array}{cc} \nabla \times E=0 & \nabla \cdot B=0 \\ \nabla \times H=J & \nabla \cdot D=\rho_{v} \end{array}
\]

 For source free like charge or current free conditions
$\nabla \times E=0 \quad \nabla \cdot B=0$
$\nabla \times H=0 \quad \nabla \cdot D=0$
 If $D=\epsilon$ E and $B=\mu$ H and $\in$ or $\mu$ are constant with  distance or direction then

\[
\begin{aligned} \nabla \cdot(\in E) &=0 & & \nabla \times \frac{B}{\mu}=0 \\ \nabla \cdot E &=0 & & \nabla \times B=0 \end{aligned}
\]

 These are called as homogeneous/isotropic medium. 
 For non-homogeneous material. However they are not valid.
\end{tcolorbox}

\vspace{1.5em}
\hrule
\vspace{1.5em}

\section*{Question 142: Antennas \& Radiating Systems (GATE EC 2008)}
\addcontentsline{toc}{section}{Question 142: Antennas \& Radiating Systems (GATE EC 2008)}
\noindent \textbf{\color{primary}MCQ} \hfill \textbf{\color{secondary}1 Mark}


\noindent For a Hertz dipole antenna, the half power beam width (HPBW) in the E-plane is


\begin{enumerate}[label=(\Alph*)]
    \item $360^{\circ}$
    \item $180^{\circ}$
    \item \textbf{(Correct)} $90^{\circ}$
    \item $45^{\circ}$
\end{enumerate}

\noindent \textbf{Correct Answer:} (C)

\begin{tcolorbox}[solbox, title=Detailed Solution -- Question 142]
E-plane mean $\theta$-plane or vertical plane 
$E \propto \sin \theta$ in Hertzian dipole

\[
\begin{array}{l} E_{\max } \text { when } \theta=90^{\circ} \\ E_{\min } \text { when } \theta=0 \end{array}
\]

 E has half power of maximum when $\theta=45^{\circ} \; and\;  180^{\circ}$
$\text { HPBW in } \theta \text { side }=135^{\circ}-45^{\circ}=90^{\circ}$
\end{tcolorbox}

\vspace{1.5em}
\hrule
\vspace{1.5em}

\section*{Question 143: Uniform Plane Waves (GATE EC 2007)}
\addcontentsline{toc}{section}{Question 143: Uniform Plane Waves (GATE EC 2007)}
\noindent \textbf{\color{primary}MCQ} \hfill \textbf{\color{secondary}1 Mark}


\noindent A right circularly polarized (RCP) plane wave is incident at an angle $60^{\circ}$ to the
normal, on an air-dielectric interface. If the reflected wave is linearly polarized,
the relative dielectric constant $\varepsilon_{r2}$
is.

\begin{center}\includegraphics[max width=0.85\linewidth]{images/img\_60.jpg}\end{center}

\begin{center}\includegraphics[max width=0.85\linewidth]{images/img\_60.jpg}\end{center}


\begin{enumerate}[label=(\Alph*)]
    \item $\sqrt{2}$
    \item $\sqrt{3}$
    \item 2
    \item \textbf{(Correct)} 3
\end{enumerate}

\noindent \textbf{Correct Answer:} (D)

\begin{tcolorbox}[solbox, title=Detailed Solution -- Question 143]
The reflected wave is linearly polarized when circularly polarized wave is incident under Brewster's angle of incidence.

\[
\begin{aligned} \tan \theta_{B} &=\sqrt{\frac{\epsilon_{r_{2}}}{\epsilon_{r_{1}}}}=\tan 60^{\circ} \\ \epsilon_{r_{2}} &=3 \end{aligned}
\]
\end{tcolorbox}

\vspace{1.5em}
\hrule
\vspace{1.5em}

\section*{Question 144: Antennas \& Radiating Systems (GATE EC 2007)}
\addcontentsline{toc}{section}{Question 144: Antennas \& Radiating Systems (GATE EC 2007)}
\noindent \textbf{\color{primary}MCQ} \hfill \textbf{\color{secondary}1 Mark}


\noindent A $\frac{\lambda }{2}$ dipole is kept horizontally at a height of $\frac{\lambda _{0}}{2}$ above a perfectly conducting infinite ground plane. The radiation pattern in the lane of the dipole ($\vec{E}$ plane) looks approximately as 

\begin{center}\includegraphics[max width=0.85\linewidth]{images/img\_87.jpg}\end{center}

\begin{center}\includegraphics[max width=0.85\linewidth]{images/img\_87.jpg}\end{center}


\begin{enumerate}[label=(\Alph*)]
    \item A
    \item \textbf{(Correct)} B
    \item C
    \item D
\end{enumerate}

\noindent \textbf{Correct Answer:} (B)

\begin{tcolorbox}[solbox, title=Detailed Solution -- Question 144]
\begin{center}\includegraphics[max width=0.85\linewidth]{images/img\_9.jpg}\end{center}

\begin{center}\includegraphics[max width=0.85\linewidth]{images/img\_9.jpg}\end{center}

A horizontally placed dipole has an image in the
earth and interference with the image formed gives
the final pattern of radiation 

\begin{center}\includegraphics[max width=0.85\linewidth]{images/img\_29.jpg}\end{center}

\begin{center}\includegraphics[max width=0.85\linewidth]{images/img\_29.jpg}\end{center}

Dipole pattern

\begin{center}\includegraphics[max width=0.85\linewidth]{images/img\_17.jpg}\end{center}

\begin{center}\includegraphics[max width=0.85\linewidth]{images/img\_17.jpg}\end{center}

Array factor between $D_{1}$ and  $D_{2}$

\[
\begin{array}{l}
\psi=\alpha+\beta d \cos \theta \\
d=\lambda, \lambda=180^{\circ} \text { (image is out of phase) } \\
\psi=\pi+2 \pi \cos \theta
\end{array}
\]

Array pattern

\begin{center}\includegraphics[max width=0.85\linewidth]{images/img\_37.jpg}\end{center}

\begin{center}\includegraphics[max width=0.85\linewidth]{images/img\_37.jpg}\end{center}

Resultant pattern can be obtained by multiplication
of patterns.
\end{tcolorbox}

\vspace{1.5em}
\hrule
\vspace{1.5em}

\section*{Question 145: Transmission Lines (GATE EC 2007)}
\addcontentsline{toc}{section}{Question 145: Transmission Lines (GATE EC 2007)}
\noindent \textbf{\color{primary}MCQ} \hfill \textbf{\color{secondary}1 Mark}


\noindent The parallel branches of a 2-wire transmission line re terminated in 100$\Omega$ and
200$\Omega$ resistors as shown in the figure. The characteristic impedance of the line is $Z_{0} =  50\Omega$ and each section has a length of $\frac{\lambda }{4}$. The voltage reflection coefficient $\Gamma$ at the input is

\begin{center}\includegraphics[max width=0.85\linewidth]{images/img\_27.jpg}\end{center}

\begin{center}\includegraphics[max width=0.85\linewidth]{images/img\_27.jpg}\end{center}


\begin{enumerate}[label=(\Alph*)]
    \item $-j\frac{7}{5}$
    \item $\frac{-5}{7}$
    \item $j \frac{5}{7}$
    \item \textbf{(Correct)} $\frac{5}{7}$
\end{enumerate}

\noindent \textbf{Correct Answer:} (D)

\begin{tcolorbox}[solbox, title=Detailed Solution -- Question 145]
\begin{center}\includegraphics[max width=0.85\linewidth]{images/img\_61.jpg}\end{center}

\begin{center}\includegraphics[max width=0.85\linewidth]{images/img\_61.jpg}\end{center}

\[
\begin{aligned} Z_{\text {in }}&=\frac{Z_{0}^{2}}{Z_{L}} ; if l=\frac{\lambda}{4} \\ R_{1} \text{ due to }100 \Omega&=\frac{50^{2}}{100}=25 \\ R_{2} \text{ due to }200 \Omega&=\frac{50^{2}}{200}=\frac{25}{2} \\ \Rightarrow \quad R_{1} \| R_{2} &=25 \text { II } \frac{25}{2}=\frac{25}{3} \\ Z_{S} &=\frac{(50)^{2}}{25 / 3}=300 \Omega \\ \therefore \quad\Gamma&=\frac{Z_{S}-Z_{0}}{Z_{S}+Z_{0}}=\frac{300-50}{300+50}=\frac{5}{7} \end{aligned}
\]
\end{tcolorbox}

\vspace{1.5em}
\hrule
\vspace{1.5em}

\section*{Question 146: miscellaneous Topics (GATE EC 2007)}
\addcontentsline{toc}{section}{Question 146: miscellaneous Topics (GATE EC 2007)}
\noindent \textbf{\color{primary}MCQ} \hfill \textbf{\color{secondary}1 Mark}


\noindent A load of 50$\Omega$ is connected in shunt in a 2-wire transmission line of $Z_{0} = 50\Omega$ as shown in the figure. The 2-port scattering parameter matrix (s-matrix) of the
shunt element is  

\begin{center}\includegraphics[max width=0.85\linewidth]{images/img\_44.jpg}\end{center}

\begin{center}\includegraphics[max width=0.85\linewidth]{images/img\_44.jpg}\end{center}


\begin{enumerate}[label=(\Alph*)]
    \item \[
\begin{bmatrix} -\frac{1}{2} &\frac{1}{2} \\ \frac{1}{2} & -\frac{1}{2} \end{bmatrix}
\]
    \item \[
\begin{bmatrix} 0 & 1\\ 1&0 \end{bmatrix}
\]
    \item \textbf{(Correct)} \[
\begin{bmatrix} -\frac{1}{3} &\frac{2}{3} \\ \frac{2}{3} & -\frac{1}{3} \end{bmatrix}
\]
    \item \[
\begin{bmatrix} \frac{1}{4} &-\frac{3}{4} \\ -\frac{3}{4} & \frac{1}{4} \end{bmatrix}
\]
\end{enumerate}

\noindent \textbf{Correct Answer:} (C)

\begin{tcolorbox}[solbox, title=Detailed Solution -- Question 146]
The line is terminated with $50 \Omega$ at the center and
so, matched on both the sides.
\end{tcolorbox}

\vspace{1.5em}
\hrule
\vspace{1.5em}

\section*{Question 147: Waveguides \& Optical Fibers (GATE EC 2007)}
\addcontentsline{toc}{section}{Question 147: Waveguides \& Optical Fibers (GATE EC 2007)}
\noindent \textbf{\color{primary}MCQ} \hfill \textbf{\color{secondary}1 Mark}


\noindent The $\vec{E}$ field in a rectangular waveguide of inner dimension $a \times b$
is given by 

\[
\vec{E}=\frac{\omega \mu }{h^{2}}(\frac{\pi}{a})H_{0}sin(\frac{2\pi x}{a})^{2}sin(\omega t-\beta z)\hat{y}
\]
 
where  $H_{0}$ is a constant, and a and b are the dimensions along the x-axis and the y-axis respectively. The mode of propagation in the waveguide is


\begin{enumerate}[label=(\Alph*)]
    \item \textbf{(Correct)} $TE_{20}$
    \item $TM_{11}$
    \item $TM_{20}$
    \item $TE_{10}$
\end{enumerate}

\noindent \textbf{Correct Answer:} (A)

\begin{tcolorbox}[solbox, title=Detailed Solution -- Question 147]
The wave is E(x,z,t) with no propagation in y-direction.
 This is possible only in TE waves when $n=0$
$\left(\cos \frac{n \pi}{b} y=1\right)$
\end{tcolorbox}

\vspace{1.5em}
\hrule
\vspace{1.5em}

\section*{Question 148: Uniform Plane Waves (GATE EC 2007)}
\addcontentsline{toc}{section}{Question 148: Uniform Plane Waves (GATE EC 2007)}
\noindent \textbf{\color{primary}MCQ} \hfill \textbf{\color{secondary}1 Mark}


\noindent The $\vec{H}$ field (in A/m) of a plane wave propagating in free space is given by 
 
\[
\vec{H}=\hat{x}\frac{5\sqrt{3}}{\eta _{0}}cos(\omega t-\beta z)+\hat{y}\frac{5}{\eta _{0}}sin(\omega t-\beta z+\frac{\pi }{2})
\]
.  
The time average power flow density in Watts is


\begin{enumerate}[label=(\Alph*)]
    \item $\frac{\eta _{0}}{100}$
    \item $\frac{100}{\eta _{0}}$
    \item $50 {\eta _{0}}^{2}$
    \item \textbf{(Correct)} $\frac{50}{\eta _{0}}$
\end{enumerate}

\noindent \textbf{Correct Answer:} (D)

\begin{tcolorbox}[solbox, title=Detailed Solution -- Question 148]
For free space,

\[
\begin{aligned} P&=\frac{E^{2}}{2 \eta_{0}}, E=\eta_{0} H \\ P&=\frac{\eta_{0}^{2} H^{2}}{2 \eta_{0}}=\frac{\eta_{0} H^{2}}{2} \\ H&=\frac{1}{\eta_{0}} \sqrt{(5 \sqrt{3})^{2}+(5)^{2}}=\frac{10}{\eta_{0}} \\ P&=\frac{\eta_{0}}{2} \times \frac{100}{\eta_{0}^{2}}=\frac{50}{\eta_{0}} \text { watts } \end{aligned}
\]
\end{tcolorbox}

\vspace{1.5em}
\hrule
\vspace{1.5em}

\section*{Question 149: Waveguides \& Optical Fibers (GATE EC 2007)}
\addcontentsline{toc}{section}{Question 149: Waveguides \& Optical Fibers (GATE EC 2007)}
\noindent \textbf{\color{primary}MCQ} \hfill \textbf{\color{secondary}1 Mark}


\noindent An air-filled rectangular waveguide has inner dimensions of 3cm x 2cm.
The wave impedance of the TE20  mode of propagation in the waveguide at a
frequency of 30 GHz is (free space impedance $\eta_{0}$ = 377 $\Omega$)


\begin{enumerate}[label=(\Alph*)]
    \item $308\Omega$
    \item $355\Omega$
    \item \textbf{(Correct)} $400\Omega$
    \item $461\Omega$
\end{enumerate}

\noindent \textbf{Correct Answer:} (C)

\begin{tcolorbox}[solbox, title=Detailed Solution -- Question 149]
\[
\begin{aligned} \text { For } TE_{20}, \quad \lambda_{C}&=a=3 \times 10^{-2} \mathrm{m} \\ f_{C}&= \frac{C}{\lambda_{C}}=\frac{3 \times 10^{8}}{3 \times 10^{-2}}=10^{10} \mathrm{Hz} \\ \eta_{T E}&=\frac{\eta}{\sqrt{1-\left(\frac{f_{c}}{f}\right)^{2}}} =\frac{377}{\sqrt{1-\left(\frac{10^{10}}{3 \times 10^{10}}\right)^{2}}}\\ \eta_{T E}&=400 \Omega \end{aligned}
\]
\end{tcolorbox}

\vspace{1.5em}
\hrule
\vspace{1.5em}

\section*{Question 150: Basics of Electromagnetics \& Maxwell's Eqns (GATE EC 2007)}
\addcontentsline{toc}{section}{Question 150: Basics of Electromagnetics \& Maxwell's Eqns (GATE EC 2007)}
\noindent \textbf{\color{primary}MCQ} \hfill \textbf{\color{secondary}1 Mark}


\noindent If C is closed curve enclosing a surface $S$, then the magnetic field intensity $\vec{H}$, the current density $\vec{J}$ and the electric flux density $\vec{D}$ are related by


\begin{enumerate}[label=(\Alph*)]
    \item \[
\int \int_{S} \vec{H}\cdot \vec{ds}= \oint_{c}(\vec{j}+\frac{\partial \vec{D}}{\partial t})\cdot d\vec{l}
\]
    \item \[
\int_{S} \vec{H}\cdot \vec{dl}=\oint \oint_{S}(\vec{j}+\frac{\partial \vec{D}}{\partial t})\cdot d\vec{s}
\]
    \item \[
\oint \oint_{S} \vec{H}\cdot \vec{ds}=\int_{C}(\vec{j}+\frac{\partial \vec{D}}{\partial t}) \cdot d\vec{l}
\]
    \item \textbf{(Correct)} \[
\oint_{C} \vec{H}\cdot \vec{dl}\oint_{c}=\int \int_{S}(\vec{j}+\frac{\partial \vec{D}}{\partial t})\cdot d\vec{s}
\]
\end{enumerate}

\noindent \textbf{Correct Answer:} (D)


\vspace{1.5em}
\hrule
\vspace{1.5em}

\section*{Question 151: Uniform Plane Waves (GATE EC 2007)}
\addcontentsline{toc}{section}{Question 151: Uniform Plane Waves (GATE EC 2007)}
\noindent \textbf{\color{primary}MCQ} \hfill \textbf{\color{secondary}1 Mark}


\noindent A plane wave of wavelength $\lambda$
is traveling in a direction making an angle $30^{\circ}$  with positive x-axis and $90^{\circ}$ with positive y-axis. The $\vec{E}$ field of the plane
wave can be represented as ($E_0$  is constant)


\begin{enumerate}[label=(\Alph*)]
    \item \textbf{(Correct)} \[
\vec{E}=\hat{y}E_0 e^{j(\omega t-\frac{\sqrt{3}\pi}{\lambda}x-\frac{\pi }{\lambda }z)}
\]
    \item \[
\vec{E}=\hat{y}E_0 e^{j(\omega t-\frac{\pi }{\lambda }x-\frac{\sqrt{3}\pi }{\lambda }z)}
\]
    \item \[
\vec{E}=\hat{y}E_0 e^{j(\omega t+\frac{\sqrt{3}\pi}{\lambda }x+\frac{\pi }{\lambda }z)}
\]
    \item \[
\vec{E}=\hat{y}E_0 e^{j(\omega t-\frac{\pi }{\lambda }x+\frac{\sqrt{3}\pi }{\lambda }z)}
\]
\end{enumerate}

\noindent \textbf{Correct Answer:} (A)

\begin{tcolorbox}[solbox, title=Detailed Solution -- Question 151]
\begin{center}\includegraphics[max width=0.85\linewidth]{images/img\_3.jpg}\end{center}

\begin{center}\includegraphics[max width=0.85\linewidth]{images/img\_3.jpg}\end{center}

Wave is travelling in the ZX-plane $\left(90^{\circ} \text { to Y-axis} \right)$ and $30^{\circ}$ to X-axis.

\[
\begin{array}{l} \beta_{x}=\beta \cos 30^{\circ}=\frac{2 \pi}{\lambda} \cdot \frac{\sqrt{3}}{2}=\frac{\sqrt{3} \pi}{\lambda} \\ \beta_{z}=\beta \sin 30^{\circ}=\frac{\pi}{\lambda} \end{array}
\]

 The two solutions possible are shown as $P_{1}$ and $P_{2}$ paths in the diagrams. Option (A) is $P_{2}$
\end{tcolorbox}

\vspace{1.5em}
\hrule
\vspace{1.5em}

\section*{Question 152: Transmission Lines (GATE EC 2006)}
\addcontentsline{toc}{section}{Question 152: Transmission Lines (GATE EC 2006)}
\noindent \textbf{\color{primary}MCQ} \hfill \textbf{\color{secondary}1 Mark}


\noindent A 30-Volts battery with zero source resistance is connected to a coaxial line of
characteristic impedance of 50 Ohms at t=0 second terminated in an unknown resistive load. The line length is that it takes 400 $\mu$s for an electromagnetic wave to travel from source end to load end and vice-versa. At  t=400 $\mu$s,  the voltage at the load end is
found to be 40 Volts. 

 The steady-state current through the load resistance is


\begin{enumerate}[label=(\Alph*)]
    \item 1.2 Amps
    \item \textbf{(Correct)} 0.3 Amps
    \item 0.6 Amps
    \item 0.4 Amps
\end{enumerate}

\noindent \textbf{Correct Answer:} (B)

\begin{tcolorbox}[solbox, title=Detailed Solution -- Question 152]
In steady state, $V_{L}=30 \mathrm{V}$
$I_{L}=\frac{30}{100}=0.3 \mathrm{A}$
\end{tcolorbox}

\vspace{1.5em}
\hrule
\vspace{1.5em}

\section*{Question 153: Transmission Lines (GATE EC 2006)}
\addcontentsline{toc}{section}{Question 153: Transmission Lines (GATE EC 2006)}
\noindent \textbf{\color{primary}MCQ} \hfill \textbf{\color{secondary}1 Mark}


\noindent A 30-Volts battery with zero source resistance is connected to a coaxial line of
characteristic impedance of 50 Ohms at t=0 second terminated in an unknown resistive load. The line length is that it takes 400 $\mu$s for an electromagnetic wave to travel from source end to load end and vice-versa. At  t=400 $\mu$s,  the voltage at the load end is
found to be 40 Volts. 

 The load resistance is


\begin{enumerate}[label=(\Alph*)]
    \item 25 Ohms
    \item 50 Ohms
    \item 75 Ohms
    \item \textbf{(Correct)} 100 Ohms
\end{enumerate}

\noindent \textbf{Correct Answer:} (D)

\begin{tcolorbox}[solbox, title=Detailed Solution -- Question 153]
\begin{center}\includegraphics[max width=0.85\linewidth]{images/img\_79.jpg}\end{center}

\begin{center}\includegraphics[max width=0.85\linewidth]{images/img\_79.jpg}\end{center}

10V reflects to add to 30V in phase and gives 40V

\[
\begin{aligned} \Gamma &=\frac{V_{r}}{V_{i}}=\frac{10}{30}=\frac{1}{3} \\ \Gamma &=\frac{Z_{L}-Z_{0}}{Z_{L}+Z_{0}} \\ \Rightarrow \quad \frac{1}{3} &=\frac{Z_{L}-50}{Z_{L}+50} \\ \Rightarrow \quad Z_{L} &=100 \Omega \end{aligned}
\]
\end{tcolorbox}

\vspace{1.5em}
\hrule
\vspace{1.5em}

\section*{Question 154: Antennas \& Radiating Systems (GATE EC 2006)}
\addcontentsline{toc}{section}{Question 154: Antennas \& Radiating Systems (GATE EC 2006)}
\noindent \textbf{\color{primary}MCQ} \hfill \textbf{\color{secondary}1 Mark}


\noindent A mast antenna consisting of a 50 meter long vertical conductor operates over a
perfectly conducting ground plane. It is base-fed at a frequency of 600 kHz. The
radiation resistance of the antenna is Ohms is


\begin{enumerate}[label=(\Alph*)]
    \item \textbf{(Correct)} $\frac{2\pi ^{2}}{5}$
    \item $\frac{\pi ^{2}}{5}$
    \item $\frac{4\pi ^{2}}{5}$
    \item $20 \pi ^{2}$
\end{enumerate}

\noindent \textbf{Correct Answer:} (A)

\begin{tcolorbox}[solbox, title=Detailed Solution -- Question 154]
A mast antenna is a tower antenna or monopole  vertically erected on the ground for low frequency  ground wave communication.
$l=50 \mathrm{m}, \quad \lambda=500$
$l=\frac{\lambda}{10}$electrically short monopole
$R_{r}=20 \pi^{2}\left(\frac{d l}{\lambda}\right)^{2} \times 2$
 This 2 in the multiplication is due to image in conducting earth.
$R_{r}=\frac{2 \pi^{2}}{5} \simeq 4 \Omega$
\end{tcolorbox}

\vspace{1.5em}
\hrule
\vspace{1.5em}

\section*{Question 155: Waveguides \& Optical Fibers (GATE EC 2006)}
\addcontentsline{toc}{section}{Question 155: Waveguides \& Optical Fibers (GATE EC 2006)}
\noindent \textbf{\color{primary}MCQ} \hfill \textbf{\color{secondary}1 Mark}


\noindent A rectangular wave guide having $TE_{10}$ mode as dominant mode is having a cut off frequency 18-GHz for the mode $TE_{30}$. The inner broad-wall dimension of the rectangular wave guide is


\begin{enumerate}[label=(\Alph*)]
    \item $\frac{5}{3}cm$
    \item 5 cm
    \item \textbf{(Correct)} $\frac{5}{2}cm$
    \item 10 cm
\end{enumerate}

\noindent \textbf{Correct Answer:} (C)

\begin{tcolorbox}[solbox, title=Detailed Solution -- Question 155]
\[
\begin{aligned} \mathrm{TE}_{30^{\prime}} \quad \lambda_{c}&=\frac{2 a}{m}=\frac{2 a}{3} \\ \lambda_{c} &=\frac{c}{t}=\frac{3 \times 10^{8}}{18 \times 10^{9}}=\frac{1}{60} \\ \frac{1}{60} &=\frac{2 a}{3} \\ \Rightarrow \quad a &=\frac{1}{40}=2.5 \mathrm{cm}=\frac{5}{2} \mathrm{cm} \end{aligned}
\]
\end{tcolorbox}

\vspace{1.5em}
\hrule
\vspace{1.5em}

\section*{Question 156: Uniform Plane Waves (GATE EC 2006)}
\addcontentsline{toc}{section}{Question 156: Uniform Plane Waves (GATE EC 2006)}
\noindent \textbf{\color{primary}MCQ} \hfill \textbf{\color{secondary}1 Mark}


\noindent When a planes wave traveling in free-space is incident normally on a medium
having $\varepsilon_{r}=4.0$, the fraction of power transmitted into the medium is given by


\begin{enumerate}[label=(\Alph*)]
    \item \textbf{(Correct)} $\frac{8}{9}$
    \item $\frac{1}{2}$
    \item $\frac{1}{3}$
    \item $\frac{5}{6}$
\end{enumerate}

\noindent \textbf{Correct Answer:} (A)

\begin{tcolorbox}[solbox, title=Detailed Solution -- Question 156]
\[
\begin{aligned} P_{t}&=\left(1-\Gamma^{2}\right) P_{t} \\ \Gamma&=\frac{\eta_{2}-\eta_{1}}{\eta_{2}+\eta_{1}}=\frac{\sqrt{\frac{\mu_{0}}{\epsilon_{0} \epsilon_{r}}}-\sqrt{\frac{\mu_{0}}{\epsilon_{0}}}}{\sqrt{\frac{\mu_{0}}{\epsilon_{0} \epsilon_{r}}}+\sqrt{\frac{\mu_{0}}{\epsilon_{0}}}} \\ \Gamma&=\frac{\frac{1}{2}-1}{\frac{1}{2}+1}=\frac{-1}{3} \\ \tau&=1-\Gamma^{2}=1-\frac{1}{9}=\frac{8}{9} \\ \therefore \quad \quad P_{t}&=\tau \cdot P_{i} \\ \frac{P_{t}}{P_{i}}&=\frac{8}{9} \end{aligned}
\]
\end{tcolorbox}

\vspace{1.5em}
\hrule
\vspace{1.5em}

\section*{Question 157: Basics of Electromagnetics \& Maxwell's Eqns (GATE EC 2006)}
\addcontentsline{toc}{section}{Question 157: Basics of Electromagnetics \& Maxwell's Eqns (GATE EC 2006)}
\noindent \textbf{\color{primary}MCQ} \hfill \textbf{\color{secondary}1 Mark}


\noindent A medium is divide into regions I and II about x = 0 plane, as shown in the
figure below. 

\begin{center}\includegraphics[max width=0.85\linewidth]{images/img\_34.jpg}\end{center}

\begin{center}\includegraphics[max width=0.85\linewidth]{images/img\_34.jpg}\end{center}
 
An electromagnetic wave with electric field $E_{1}=4\hat{a}_{x}+3\hat{a}_{y}+5\hat{a}_{z}$ is incident normally on the interface from region I . The electric file $E_2$ in region II at the interface is


\begin{enumerate}[label=(\Alph*)]
    \item $E_{2}=E_{1}$
    \item $4\hat{a}_{x}+0.75\hat{a}_{y}-1.25\hat{a}_{z}$
    \item \textbf{(Correct)} $3\hat{a}_{x}+3\hat{a}_{y}+5\hat{a}_{z}$
    \item $-3\hat{a}_{x}+3\hat{a}_{y}+5\hat{a}_{z}$
\end{enumerate}

\noindent \textbf{Correct Answer:} (C)

\begin{tcolorbox}[solbox, title=Detailed Solution -- Question 157]
\[
\begin{aligned} \vec{E}_{t_{1}} &=3 \hat{a}_{y}+5 \hat{a}_{z}=\vec{E}_{t_{2}} \\ \vec{D}_{n_{1}} &=\vec{D}_{n_{2}} \\ \epsilon_{2} \vec{E}_{n_{2}} &=\epsilon_{1} \vec{E}_{n_{1}} \\ \vec{E}_{n_{2}} &=\frac{3 \times 4 \hat{a}_{x}}{4}=3 \hat{a}_{x} \\ \therefore \quad \vec{E}_{2} &=\vec{E}_{t_{2}}+\vec{E}_{n_{2}} \\ &=3 \hat{a}_{y}+5 \hat{a}_{z}+3 \hat{a}_{x} \end{aligned}
\]
\end{tcolorbox}

\vspace{1.5em}
\hrule
\vspace{1.5em}

\section*{Question 158: Uniform Plane Waves (GATE EC 2006)}
\addcontentsline{toc}{section}{Question 158: Uniform Plane Waves (GATE EC 2006)}
\noindent \textbf{\color{primary}MCQ} \hfill \textbf{\color{secondary}1 Mark}


\noindent A medium of relative permittivity $\varepsilon _{r2}=2$ forms an interface with free-space. A point source of electromagnetic energy is located in the medium at a depth of 1 meter from the interface. Due to the total internal reflection, the transmitted beam has a circular cross-section over the interface. The area of the beam crosssection at the interface is given by


\begin{enumerate}[label=(\Alph*)]
    \item $2\pi m^{2}$
    \item $\pi^{2} m^{2}$
    \item $\frac{\pi }{2} m^{2}$
    \item \textbf{(Correct)} $\pi m^{2}$
\end{enumerate}

\noindent \textbf{Correct Answer:} (D)

\begin{tcolorbox}[solbox, title=Detailed Solution -- Question 158]
\begin{center}\includegraphics[max width=0.85\linewidth]{images/img\_2.jpg}\end{center}

\begin{center}\includegraphics[max width=0.85\linewidth]{images/img\_2.jpg}\end{center}

Circular cross-section beam for rays incident 
 below critical angle. 
 Calculating critical angle $\sin \theta_{c}=\sqrt{\frac{\epsilon_{o}}{2 \epsilon_{o}}}$
$\theta_{c}=45^{\circ}$
$1 \mathrm{m}-\mathrm{depth} ; 1 \mathrm{m}-\mathrm{radius}$
$\text { Area }=\pi R^{2}=\pi$
\end{tcolorbox}

\vspace{1.5em}
\hrule
\vspace{1.5em}

\section*{Question 159: Antennas \& Radiating Systems (GATE EC 2006)}
\addcontentsline{toc}{section}{Question 159: Antennas \& Radiating Systems (GATE EC 2006)}
\noindent \textbf{\color{primary}MCQ} \hfill \textbf{\color{secondary}1 Mark}


\noindent A transmission line is feeding 1 watt of power to a horn antenna having a gain
of 10 dB. The antenna is matched to the transmission line. The total power
radiated by the horn antenna into the free space is


\begin{enumerate}[label=(\Alph*)]
    \item 10 Watts
    \item \textbf{(Correct)} 1 Watts
    \item 0.1 Watts
    \item 0.01 Watts
\end{enumerate}

\noindent \textbf{Correct Answer:} (B)

\begin{tcolorbox}[solbox, title=Detailed Solution -- Question 159]
Gain of antenna is directive gain i.e. different from
that of amplifier. It radiates same power. So 1 W.
\end{tcolorbox}

\vspace{1.5em}
\hrule
\vspace{1.5em}

\section*{Question 160: Uniform Plane Waves (GATE EC 2006)}
\addcontentsline{toc}{section}{Question 160: Uniform Plane Waves (GATE EC 2006)}
\noindent \textbf{\color{primary}MCQ} \hfill \textbf{\color{secondary}1 Mark}


\noindent The electric field of an electromagnetic wave propagation in the positive
direction is given by 
 $E=\hat{a}_{x}sin(\omega t-\beta z)+\hat{a}_{y}sin(\omega t-\beta z+\pi /2)$. 
 The wave is


\begin{enumerate}[label=(\Alph*)]
    \item Linearly polarized in the z -direction
    \item Elliptically polarized
    \item \textbf{(Correct)} Left-hand circularly polarized
    \item Right-hand circularly polarized
\end{enumerate}

\noindent \textbf{Correct Answer:} (C)

\begin{tcolorbox}[solbox, title=Detailed Solution -- Question 160]
Equal amplitudes and out of phase by $90^{\circ}$, wave
is circularly polarized. Tracing the time advancing
E on the z = 0 plane at t= 0, t= T/4, t = T/2 gives
left/right notations
\end{tcolorbox}

\vspace{1.5em}
\hrule
\vspace{1.5em}

\section*{Question 161: Transmission Lines (GATE EC 2005)}
\addcontentsline{toc}{section}{Question 161: Transmission Lines (GATE EC 2005)}
\noindent \textbf{\color{primary}MCQ} \hfill \textbf{\color{secondary}1 Mark}


\noindent Voltage standing wave pattern in a lossless transmission line with characteristic
impedance 50$\Omega$ and a resistive load is shown in figure. 

\begin{center}\includegraphics[max width=0.85\linewidth]{images/img\_93.jpg}\end{center}

\begin{center}\includegraphics[max width=0.85\linewidth]{images/img\_93.jpg}\end{center}
 
 The reflection coefficient is given by


\begin{enumerate}[label=(\Alph*)]
    \item \textbf{(Correct)} -0.6
    \item -1
    \item 0.6
    \item 0
\end{enumerate}

\noindent \textbf{Correct Answer:} (A)

\begin{tcolorbox}[solbox, title=Detailed Solution -- Question 161]
$\Gamma=\frac{Z_{L}-Z_{0}}{Z_{L}+Z_{0}}=\frac{12.5-50}{12.5+50}=-0.6$
\end{tcolorbox}

\vspace{1.5em}
\hrule
\vspace{1.5em}

\section*{Question 162: Transmission Lines (GATE EC 2005)}
\addcontentsline{toc}{section}{Question 162: Transmission Lines (GATE EC 2005)}
\noindent \textbf{\color{primary}MCQ} \hfill \textbf{\color{secondary}1 Mark}


\noindent Voltage standing wave pattern in a lossless transmission line with characteristic
impedance 50$\Omega$ and a resistive load is shown in figure. 

\begin{center}\includegraphics[max width=0.85\linewidth]{images/img\_93.jpg}\end{center}

\begin{center}\includegraphics[max width=0.85\linewidth]{images/img\_93.jpg}\end{center}
 
 The value of the load resistance is


\begin{enumerate}[label=(\Alph*)]
    \item 50 $\Omega$
    \item 200 $\Omega$
    \item \textbf{(Correct)} 12.5 $\Omega$
    \item 0
\end{enumerate}

\noindent \textbf{Correct Answer:} (C)

\begin{tcolorbox}[solbox, title=Detailed Solution -- Question 162]
\[
\begin{aligned} S &=\frac{V_{\max }}{V_{\min }}=\frac{4}{1}=4 \\ S &=\frac{1+\Gamma}{1-\Gamma} \\ Z_{\max } &=Z_{0} \cdot S \\ Z_{\min } &=Z_{0} / S \end{aligned}
\]

 The digram shows that minima is at load
$\therefore \quad Z_{L}=Z_{\min }=\frac{50}{4}=12.5 \Omega$
\end{tcolorbox}

\vspace{1.5em}
\hrule
\vspace{1.5em}

\section*{Question 163: Antennas \& Radiating Systems (GATE EC 2005)}
\addcontentsline{toc}{section}{Question 163: Antennas \& Radiating Systems (GATE EC 2005)}
\noindent \textbf{\color{primary}MCQ} \hfill \textbf{\color{secondary}1 Mark}


\noindent Two identical and parallel dipole antennas are kept apart by a distance of $\frac{\lambda }{4}$ in the H-plane. They are fed with equal currents but the right most antenna has a phase shift of +$90^{\circ}$. The radiation pattern is given as.

\begin{center}\includegraphics[max width=0.85\linewidth]{images/img\_84.jpg}\end{center}

\begin{center}\includegraphics[max width=0.85\linewidth]{images/img\_84.jpg}\end{center}


\begin{enumerate}[label=(\Alph*)]
    \item A
    \item B
    \item \textbf{(Correct)} C
    \item D
\end{enumerate}

\noindent \textbf{Correct Answer:} (C)

\begin{tcolorbox}[solbox, title=Detailed Solution -- Question 163]
H-plane means $\phi$ pattern. 
 Both dipoles are horizontally isotropic in the  H-plane.

\[
\begin{aligned} E &=2 E_{0} \cos (\psi / 2) \\ \psi &=\alpha+\beta d \cos \theta \\ &=\frac{\pi}{2}+\frac{2 \pi}{\lambda} \cdot \frac{\lambda}{4} \cos \theta \\ \psi &=\frac{\pi}{2}+\frac{\pi}{2} \cos \theta \\ \psi=0 & \; \cos \theta=-1 & \theta_{\max }=180^{\circ} \\ \psi=\pi &\; \cos \theta=1 & \theta_{N P}=0^{\circ} \\ \psi=\frac{\pi}{2}& \; \cos \theta=0 &\theta_{\text{HPP}=90^{\circ} / 270^{\circ} \end{aligned}}
\]

Option (C) when $\theta$ is considered as shown in array calculations 

\begin{center}\includegraphics[max width=0.85\linewidth]{images/img\_36.jpg}\end{center}

\begin{center}\includegraphics[max width=0.85\linewidth]{images/img\_36.jpg}\end{center}
\end{tcolorbox}

\vspace{1.5em}
\hrule
\vspace{1.5em}

\section*{Question 164: Transmission Lines (GATE EC 2005)}
\addcontentsline{toc}{section}{Question 164: Transmission Lines (GATE EC 2005)}
\noindent \textbf{\color{primary}MCQ} \hfill \textbf{\color{secondary}1 Mark}


\noindent Characteristic impedance of a transmission line is 50 $\Omega$. Input impedance of the open-circuited line $Z_{oc}=100+j150\Omega$.
when the transmission line a short circuited, then value of the input impedance will be.


\begin{enumerate}[label=(\Alph*)]
    \item 50 $\Omega$
    \item 100+j150 $\Omega$
    \item 7.69+j11.54 $\Omega$
    \item \textbf{(Correct)} 7.69-j11.54 $\Omega$
\end{enumerate}

\noindent \textbf{Correct Answer:} (D)

\begin{tcolorbox}[solbox, title=Detailed Solution -- Question 164]
\[
\begin{aligned} z_{0}^{2}&=z_{o c} \cdot z_{s c} \\ z_{s c}&=\frac{50 \times 50}{100+j 150}=\frac{50}{2+3 j}=\frac{50(2-3 j)}{13} \\ z_{S c}&=7.69- j11.54  \end{aligned}
\]
\end{tcolorbox}

\vspace{1.5em}
\hrule
\vspace{1.5em}

\section*{Question 165: Waveguides \& Optical Fibers (GATE EC 2005)}
\addcontentsline{toc}{section}{Question 165: Waveguides \& Optical Fibers (GATE EC 2005)}
\noindent \textbf{\color{primary}MCQ} \hfill \textbf{\color{secondary}1 Mark}


\noindent Which one of the following does represent the electric field lines for the $TE_{02}$ mode in
the cross-section of a hollow rectangular metallic waveguide ?

\begin{center}\includegraphics[max width=0.85\linewidth]{images/img\_96.jpg}\end{center}

\begin{center}\includegraphics[max width=0.85\linewidth]{images/img\_96.jpg}\end{center}


\begin{enumerate}[label=(\Alph*)]
    \item A
    \item B
    \item C
    \item \textbf{(Correct)} D
\end{enumerate}

\noindent \textbf{Correct Answer:} (D)

\begin{tcolorbox}[solbox, title=Detailed Solution -- Question 165]
Two maximas in the field along y-direction lndentifies $TE_{02}$
\end{tcolorbox}

\vspace{1.5em}
\hrule
\vspace{1.5em}

\section*{Question 166: Uniform Plane Waves (GATE EC 2005)}
\addcontentsline{toc}{section}{Question 166: Uniform Plane Waves (GATE EC 2005)}
\noindent \textbf{\color{primary}MCQ} \hfill \textbf{\color{secondary}1 Mark}


\noindent Refractive index of glass is 1.5. Find the wavelength of a beam of light with
frequency of $10^{14}$ Hz in glass. Assume velocity of light is 3 $\times$  $10^{8}$ m/s in vacuum


\begin{enumerate}[label=(\Alph*)]
    \item 3 $\mu$m
    \item 3 mm
    \item \textbf{(Correct)} 2 $\mu$m
    \item 1 $\mu$m
\end{enumerate}

\noindent \textbf{Correct Answer:} (C)

\begin{tcolorbox}[solbox, title=Detailed Solution -- Question 166]
$\text { Refractive index }=\sqrt{\frac{\epsilon_{2}}{\epsilon_{1}}}$
 For glass with respective air $=\sqrt{\epsilon_{R}}$

\[
V_{p}=\frac{3 \times 10^{8}}{\sqrt{\epsilon_{R}} \times f}=\frac{3 \times 10^{8}}{1.5 \times 10^{14}}=2 \mu \mathrm{m}
\]
\end{tcolorbox}

\vspace{1.5em}
\hrule
\vspace{1.5em}

\section*{Question 167: Transmission Lines (GATE EC 2005)}
\addcontentsline{toc}{section}{Question 167: Transmission Lines (GATE EC 2005)}
\noindent \textbf{\color{primary}MCQ} \hfill \textbf{\color{secondary}1 Mark}


\noindent Many circles are drawn in a Smith Chart used for transmission line calculations. The circles shown in the figure represent 

\begin{center}\includegraphics[max width=0.85\linewidth]{images/img\_90.jpg}\end{center}

\begin{center}\includegraphics[max width=0.85\linewidth]{images/img\_90.jpg}\end{center}


\begin{enumerate}[label=(\Alph*)]
    \item Unit circles
    \item Constant resistance circles
    \item \textbf{(Correct)} Constant reactance circles
    \item Constant reflection coefficient circles.
\end{enumerate}

\noindent \textbf{Correct Answer:} (C)


\vspace{1.5em}
\hrule
\vspace{1.5em}

\section*{Question 168: Uniform Plane Waves (GATE EC 2005)}
\addcontentsline{toc}{section}{Question 168: Uniform Plane Waves (GATE EC 2005)}
\noindent \textbf{\color{primary}MCQ} \hfill \textbf{\color{secondary}1 Mark}


\noindent The magnetic field intensity vector of a plane wave is given by 
$\bar{H}(x,y,z,t)=10sin(50000t+0.004x+30)\hat{a}_{y}$
 where $\hat{a}_{y}$, denotes the unit vector in y direction. The wave is propagating with a phase velocity.


\begin{enumerate}[label=(\Alph*)]
    \item $5 \times 10^{4}m/s$
    \item $-3 \times 10^{8}m/s$
    \item \textbf{(Correct)} $-1.25 \times 10^{7}m/s$
    \item $3 \times 10^{8}m/s$
\end{enumerate}

\noindent \textbf{Correct Answer:} (C)

\begin{tcolorbox}[solbox, title=Detailed Solution -- Question 168]
\[
\begin{aligned} \omega&=50,000 ; \beta=0.004 \\ v_{p}&=\frac{\omega}{\beta}=\frac{5 \times 10^{4}}{4 \times 10^{-3}}=1.25 \times 10^{7} \mathrm{m} / \mathrm{s} \end{aligned}
\]
\end{tcolorbox}

\vspace{1.5em}
\hrule
\vspace{1.5em}

\section*{Question 169: Transmission Lines (GATE EC 2004)}
\addcontentsline{toc}{section}{Question 169: Transmission Lines (GATE EC 2004)}
\noindent \textbf{\color{primary}MCQ} \hfill \textbf{\color{secondary}1 Mark}


\noindent A lossless transmission line is terminated in a load which reflects a part of the incident power. The measured VSWR is 2. The percentage of the power that is reflected back is


\begin{enumerate}[label=(\Alph*)]
    \item 57.73
    \item 33.33
    \item 0.11
    \item \textbf{(Correct)} 11.11
\end{enumerate}

\noindent \textbf{Correct Answer:} (D)

\begin{tcolorbox}[solbox, title=Detailed Solution -- Question 169]
$\mathrm{SWR}=\frac{1+\Gamma}{1-\Gamma}$
$\Rightarrow \quad 2=\frac{1+\Gamma}{1-\Gamma}$
 Voltage reflection coefficient
$\Rightarrow \quad \Gamma=\frac{1}{3}$
 Power reflection coefficient

\[
\therefore \quad \frac{P_{\text {ref }}}{P_{\text {inc }}}=\Gamma^{2}=\frac{1}{9}
\]

$\Rightarrow 11.11 \%$ of $P_{i}$ is reflected
\end{tcolorbox}

\vspace{1.5em}
\hrule
\vspace{1.5em}

\section*{Question 170: Transmission Lines (GATE EC 2004)}
\addcontentsline{toc}{section}{Question 170: Transmission Lines (GATE EC 2004)}
\noindent \textbf{\color{primary}MCQ} \hfill \textbf{\color{secondary}1 Mark}


\noindent A plane electromagnetic wave propagating in free space is incident normally on a large slab of loss-less, non-magnetic, dielectric material with $\varepsilon >\varepsilon _{0}$. Maxima and minima are observed when the electric field is measured in front of the slab. The maximum electric field is found to be 5 times the minimum field. The intrinsic impedance of the medium should be


\begin{enumerate}[label=(\Alph*)]
    \item $120 \pi \Omega$
    \item $60 \pi \Omega$
    \item $600 \pi \Omega$
    \item \textbf{(Correct)} $24 \pi \Omega$
\end{enumerate}

\noindent \textbf{Correct Answer:} (D)

\begin{tcolorbox}[solbox, title=Detailed Solution -- Question 170]
\[
\begin{aligned} \text{SWR} &=\frac{E_{\max }}{E_{\min }} \\ \Rightarrow \quad s &=\frac{5 E_{\min }}{E_{\min }}=5 \\ & 5=\frac{1+|\Gamma|}{1-|\Gamma|}\;|\Gamma|=\frac{2}{3} \\ & \eta_{2}=\frac{120 \pi}{\sqrt{\epsilon_{r}}} \text { and } \eta_{2}=120 \pi \\ \text { As } & \eta_{2} < \eta_{1} \\ \Rightarrow -\frac{2}{3}&=\frac{\eta_{2}-120 \pi}{\eta+120 \pi} \\ \Rightarrow \eta_{2}&=24 \pi \end{aligned}
\]
\end{tcolorbox}

\vspace{1.5em}
\hrule
\vspace{1.5em}

\section*{Question 171: Transmission Lines (GATE EC 2004)}
\addcontentsline{toc}{section}{Question 171: Transmission Lines (GATE EC 2004)}
\noindent \textbf{\color{primary}MCQ} \hfill \textbf{\color{secondary}1 Mark}


\noindent Consider an impedance Z = R + jX marked with point P in an impedance Smith chart as shown in Fig. The movement from point P along a constant resistance circle in the clockwise direction by an angle 45$^{\circ}$  is equivalent to 

\begin{center}\includegraphics[max width=0.85\linewidth]{images/img\_35.jpg}\end{center}

\begin{center}\includegraphics[max width=0.85\linewidth]{images/img\_35.jpg}\end{center}


\begin{enumerate}[label=(\Alph*)]
    \item \textbf{(Correct)} adding an inductance in series with Z
    \item adding a capacitance in series with Z
    \item adding an inductance in shunt across Z
    \item adding a capacitance in shunt across Z
\end{enumerate}

\noindent \textbf{Correct Answer:} (A)

\begin{tcolorbox}[solbox, title=Detailed Solution -- Question 171]
Movement on constant resistance circle by an $\angle 45^{\circ}$
in clockwise direction.

R is same and reactance increases i.e. addition of
an inductance in series with Z.
\end{tcolorbox}

\vspace{1.5em}
\hrule
\vspace{1.5em}

\section*{Question 172: Uniform Plane Waves (GATE EC 2004)}
\addcontentsline{toc}{section}{Question 172: Uniform Plane Waves (GATE EC 2004)}
\noindent \textbf{\color{primary}MCQ} \hfill \textbf{\color{secondary}1 Mark}


\noindent If $\vec{E}=(\hat{a}_{x}+j\hat{a}_{y})e^{jkz-k\omega t}$ and $\vec{H}=(k/\omega \mu )(\hat{a}_{y}+j\hat{a}_{x})e^{jkz-j\omega t}$, the time-averaged Poynting vector is


\begin{enumerate}[label=(\Alph*)]
    \item \textbf{(Correct)} null vector
    \item $(k/\omega \mu )\hat{a}_{z}$
    \item $(2k/\omega \mu )\hat{a}_{z}$
    \item $(k/2\omega \mu )\hat{a}_{z}$
\end{enumerate}

\noindent \textbf{Correct Answer:} (A)

\begin{tcolorbox}[solbox, title=Detailed Solution -- Question 172]
\[
\begin{aligned} \vec{P}_{\mathrm{avg}}=& \frac{1}{2} \text{Re}\left[\vec{E} \times \vec{H}^{*}\right] \\ =& \frac{1}{2} \text{Re}\left[\left(\hat{a}_{x}+j \hat{a}_{y}\right) e^{j k z-j \omega t}\right.\\ &\left.\times\frac{k}{ \omega\mu}\left(\hat{a}_{x}+j \hat{a}_{y} \right)e^{j k z-j \omega t} \right ]\\ \vec{P}_{\mathrm{avg}}=& \frac{k}{2 \mu \omega}[0]=0 \\ \left[\begin{array}{cc}1 &j\\-j&1\end{array}\right]=&1+j^{2}=1-1=0 \end{aligned}
\]
\end{tcolorbox}

\vspace{1.5em}
\hrule
\vspace{1.5em}

\section*{Question 173: Miscellaneous Topics (GATE EC 2004)}
\addcontentsline{toc}{section}{Question 173: Miscellaneous Topics (GATE EC 2004)}
\noindent \textbf{\color{primary}MCQ} \hfill \textbf{\color{secondary}1 Mark}


\noindent In a microwave test bench, why is the microwave signal amplitude modulated at 1 kHz


\begin{enumerate}[label=(\Alph*)]
    \item To increase the sensitivity of measurement
    \item To transmit the signal to a far-off place
    \item To study amplitude modulations
    \item \textbf{(Correct)} Because crystal detector fails at microwave frequencies
\end{enumerate}

\noindent \textbf{Correct Answer:} (D)


\vspace{1.5em}
\hrule
\vspace{1.5em}

\section*{Question 174: Transmission Lines (GATE EC 2004)}
\addcontentsline{toc}{section}{Question 174: Transmission Lines (GATE EC 2004)}
\noindent \textbf{\color{primary}MCQ} \hfill \textbf{\color{secondary}1 Mark}


\noindent Consider a 300$\Omega$, quarter-wave long (at 1 GHz) transmission line as shown in Fig. It is connected to a 10V, 50$\Omega$ source at one end and is left open circuited at the other end. The magnitude of the voltage at the open circuit end of the line is

\begin{center}\includegraphics[max width=0.85\linewidth]{images/img\_19.jpg}\end{center}

\begin{center}\includegraphics[max width=0.85\linewidth]{images/img\_19.jpg}\end{center}


\begin{enumerate}[label=(\Alph*)]
    \item 10 V
    \item 5 V
    \item \textbf{(Correct)} 60 V
    \item 60/7 V
\end{enumerate}

\noindent \textbf{Correct Answer:} (C)

\begin{tcolorbox}[solbox, title=Detailed Solution -- Question 174]
The $\lambda / 4$ line has $Z_{\text {in }}^{\prime}=0 as Z_{L}=\infty$
 Current $I_{\text {in }}$ into the line $=\frac{V}{50}=0.2 \mathrm{A}$
 Using current equation on the line
$I(x)=I_{\text {in }}=I_{L} \cos \beta x+\frac{j V_{L}}{Z_{0}} \sin \beta x$
 with $x=\lambda / 4, I_{L}=0  \; as \;   Z_{L}=\infty$
 Open circuit end
$V_{L}=60 \mathrm{V}$
\end{tcolorbox}

\vspace{1.5em}
\hrule
\vspace{1.5em}

\section*{Question 175: Uniform Plane Waves (GATE EC 2004)}
\addcontentsline{toc}{section}{Question 175: Uniform Plane Waves (GATE EC 2004)}
\noindent \textbf{\color{primary}MCQ} \hfill \textbf{\color{secondary}1 Mark}


\noindent A parallel plate air-filled capacitor has plate area of $10^{-4} m^{2}$ and plate
separation of $10^{-3}$ m. It is connected to a 0.5V, 3.6 GHz source. The
magnitude of the displacement current is ($\varepsilon =\frac{1}{36\pi}10^{-9} F/m$)


\begin{enumerate}[label=(\Alph*)]
    \item \textbf{(Correct)} 10 mA
    \item 100 mA
    \item 10A
    \item 1.59 mA
\end{enumerate}

\noindent \textbf{Correct Answer:} (A)

\begin{tcolorbox}[solbox, title=Detailed Solution -- Question 175]
Displacement current,

\[
\begin{aligned} I_{d} &=A J_{d} \\ &=A \frac{\partial D}{\partial t}=A \in \frac{\partial E}{\partial t} \\ \left|I_{d}\right| &=|A \in \omega E|=A \omega \in \frac{V}{d} \end{aligned}
\]

 After putting values, we get,
$I_{d}=10 \mathrm{mA}$
\end{tcolorbox}

\vspace{1.5em}
\hrule
\vspace{1.5em}

\section*{Question 176: Antennas \& Radiating Systems (GATE EC 2004)}
\addcontentsline{toc}{section}{Question 176: Antennas \& Radiating Systems (GATE EC 2004)}
\noindent \textbf{\color{primary}MCQ} \hfill \textbf{\color{secondary}1 Mark}


\noindent Consider a lossless antenna with a directive gain of +6dB. If 1 mW of power is fed to it the total power radiated by the antenna will be


\begin{enumerate}[label=(\Alph*)]
    \item \textbf{(Correct)} 4 mW
    \item 1 mW
    \item 7 mW
    \item 1/4 mW
\end{enumerate}

\noindent \textbf{Correct Answer:} (A)

\begin{tcolorbox}[solbox, title=Detailed Solution -- Question 176]
When an antenna is lossless then the antenna
efficiency is 100%. In such case whole of the
power fed to antenna is radiated in the form of
radiation. Therefore, if power fed to antenna is
1 mW then total radiated power is also 1 mW.
\end{tcolorbox}

\vspace{1.5em}
\hrule
\vspace{1.5em}

\section*{Question 177: Waveguides \& Optical Fibers (GATE EC 2004)}
\addcontentsline{toc}{section}{Question 177: Waveguides \& Optical Fibers (GATE EC 2004)}
\noindent \textbf{\color{primary}MCQ} \hfill \textbf{\color{secondary}1 Mark}


\noindent The phase velocity of an electromagnetic wave propagating in a hollow metallic rectangular waveguide in the $TE_{10}$
mode is


\begin{enumerate}[label=(\Alph*)]
    \item equal to its group velocity
    \item less than the velocity of light in free space
    \item equal to the velocity of light in free space
    \item \textbf{(Correct)} greater than the velocity of light in free space
\end{enumerate}

\noindent \textbf{Correct Answer:} (D)


\vspace{1.5em}
\hrule
\vspace{1.5em}

\section*{Question 178: Antennas \& Radiating Systems (GATE EC 2003)}
\addcontentsline{toc}{section}{Question 178: Antennas \& Radiating Systems (GATE EC 2003)}
\noindent \textbf{\color{primary}MCQ} \hfill \textbf{\color{secondary}1 Mark}


\noindent Two identical antennas are placed in the $\theta =\pi /2$ plane as shown in Fig. The elements have equal amplitude excitation with $180^{\circ}$ polarity difference, operating at wavelength $\lambda$. The correct value of the magnitude of the far-zone resultant electric field strength normalized with that of a single element, both computed for $\phi$ = 0, is  

\begin{center}\includegraphics[max width=0.85\linewidth]{images/img\_51.jpg}\end{center}

\begin{center}\includegraphics[max width=0.85\linewidth]{images/img\_51.jpg}\end{center}


\begin{enumerate}[label=(\Alph*)]
    \item $2\cos (\frac{2\pi s}{\lambda })$
    \item $2\sin (\frac{2\pi s}{\lambda })$
    \item $2\cos (\frac{\pi s}{\lambda })$
    \item \textbf{(Correct)} $2\sin (\frac{\pi s}{\lambda })$
\end{enumerate}

\noindent \textbf{Correct Answer:} (D)

\begin{tcolorbox}[solbox, title=Detailed Solution -- Question 178]
$0 = 90^{\circ}$ means XY plane means Z = O plane
 Far zone calculations of field at $\phi$ = O means point
 on X-axis.
 Angle made by the array axis to the point of
 consideration, if taken as $\theta$.

\begin{center}\includegraphics[max width=0.85\linewidth]{images/img\_86.jpg}\end{center}

\begin{center}\includegraphics[max width=0.85\linewidth]{images/img\_86.jpg}\end{center}

\[
\begin{aligned} \text { Normalized array factor }&=2\left|\cos \frac{\psi}{2}\right| \\ \psi&=\beta d \cos \theta+\alpha \\ d&=\sqrt{2} s, \theta=45^{\circ}, \alpha=180^{\circ} \\ \Rightarrow 2\left|\cos \frac{\psi}{2}\right| &=2\cos\left[\frac{\beta d \cos \theta+\alpha}{2} \right ]\\ =& 2 \cos \left[\frac{2 \pi}{\lambda .2} \sqrt{2} \operatorname{scos} 45^{\circ}+\frac{180^{\circ}}{2}\right] \\ =& 2 \cos \left[\frac{\pi s}{\lambda}+90^{\circ}\right]\\ = & 2 \sin \left(\frac{\pi \mathrm{s}}{\lambda}\right) \end{aligned}
\]
\end{tcolorbox}

\vspace{1.5em}
\hrule
\vspace{1.5em}

\section*{Question 179: Waveguides \& Optical Fibers (GATE EC 2003)}
\addcontentsline{toc}{section}{Question 179: Waveguides \& Optical Fibers (GATE EC 2003)}
\noindent \textbf{\color{primary}MCQ} \hfill \textbf{\color{secondary}1 Mark}


\noindent A rectangular metal wave guide filled with a dielectric material of relative permittivity $\varepsilon _{r}$ = 4 has the inside dimensions 3.0 cm $\times$ 1.2 cm. The cut-off frequency for the dominant mode is


\begin{enumerate}[label=(\Alph*)]
    \item \textbf{(Correct)} 2.5 GHz
    \item 5.0 GHz
    \item 10.0 GHz
    \item 12.5 GHz
\end{enumerate}

\noindent \textbf{Correct Answer:} (A)

\begin{tcolorbox}[solbox, title=Detailed Solution -- Question 179]
\[
\begin{aligned} v &=\frac{c}{\sqrt{\epsilon_{r}}}=\frac{3 \times 10^{8}}{2}=1.5 \times 10^{8} \\ f_{c} &=\frac{v}{\lambda_{c}}=\frac{1.5 \times 10^{8}}{6 \times 10^{-2}}=2.5 \mathrm{GHz} \\ \text { or, } f_{c} &=\frac{v}{2 a}=\frac{1.5 \times 10^{8}}{6 \times 10^{-2}}=2.5 \mathrm{GHz} \end{aligned}
\]
\end{tcolorbox}

\vspace{1.5em}
\hrule
\vspace{1.5em}

\section*{Question 180: Transmission Lines (GATE EC 2003)}
\addcontentsline{toc}{section}{Question 180: Transmission Lines (GATE EC 2003)}
\noindent \textbf{\color{primary}MCQ} \hfill \textbf{\color{secondary}1 Mark}


\noindent A short - circuited stub is shunt connected to a transmission line as shown in fig. If $Z_{0}$ = 50 ohm, the admittance Y seen at the junction of the stub and the transmission line is 

\begin{center}\includegraphics[max width=0.85\linewidth]{images/img\_69.jpg}\end{center}

\begin{center}\includegraphics[max width=0.85\linewidth]{images/img\_69.jpg}\end{center}


\begin{enumerate}[label=(\Alph*)]
    \item \textbf{(Correct)} (0.01 - j0.02) mho
    \item (0.02 - j0.01) mho
    \item (0.04 - j0.02) mho
    \item (0.02 + j0) mho
\end{enumerate}

\noindent \textbf{Correct Answer:} (A)

\begin{tcolorbox}[solbox, title=Detailed Solution -- Question 180]
\begin{center}\includegraphics[max width=0.85\linewidth]{images/img\_62.jpg}\end{center}

\begin{center}\includegraphics[max width=0.85\linewidth]{images/img\_62.jpg}\end{center}

\[
\begin{aligned} \beta &=\frac{2 \pi}{\lambda} \\ \text { For } Y_{d} & \beta d=\frac{2 \pi}{\lambda} \times \frac{\lambda}{2}=\pi \\ & Z_{d}=\frac{Z_{0}\left[Z_{L}+j Z_{0} \tan \beta d\right.}{Z_{0}+j Z_{L} \tan \beta d} \\ & Z_{d}=\frac{50[100+j 50 \tan \pi]}{(50+j 100 \tan \pi)}=100 \\ Y_{d}&=\frac{1}{Z_{d}}=\frac{1}{100}=0.01\\ For Y_{s} \quad \beta d&=\frac{2 \pi}{\lambda} \times \frac{\lambda}{8}=\frac{\pi}{4}, Z_{L} \rightarrow 0 \\ Z_{s}&=\frac{Z_{0}\left[Z_{L}+j Z_{0} \tan \frac{\pi}{4}\right]}{\left[Z_{0}+j Z_{L} \tan \frac{\pi}{4}\right]}=j Z_{0} \\ Y_{s}&=\frac{1}{Z_{s}}=\frac{1}{j Z_{0}}=-0.02 j\\ \therefore\quad Y&=Y_{d}+Y_{s}=0.01-0.02 j\end{aligned}
\]
\end{tcolorbox}

\vspace{1.5em}
\hrule
\vspace{1.5em}

\section*{Question 181: Uniform Plane Waves (GATE EC 2003)}
\addcontentsline{toc}{section}{Question 181: Uniform Plane Waves (GATE EC 2003)}
\noindent \textbf{\color{primary}MCQ} \hfill \textbf{\color{secondary}1 Mark}


\noindent If the electric field intensity associated with a uniform plane electromagnetic wave traveling in a perfect dielectric medium is given by
 $E(z,t)=10\cos (2\pi 10^{7}t-0.1\pi z)$ V/m, then the velocity of the traveling wave is


\begin{enumerate}[label=(\Alph*)]
    \item $3.00\times 10^{8}$ m/sec
    \item \textbf{(Correct)} $2.00\times 10^{8}$ m/sec
    \item $6.28\times 10^{7}$ m/sec
    \item $2.00\times 10^{7}$ m/sec
\end{enumerate}

\noindent \textbf{Correct Answer:} (B)

\begin{tcolorbox}[solbox, title=Detailed Solution -- Question 181]
\[
\begin{array}{l} \omega=2 \pi \times 10^{7} \\ \beta=0.1 \pi \end{array}
\]

 Comparing with $A \cos (\omega t-\beta z)$

\[
v=\frac{\omega}{\beta}=\frac{2 \pi \times 10^{7}}{0.1 \pi}=2 \times 10^{8} \mathrm{m}
\]
\end{tcolorbox}

\vspace{1.5em}
\hrule
\vspace{1.5em}

\section*{Question 182: Uniform Plane Waves (GATE EC 2003)}
\addcontentsline{toc}{section}{Question 182: Uniform Plane Waves (GATE EC 2003)}
\noindent \textbf{\color{primary}MCQ} \hfill \textbf{\color{secondary}1 Mark}


\noindent A uniform plane wave traveling in air is incident on the plane boundary between air and another dielectric medium with $\varepsilon _{r}$ = 4. The reflection coefficient for the normal incidence, is


\begin{enumerate}[label=(\Alph*)]
    \item zero
    \item $0.5\angle 180^{\circ}$
    \item $0.333\angle 0^{\circ}$
    \item \textbf{(Correct)} $0.333\angle 180^{\circ}$
\end{enumerate}

\noindent \textbf{Correct Answer:} (D)

\begin{tcolorbox}[solbox, title=Detailed Solution -- Question 182]
\[
\begin{aligned} \Gamma &=\frac{\eta_{2}-\eta_{1}}{\eta_{z}+\eta_{1}} \\ &=\frac{\sqrt{\frac{\mu_{0}}{\epsilon_{0} \epsilon_{r}}}-\sqrt{\frac{\mu_{0}}{E_{0}}}}{\sqrt{\frac{\mu_{0}}{\epsilon_{0} \epsilon_{r}}}+\sqrt{\frac{\mu_{0}}{\epsilon_{0}}}}=\frac{\frac{1}{\sqrt{\epsilon_{r}}}-1}{\frac{1}{\sqrt{\epsilon_{r}}}+1},\epsilon_{r}=4\\ &=\frac{-1 / 2}{3 / 2}=\frac{-1}{3} \Rightarrow 0.333 \angle 180^{\circ} \end{aligned}
\]
\end{tcolorbox}

\vspace{1.5em}
\hrule
\vspace{1.5em}

\section*{Question 183: Basics of Electromagnetics \& Maxwell's Eqns (GATE EC 2003)}
\addcontentsline{toc}{section}{Question 183: Basics of Electromagnetics \& Maxwell's Eqns (GATE EC 2003)}
\noindent \textbf{\color{primary}MCQ} \hfill \textbf{\color{secondary}1 Mark}


\noindent If the electric field intensity is given by $E = (xu_{x}+yu_{y}+zu_{z})$ volt/m, the potential difference between X(2,0,0) and Y(1,2,3) is


\begin{enumerate}[label=(\Alph*)]
    \item $+1$ volt
    \item $-1$ volt
    \item \textbf{(Correct)} $+5$ volt
    \item $+6$ volt
\end{enumerate}

\noindent \textbf{Correct Answer:} (C)

\begin{tcolorbox}[solbox, title=Detailed Solution -- Question 183]
\[
\begin{aligned} V &=-\int \vec{E} \cdot \vec{d} \mid \\ &=-\left[\int_{1}^{2} x d x \vec{u}_{x}+\int_{2}^{0} y d y \vec{u}_{y}+\int_{3}^{0} z d z \vec{u}_{z}\right] \\ &=-\left[\left.\frac{x^{2}}{2}\right|_{1} ^{2}+\left.\frac{y^{2}}{2}\right|_{2} ^{0}+\left.\frac{z^{2}}{2}\right|_{3}\right] \\ &=-\frac{1}{2}\left[2^{2}-1^{2}+0^{2}-2^{2}+0^{2}-3^{2}\right] \\ &=-\frac{1}{2} x-10=5 \mathrm{V} \end{aligned}
\]
\end{tcolorbox}

\vspace{1.5em}
\hrule
\vspace{1.5em}

\section*{Question 184: Basics of Electromagnetics \& Maxwell's Eqns (GATE EC 2003)}
\addcontentsline{toc}{section}{Question 184: Basics of Electromagnetics \& Maxwell's Eqns (GATE EC 2003)}
\noindent \textbf{\color{primary}MCQ} \hfill \textbf{\color{secondary}1 Mark}


\noindent Medium 1 has the electrical permittivity $\varepsilon _{1}=1.5\varepsilon _{0}$ farad/m and occupies the region to the left of x = 0 plane. Medium 2 has the electrical permittivity $\varepsilon _{2}=2.5\varepsilon _{0}$ farad/m and occupies the region to the right of x = 0 plane. If $E_{1}$ in medium 1 is $E_{1}$ = $(2u_{x}-3u_{y}+1u_{z})$ volt/m, then  $E_{2}$ in medium 2 is


\begin{enumerate}[label=(\Alph*)]
    \item $(2.0u_{x}-7.5u_{y}+2.5u_{z})$ volt/m
    \item $(2.0u_{x}-2.0u_{y}+0.6u_{z})$ volt/m
    \item \textbf{(Correct)} $(1.2u_{x}-3.0u_{y}+1.0u_{z})$ volt/m
    \item $(1.2u_{x}-2.0u_{y}+0.6u_{z})$ volt/m
\end{enumerate}

\noindent \textbf{Correct Answer:} (C)

\begin{tcolorbox}[solbox, title=Detailed Solution -- Question 184]
\[
\begin{aligned} \vec{E}_{1} &=2 u_{x}-3 u_{y}+1 u_{z} \\ \vec{E}_{1 t} &=-3 u_{y}+u_{z}=E_{2 t} &(x=0 \text { plane}) \\ \vec{E}_{1 n} &=2 u_{x} \\ \vec{D}_{1 n} &=\vec{D}_{2 n}=\epsilon \vec{E} \\ \Rightarrow \epsilon_{1} \vec{E}_{1 n} &=\epsilon_{2} \vec{E}_{2 n} \\ &=1.5 \epsilon_{0} \cdot 2 u_{x}=2.5 \epsilon_{0} \cdot \vec{E}_{2 n} \\ \vec{E}_{2 n} &=\frac{3}{2.5} u_{x}=1.2 u_{x} \\ \vec{E}_{2} &=\vec{E}_{2 t}+\vec{E}_{2 n} \\ \vec{E}_{2} &=-3 u_{y}+u_{z}+1.2 u_{x} \\ &=1.2u_{x}-3.0u_{y}+1.0u_{z} \end{aligned}
\]
\end{tcolorbox}

\vspace{1.5em}
\hrule
\vspace{1.5em}

\section*{Question 185: Uniform Plane Waves (GATE EC 2003)}
\addcontentsline{toc}{section}{Question 185: Uniform Plane Waves (GATE EC 2003)}
\noindent \textbf{\color{primary}MCQ} \hfill \textbf{\color{secondary}1 Mark}


\noindent The depth of penetration of electromagnetic wave in a medium having conductivity $\sigma$ at a frequency of 1 MHz is 25 cm. The depth of penetration at a frequency of 4 MHz will be


\begin{enumerate}[label=(\Alph*)]
    \item 6.25 dm
    \item \textbf{(Correct)} 12.50 cm
    \item 50.00 cm
    \item 100.00 cm
\end{enumerate}

\noindent \textbf{Correct Answer:} (B)

\begin{tcolorbox}[solbox, title=Detailed Solution -- Question 185]
\[
\begin{array}{l} \text { For conductors, } \delta=\sqrt{\frac{2}{\omega \mu \sigma}} \\ \qquad \begin{aligned} \delta \propto & \frac{1}{\sqrt{f}} \\ \frac{\delta_{2}}{\delta_{1}}=& \sqrt{\frac{f_{1}}{f_{2}}} \\ \delta_{2}=& \sqrt{\frac{1 \mathrm{MHz}}{4 \mathrm{MHz}}} \times 25 \mathrm{cm}=12.5 \mathrm{cm} \end{aligned} \end{array}
\]
\end{tcolorbox}

\vspace{1.5em}
\hrule
\vspace{1.5em}

\section*{Question 186: Basics of Electromagnetics \& Maxwell's Eqns (GATE EC 2003)}
\addcontentsline{toc}{section}{Question 186: Basics of Electromagnetics \& Maxwell's Eqns (GATE EC 2003)}
\noindent \textbf{\color{primary}MCQ} \hfill \textbf{\color{secondary}1 Mark}


\noindent The unit of $\bigtriangledown \times H$ is


\begin{enumerate}[label=(\Alph*)]
    \item Ampere
    \item Ampere/meter
    \item \textbf{(Correct)} $Ampere/meter^{2}$
    \item Ampere-meter
\end{enumerate}

\noindent \textbf{Correct Answer:} (C)

\begin{tcolorbox}[solbox, title=Detailed Solution -- Question 186]
\[
\begin{aligned} \nabla \times \vec{H} &=\frac{\partial \bar{D}}{\partial t}+\vec{J} \\ \frac{1}{m} \times \frac{\text { Ampere }}{m} &=\frac{A}{m^{2}} \end{aligned}
\]
\end{tcolorbox}

\vspace{1.5em}
\hrule
\vspace{1.5em}

\section*{Question 187: Antennas \& Radiating Systems (GATE EC 2002)}
\addcontentsline{toc}{section}{Question 187: Antennas \& Radiating Systems (GATE EC 2002)}
\noindent \textbf{\color{primary}MCQ} \hfill \textbf{\color{secondary}1 Mark}


\noindent A person with a receiver is 5 Km away from the transmitter. What is the distance
that this person must move further to detect a 3-dB decrease in signal strength?


\begin{enumerate}[label=(\Alph*)]
    \item 942m
    \item \textbf{(Correct)} 2070m
    \item 4978m
    \item 5320m
\end{enumerate}

\noindent \textbf{Correct Answer:} (B)

\begin{tcolorbox}[solbox, title=Detailed Solution -- Question 187]
\[
\begin{array}{l} \text { Signal strength }=\frac{P}{4 \pi R^{2}} \\ P_{1}=\frac{P}{4 \pi(5000)^{2}} \text { and } P_{2}=\frac{P}{4 \pi(x)^{2}} \end{array}
\]

 Given, $\quad \frac{P_{2}}{P_{1}}=\frac{1}{2}$ (3 dB decrease)

\[
\begin{aligned} \Rightarrow\left(\frac{5000}{x}\right)^{2}&=\frac{1}{2} \\ \Rightarrow \quad x&=5000 \sqrt{2} \end{aligned}
\]

$\therefore$Required distance
$=5000 \sqrt{2}-5000=2071 \mathrm{m}$
\end{tcolorbox}

\vspace{1.5em}
\hrule
\vspace{1.5em}

\section*{Question 188: Basics of Electromagnetics \& Maxwell's Eqns (GATE EC 2002)}
\addcontentsline{toc}{section}{Question 188: Basics of Electromagnetics \& Maxwell's Eqns (GATE EC 2002)}
\noindent \textbf{\color{primary}MCQ} \hfill \textbf{\color{secondary}1 Mark}


\noindent The electric field on the surface of a perfect conductor is 2 V/m. The conductor is
immersed in water with $\epsilon =80\epsilon_{0}$. The surface charge density on the conductor is ($\epsilon =10^{-9}/(36\pi)F/m$ )


\begin{enumerate}[label=(\Alph*)]
    \item $0 \; C/m^{2}$
    \item $2 \; C/m^{2}$
    \item $1.8 \times 10^{-11} \; C/m^{2}$
    \item \textbf{(Correct)} $1.14 \times 10^{-9} \; C/m^{2}$
\end{enumerate}

\noindent \textbf{Correct Answer:} (D)

\begin{tcolorbox}[solbox, title=Detailed Solution -- Question 188]
The surface change density is equally to normal
 flux density 

\[
\begin{aligned} D_{\text {normal }} &=\rho_{s} \\ D &=\rho_{s}=\epsilon E=80 \cdot \epsilon \cdot 2 \end{aligned}
\]
\end{tcolorbox}

\vspace{1.5em}
\hrule
\vspace{1.5em}

\section*{Question 189: Uniform Plane Waves (GATE EC 2002)}
\addcontentsline{toc}{section}{Question 189: Uniform Plane Waves (GATE EC 2002)}
\noindent \textbf{\color{primary}MCQ} \hfill \textbf{\color{secondary}1 Mark}


\noindent Distilled water at 25$^{\circ}$ C is characterized by $\sigma$ = 1.7 $\times$ $10^{-4}$ mho/m and $\epsilon =78\epsilon_{0}$ at a frequency of 3 GHz. Its loss tangent $tan\delta$
is (  $\epsilon=10^{-9} /(36\pi)$ F/m)


\begin{enumerate}[label=(\Alph*)]
    \item \textbf{(Correct)} $1.3 \times 10^{-5}$
    \item $1.3 \times 10^{-3}$
    \item $1.7 \times 10^{-4}$/78
    \item $1.7 \times 10^{-4}/(78\epsilon_{0})$
\end{enumerate}

\noindent \textbf{Correct Answer:} (A)

\begin{tcolorbox}[solbox, title=Detailed Solution -- Question 189]
\[
\begin{aligned} \operatorname{Loss} \text { tangent } &=\frac{\sigma}{\omega \in} \\ &=\frac{1.7 \times 10^{-4}}{2 \pi \times 3 \times 10^{9} \times 78 \epsilon_{0}} \\ &=\frac{1.7 \times 10^{-4} \times 9 \times 10^{9}}{3 \times 10^{9} \times 39} \\ &=0.13 \times 10^{-4}=1.3 \times 10^{-5} \end{aligned}
\]
\end{tcolorbox}

\vspace{1.5em}
\hrule
\vspace{1.5em}

\section*{Question 190: Uniform Plane Waves (GATE EC 2002)}
\addcontentsline{toc}{section}{Question 190: Uniform Plane Waves (GATE EC 2002)}
\noindent \textbf{\color{primary}MCQ} \hfill \textbf{\color{secondary}1 Mark}


\noindent A plane wave is characterized by 
$\vec{E}=(0.5\hat{x}+\hat{y}e^{\frac{j\pi }{2}})e^{jwt-jkz}$.
 This wave is


\begin{enumerate}[label=(\Alph*)]
    \item linearly polarized
    \item circularly polarized
    \item \textbf{(Correct)} elliptically polarized
    \item unpolarized
\end{enumerate}

\noindent \textbf{Correct Answer:} (C)

\begin{tcolorbox}[solbox, title=Detailed Solution -- Question 190]
$E_{x}$ and $E_{y}$ out of phase by $90^{\circ}$, unequal amplitudes.
 This is elliptical polarization.
\end{tcolorbox}

\vspace{1.5em}
\hrule
\vspace{1.5em}

\section*{Question 191: Waveguides \& Optical Fibers (GATE EC 2002)}
\addcontentsline{toc}{section}{Question 191: Waveguides \& Optical Fibers (GATE EC 2002)}
\noindent \textbf{\color{primary}MCQ} \hfill \textbf{\color{secondary}1 Mark}


\noindent The phase velocity for the TE10-mode in an air-filled rectangular waveguide is (c
is the velocity of plane waves in free space)


\begin{enumerate}[label=(\Alph*)]
    \item less than c
    \item equal to c
    \item \textbf{(Correct)} greater than c
    \item none of the above
\end{enumerate}

\noindent \textbf{Correct Answer:} (C)

\begin{tcolorbox}[solbox, title=Detailed Solution -- Question 191]
\[
\begin{aligned} V_{p}&=\frac{c}{\cos \theta}=\frac{c}{\sqrt{1-\left(\frac{f_{c}}{f}\right)^{2}}} \\ \text { with } f&>f_{c^{\prime}} V_{p}>c \text { always. } \end{aligned}
\]
\end{tcolorbox}

\vspace{1.5em}
\hrule
\vspace{1.5em}

\section*{Question 192: Transmission Lines (GATE EC 2002)}
\addcontentsline{toc}{section}{Question 192: Transmission Lines (GATE EC 2002)}
\noindent \textbf{\color{primary}MCQ} \hfill \textbf{\color{secondary}1 Mark}


\noindent In an impedance Smith chart, a clockwise movement along a constant resistance
circle gives rise to


\begin{enumerate}[label=(\Alph*)]
    \item a decrease in the value of reactance
    \item \textbf{(Correct)} an increase in the value of reactance
    \item no change in the reactance value
    \item no change in the impedance value
\end{enumerate}

\noindent \textbf{Correct Answer:} (B)

\begin{tcolorbox}[solbox, title=Detailed Solution -- Question 192]
Moving along clockwise direction is an inductive
increase in the load impedance
\end{tcolorbox}

\vspace{1.5em}
\hrule
\vspace{1.5em}

\section*{Question 193: Transmission Lines (GATE EC 2002)}
\addcontentsline{toc}{section}{Question 193: Transmission Lines (GATE EC 2002)}
\noindent \textbf{\color{primary}MCQ} \hfill \textbf{\color{secondary}1 Mark}


\noindent The VSWR can have any value between


\begin{enumerate}[label=(\Alph*)]
    \item 0 and 1
    \item -1 and +1
    \item 0 and $\infty$
    \item \textbf{(Correct)} 1 and $\infty$
\end{enumerate}

\noindent \textbf{Correct Answer:} (D)


\vspace{1.5em}
\hrule
\vspace{1.5em}

\section*{Question 194: Antennas \& Radiating Systems (GATE EC 2002)}
\addcontentsline{toc}{section}{Question 194: Antennas \& Radiating Systems (GATE EC 2002)}
\noindent \textbf{\color{primary}MCQ} \hfill \textbf{\color{secondary}1 Mark}


\noindent The line-of-sight communication requires the transmit and receive antenna to
face each other. If the transmit antenna is vertically polarized, for best reception
the receive antenna should be


\begin{enumerate}[label=(\Alph*)]
    \item horizontally polarized
    \item \textbf{(Correct)} vertically polarized
    \item at $45^{\circ}$ with respect to horizontal polarization
    \item at $45^{\circ}$ with respect to vertical polarization
\end{enumerate}

\noindent \textbf{Correct Answer:} (B)

\begin{tcolorbox}[solbox, title=Detailed Solution -- Question 194]
Both transmitting and receiving antennas should
be oriented correctly for maximum induction in the
receiver.
Transmitting antenna is vertical produced E field
is vertical inducing Efield is vertical so receiving
antenna has to be vertical.
\end{tcolorbox}

\vspace{1.5em}
\hrule
\vspace{1.5em}

\section*{Question 195: Antennas \& Radiating Systems (GATE EC 2001)}
\addcontentsline{toc}{section}{Question 195: Antennas \& Radiating Systems (GATE EC 2001)}
\noindent \textbf{\color{primary}MCQ} \hfill \textbf{\color{secondary}1 Mark}


\noindent In a uniform linear array, four isotropic radiating elements are spaced $\frac{\lambda }{4}$ apart.
The progressive phase shift between the elements required for forming the main
beam at $60^{\circ}$ off the end-fire is:


\begin{enumerate}[label=(\Alph*)]
    \item $- \pi$ radians
    \item $-\frac{\pi }{2}$ radians
    \item \textbf{(Correct)} $-\frac{\pi }{4}$ radians
    \item $-\frac{\pi }{8}$ radians
\end{enumerate}

\noindent \textbf{Correct Answer:} (C)

\begin{tcolorbox}[solbox, title=Detailed Solution -- Question 195]
$\theta_{\max }=60^{\circ}$ from end fire means $60^{\circ}$ from array axis.\\ Condition for maximum, $\psi \rightarrow 0$

\[
\begin{array}{l} E_{T}=E_{O} \frac{\sin \left(N_{\psi} / 2\right)}{\sin (\psi / 2)} \\ E_{T}=N E_{0}=E_{\max } \\ \psi=\alpha+\beta d \cos 60^{\circ}=0 \\ \alpha+\frac{2 \pi}{\lambda} \cdot \frac{\lambda}{4} \cdot \frac{1}{2}=0 \\ \alpha=-\frac{\pi}{4} \end{array}
\]
\end{tcolorbox}

\vspace{1.5em}
\hrule
\vspace{1.5em}

\section*{Question 196: Antennas \& Radiating Systems (GATE EC 2001)}
\addcontentsline{toc}{section}{Question 196: Antennas \& Radiating Systems (GATE EC 2001)}
\noindent \textbf{\color{primary}MCQ} \hfill \textbf{\color{secondary}1 Mark}


\noindent A medium wave radio transmitter operating at a wavelength of 492 m has a
tower antenna of height 124m. What is the radiation resistance of the antenna?


\begin{enumerate}[label=(\Alph*)]
    \item $25\Omega$
    \item \textbf{(Correct)} $36.5 \Omega$
    \item $50\Omega$
    \item $73 \Omega$
\end{enumerate}

\noindent \textbf{Correct Answer:} (B)

\begin{tcolorbox}[solbox, title=Detailed Solution -- Question 196]
$\frac{\lambda}{4}=\frac{492}{4}=123 \simeq L$
$\therefore \quad$It is a quarter wave monopole antenna.
 So, $\quad R_{a}=36.5 \Omega$
\end{tcolorbox}

\vspace{1.5em}
\hrule
\vspace{1.5em}

\section*{Question 197: Transmission Lines (GATE EC 2001)}
\addcontentsline{toc}{section}{Question 197: Transmission Lines (GATE EC 2001)}
\noindent \textbf{\color{primary}MCQ} \hfill \textbf{\color{secondary}1 Mark}


\noindent A uniform plane electromagnetic wave incident normally on a plane surface of a
dielectric material is reflected with a VSWR of 3. What is the percentage of
incident power that is reflected?


\begin{enumerate}[label=(\Alph*)]
    \item 10%
    \item \textbf{(Correct)} 25%
    \item 50%
    \item 75%
\end{enumerate}

\noindent \textbf{Correct Answer:} (B)

\begin{tcolorbox}[solbox, title=Detailed Solution -- Question 197]
\[
\begin{aligned} s&=\frac{1+|\Gamma|}{1-|\Gamma|} \\ \Rightarrow 3&=\frac{1+|\Gamma|}{1-|\Gamma|} \\ |\Gamma|&=0.5 \\ \frac{P_{r}}{P_{i}}&=|\Gamma|^{2}=0.25 \end{aligned}
\]

$\therefore \quad 25 \%$ of incident power is reflected.
\end{tcolorbox}

\vspace{1.5em}
\hrule
\vspace{1.5em}

\section*{Question 198: Uniform Plane Waves (GATE EC 2001)}
\addcontentsline{toc}{section}{Question 198: Uniform Plane Waves (GATE EC 2001)}
\noindent \textbf{\color{primary}MCQ} \hfill \textbf{\color{secondary}1 Mark}


\noindent A material has conductivity of $10^{-2}$
mho/m and a relative permittivity of 4. The
frequency at which the conduction current in the medium is equal to the
displacement current is


\begin{enumerate}[label=(\Alph*)]
    \item \textbf{(Correct)} 45 MHz
    \item 90 MHz
    \item 450 MHz
    \item 900 MHz
\end{enumerate}

\noindent \textbf{Correct Answer:} (A)

\begin{tcolorbox}[solbox, title=Detailed Solution -- Question 198]
\[
\begin{aligned} |\sigma E| &=|j \omega \in E| \\ \sigma &=2 \pi f \epsilon_{0} \epsilon_{r} \\ \Rightarrow \quad f &=\frac{\sigma}{2 \pi \times \epsilon_{0} \epsilon_{r}} \\ &=\frac{9 \times 10^{9} \times 2 \times 10^{-2}}{4} \\ f &=45 \times 10^{6}=45 \mathrm{MHz} \end{aligned}
\]
\end{tcolorbox}

\vspace{1.5em}
\hrule
\vspace{1.5em}

\section*{Question 199: Waveguides \& Optical Fibers (GATE EC 2001)}
\addcontentsline{toc}{section}{Question 199: Waveguides \& Optical Fibers (GATE EC 2001)}
\noindent \textbf{\color{primary}MCQ} \hfill \textbf{\color{secondary}1 Mark}


\noindent The dominant mode in a rectangular waveguide is $TE_{10}$, because this mode has


\begin{enumerate}[label=(\Alph*)]
    \item no attenuation
    \item no cut-off
    \item no magnetic field component
    \item \textbf{(Correct)} the highest cut-off wavelength
\end{enumerate}

\noindent \textbf{Correct Answer:} (D)

\begin{tcolorbox}[solbox, title=Detailed Solution -- Question 199]
Dominant mode in waveguide has lowest cut-off
frequency and hence the highest cut-off
wavelength
\end{tcolorbox}

\vspace{1.5em}
\hrule
\vspace{1.5em}

\section*{Question 200: Waveguides \& Optical Fibers (GATE EC 2001)}
\addcontentsline{toc}{section}{Question 200: Waveguides \& Optical Fibers (GATE EC 2001)}
\noindent \textbf{\color{primary}MCQ} \hfill \textbf{\color{secondary}1 Mark}


\noindent The phase velocity of waves propagating in a hollow metal waveguide is


\begin{enumerate}[label=(\Alph*)]
    \item \textbf{(Correct)} greater than the velocity of light in free space.
    \item less than the velocity of light in free space.
    \item equal to the velocity of light in free space.
    \item equal to the group velocity.
\end{enumerate}

\noindent \textbf{Correct Answer:} (A)

\begin{tcolorbox}[solbox, title=Detailed Solution -- Question 200]
\[
\begin{aligned} V_{p}&=\frac{c}{\cos \theta}=\frac{c}{\sqrt{1-\left(\frac{f_{c}}{f}\right)^{2}}} \\ \text { with } f& > f_{c^{\prime}} V_{p} > c \text { always. } \end{aligned}
\]
\end{tcolorbox}

\vspace{1.5em}
\hrule
\vspace{1.5em}

\section*{Question 201: Uniform Plane Waves (GATE EC 2001)}
\addcontentsline{toc}{section}{Question 201: Uniform Plane Waves (GATE EC 2001)}
\noindent \textbf{\color{primary}MCQ} \hfill \textbf{\color{secondary}1 Mark}


\noindent If a plane electromagnetic wave satisfies the equation $\frac{\partial^2 E_{X}}{\partial x^2}=C^{2}\frac{\partial^2X }{\partial t^2}$ the wave propagates in the


\begin{enumerate}[label=(\Alph*)]
    \item x-direction
    \item \textbf{(Correct)} z-direction
    \item y-direction
    \item xy plane at an angle of $45^{\circ}$
between the x and z directions
\end{enumerate}

\noindent \textbf{Correct Answer:} (B)

\begin{tcolorbox}[solbox, title=Detailed Solution -- Question 201]
The E field is a function of z and directed in X. i.e.

harmonic function of z or travelling in z direction.
\end{tcolorbox}

\vspace{1.5em}
\hrule
\vspace{1.5em}

\section*{Question 202: Transmission Lines (GATE EC 2001)}
\addcontentsline{toc}{section}{Question 202: Transmission Lines (GATE EC 2001)}
\noindent \textbf{\color{primary}MCQ} \hfill \textbf{\color{secondary}1 Mark}


\noindent A transmission line is distortion-less if


\begin{enumerate}[label=(\Alph*)]
    \item $RL=\frac{1}{GC}$
    \item RL=GC
    \item \textbf{(Correct)} LG=RC
    \item RG=LC
\end{enumerate}

\noindent \textbf{Correct Answer:} (C)


\vspace{1.5em}
\hrule
\vspace{1.5em}

\end{document}